\documentclass[11pt,letterpaper]{article}
\usepackage[T1]{fontenc}
\usepackage[utf8]{inputenc}
\usepackage{lmodern}
\usepackage[margin=1in,headheight=15pt]{geometry}
\usepackage{amsmath,amssymb,amsthm,mathtools}
\usepackage{microtype,booktabs,longtable,tabularx,array,enumitem}
\usepackage{xcolor,fancyhdr,needspace}
\definecolor{Ink}{HTML}{1F3648}
\definecolor{Accent}{HTML}{2A5F73}
\definecolor{TierE}{HTML}{1B5E20}
\definecolor{TierC}{HTML}{0D47A1}
\definecolor{TierA}{HTML}{8D4004}
\definecolor{TierI}{HTML}{4A4A4A}
\usepackage[colorlinks=true,linkcolor=Ink,citecolor=Accent,urlcolor=Accent]{hyperref}
\usepackage{xurl}
\usepackage[nameinlink,noabbrev]{cleveref}
\hypersetup{%
  pdftitle={Surreal Scalars in Theoretical Physics},
  pdfsubject={Merged research report: black-hole singularities, asymptotic structure,
    and the limits of replacing the scalar field},
  pdfkeywords={surreal numbers, surcomplex numbers, black-hole singularity,
    Kretschmann scalar, Hahn series, transseries, standard part}}
\pagestyle{fancy}\fancyhf{}
\fancyhead[L]{\small Surreal scalars in theoretical physics}
\fancyhead[R]{\small\thepage}
\renewcommand{\headrulewidth}{0.3pt}
\setlength{\parskip}{0.25em}
\setlength{\emergencystretch}{3em}
\widowpenalty=10000 \clubpenalty=10000 \displaywidowpenalty=10000
\setlist{itemsep=0.2em,topsep=0.3em}
\setcounter{tocdepth}{2}
\setcounter{secnumdepth}{3}
\numberwithin{equation}{section}

\newtheorem{theorem}{Theorem}[section]
\newtheorem{proposition}[theorem]{Proposition}
\newtheorem{lemma}[theorem]{Lemma}
\newtheorem{corollary}[theorem]{Corollary}
\theoremstyle{definition}
\newtheorem{definition}[theorem]{Definition}
\newtheorem{example}[theorem]{Example}
\newtheorem{computation}[theorem]{Computation}
\newtheorem{convention}[theorem]{Convention}
\newtheorem{assessment}[theorem]{Assessment}
\theoremstyle{remark}
\newtheorem{remark}[theorem]{Remark}
\newtheorem{warning}[theorem]{Warning}
\crefname{theorem}{Theorem}{Theorems}
\crefname{proposition}{Proposition}{Propositions}
\crefname{lemma}{Lemma}{Lemmas}
\crefname{corollary}{Corollary}{Corollaries}
\crefname{definition}{Definition}{Definitions}
\crefname{example}{Example}{Examples}
\crefname{computation}{Computation}{Computations}
\crefname{convention}{Convention}{Conventions}
\crefname{assessment}{Assessment}{Assessments}
\crefname{remark}{Remark}{Remarks}
\crefname{warning}{Warning}{Warnings}
\crefname{section}{Section}{Sections}
\crefname{table}{Table}{Tables}
\crefname{part}{Part}{Parts}
\crefname{appendix}{Appendix}{Appendices}

\newcommand{\No}{\mathbf{No}}
\newcommand{\SC}{\mathbf{No}[i]}
\newcommand{\R}{\mathbb R}
\newcommand{\C}{\mathbb C}
\newcommand{\Q}{\mathbb Q}
\newcommand{\Z}{\mathbb Z}
\newcommand{\N}{\mathbb N}
\newcommand{\eps}{\varepsilon}
\newcommand{\Ofin}{\mathcal O}
\newcommand{\mm}{\mathfrak m}
\newcommand{\Alg}{\mathcal A}
\newcommand{\Kret}{\mathcal K}
\newcommand{\lP}{\ell_{\mathrm P}}
\DeclareMathOperator{\st}{st}
\DeclareMathOperator{\supp}{supp}
\DeclareMathOperator{\val}{v}
\DeclareMathOperator{\Ric}{Ric}
\DeclareMathOperator{\Scal}{Scal}
\DeclareMathOperator{\FP}{FP}
\DeclareMathOperator{\cis}{cis}
\DeclareMathOperator{\diag}{diag}
\DeclareMathOperator{\PV}{PV}
\newcommand{\Oz}{\mathbf{Oz}}
\newcommand{\dd}{\mathrm d}
\newcommand{\ee}{\mathrm e}
\newcommand{\ii}{\mathrm i}
\newcommand{\ExpEK}{\operatorname{Exp}_{\mathrm{EK}}}
\newcommand{\sinEK}{\operatorname{Sin}_{\mathrm{EK}}}
\newcommand{\cosEK}{\operatorname{Cos}_{\mathrm{EK}}}
\newcommand{\Emm}{\mathrm{E}}
\newcommand{\abs}[1]{\lvert#1\rvert}
\newcommand{\Repo}{\texttt{VladimirReshetnikov/Surreal}}
\newcommand{\Commit}{\nolinkurl{39f2be6667ade51bca2b45daa47e289d69c09764}}
\newcommand{\src}[1]{{\small\textsf{[#1]}}}
\newcolumntype{L}[1]{>{\raggedright\arraybackslash}p{#1}}
\newcolumntype{Y}{>{\raggedright\arraybackslash}X}

% ---- the claim-status firewall tags -------------------------------------
\newcommand{\TE}{\textcolor{TierE}{\textsf{\bfseries[E]}}}
\newcommand{\TC}{\textcolor{TierC}{\textsf{\bfseries[C]}}}
\newcommand{\TA}{\textcolor{TierA}{\textsf{\bfseries[A]}}}
\newcommand{\TI}{\textcolor{TierI}{\textsf{\bfseries[I]}}}
\newcommand{\status}[1]{\par\smallskip\noindent{\small\textbf{Status.} #1}\par\smallskip}

\begin{document}
\hypersetup{pageanchor=false}
\begin{titlepage}
\centering
\vspace*{0.35in}
{\large\scshape A merged research report\par}
\vspace{0.4in}
{\LARGE\bfseries\color{Ink} Surreal Scalars in\\[0.2em]
Theoretical Physics\par}
\vspace{0.3in}
{\Large Black-Hole Singularities, Asymptotic Structure,\\[0.15em]
and the Limits of Replacing the Scalar Field\par}
\vspace{0.4in}
{\large September 2026\par}
\vspace{0.35in}
\begin{minipage}{0.93\textwidth}
\begin{abstract}
\noindent
Three independently written manuscripts asked whether Conway's surreal field
$\No$, or its algebraic closure $\SC$, can replace $\R$ and $\C$ in physics,
and in particular whether admitting infinite scalars removes a black-hole
singularity. This report is their union. All three reach the same verdict, and
the verdict is negative.

Surreal and surcomplex scalars make the hierarchy of quantities near a singular
regime \emph{exactly representable}. They do not make the spacetime regular.
For Schwarzschild the Kretschmann invariant is exactly $48m^{2}/r^{6}$; under a
formal ramification $r=s^{q}w(s)$ with $w(0)\neq0$ the pole order becomes
exactly $6q$ with the leading coefficient unchanged, so no radial
reparametrization removes the blow-up; the proper time from rest at $R$ to
$r=0$ is $(\pi/2)\sqrt{R^{3}/2M}$, finite --- and infinitesimal, hence still
finite, for infinitesimal $R$; the tidal component $2M(2M-r)/r^{4}$ diverges;
and the exact identity $r^{6}\Kret=48m^{2}$ has no field-valued solution at
$r=0$ in any characteristic-zero field. Evaluating at a nonzero infinitesimal
radius names another point of a punctured domain. It does not supply the
missing endpoint, a continuation law, or an observable rule.

Because the subject is empirical and the danger is overclaim, the report is
built around a firewall. Every substantive statement carries one of three
visible tags: \TE\ an \emph{exact symbolic identity}, which is all that the
three accompanying verification programs establish; \TC\ a \emph{conditional
mathematical theorem}, whose hypotheses are strong and are established for no
physical model; or \TA\ an \emph{assessment}, which is an argument and not a
result. Imported literature carries \TI. Flattening these tiers would
manufacture exactly the overclaim all three manuscripts were written to refuse.

What survives on the positive side is substantial and precisely delimited:
exact curvature and infall identities, a small-angular-momentum crossover with
an exact quadrature, inner and outer Hahn support criteria for a regular core
that are sharp at $\alpha=2/3$, formal and coefficientwise Einstein reduction
theorems with counterexamples showing their hypotheses cannot be dropped,
finite-dimensional surcomplex Hermitian and quantum-shadow theorems with their
conditioning boundary, and an explicit incompatibility between the canonical
integer-part global phase and a normalized scalar derivation. None of it is a
singularity resolution, and none of it is claimed to be.
\end{abstract}
\end{minipage}
\vfill
\begin{minipage}{0.93\textwidth}\small
\textbf{Status.} A merged exposition assembled from three independent
manuscripts, all pinned to \Repo\ at commit \Commit. Prepared with AI
assistance; not refereed; no proof-assistant verification of any statement in
this report. The three companion programs check $45+50+37=132$ finite symbolic
identities and nothing else. No priority claim is made for any proposition,
computation or benchmark. \Cref{phys:app:nonclaims} is the consolidated
register of what is \emph{not} claimed; it is part of the result, not a
disclaimer appended to it. The single sentence this report exists to refuse, in
any form, is that surreal numbers resolve the black-hole singularity.
\end{minipage}
\end{titlepage}
\hypersetup{pageanchor=true}
\pagenumbering{roman}
\begingroup\small\setlength{\parskip}{0pt}
\tableofcontents
\endgroup
\clearpage
\pagenumbering{arabic}
\part{The reading contract}
\label{phys:part:contract}

\section{What this report is}
\label{phys:sec:intro}

\subsection{Three manuscripts, one verdict}
\label{phys:sub:three}

This report is the union of three independently written manuscripts on the
same commission: assess whether surreal and surcomplex scalars can replace the
real and complex numbers in theoretical physics, with black-hole singularities
as the test case. They were prepared separately, they are pinned to the same
repository snapshot, and they reach the same conclusion by the same worked
examples. Each of the three independently rebuilds the connection and the full
Riemann tensor of the general static spherical metric in a computer algebra
system and re-derives the same Schwarzschild quantities from the metric alone.
Two of the three print the same Planck-curvature radius
$r_{\mathrm Q}=48^{1/6}(m\lP^{2})^{1/3}$ with the same disclaimer that it is
dimensional power counting and not a derived threshold.

All three carry the following spine:
\begin{enumerate}[label=\textup{(\roman*)},leftmargin=2.6em]
\item what a black-hole singularity actually is --- coordinate divergence,
curvature divergence, causal geodesic incompleteness and inextendibility at a
named regularity are four different defects;
\item Schwarzschild evaluated at an infinitesimal radius, with the exact
proper-time infall law and the exact tidal Jacobi exponents;
\item a conditional standard-part or coefficientwise reduction of the Einstein
equations, with the observation that its hypotheses fail exactly where a
singularity lives;
\item a regular-core comparison in which the changed \emph{dynamics}, not the
changed scalar field, does the regularizing work;
\item finite-dimensional surcomplex quantum algebra, together with an explicit
list of what does not extend;
\item quantum fields, regulators, and transseries or renormalization examples
showing that field enlargement does not fix a prescription;
\item a comparison with other non-Archimedean and singular methods;
\item a staged development program for the repository;
\item a claim-status ledger and a curvature appendix.
\end{enumerate}
And all three end at the same assessment: a strong exact asymptotic
coefficient language, but extending the scalar field alone does not remove a
singularity, define division by zero, fix ultraviolet dynamics, or supply an
observable rule.

Duplication between the three is very high --- on the order of two thirds. The
curvature and infall computations, the Planck power counting with its identical
$r_{\mathrm Q}$, the regular-core example, the finite-dimensional quantum
section, the transseries and regulator material, the framework comparison and
the research program are common to all three, and each ships its own
independent symbolic re-derivation of the same connection and curvature
($45$, $50$ and $37$ checks respectively). Whatever is shared is printed
\emph{once}. Where the same conclusion is reached by genuinely different
arguments, both arguments are printed and marked as such. Nothing that any one
manuscript reached and the others did not has been dropped, and every honest
limitation any of them states survives, in \cref{phys:app:nonclaims}.

\subsection{What each member contributed}
\label{phys:sub:contrib}

\paragraph{Member 12 --- scalar-extension assessment (the base).}
A $34$-page critical assessment with $25$ references. It is the most complete
of the three and supplies: the pole-persistence proposition invariant under
formal ramification; the boundary-layer warning about differentiating standard
parts; the exact inner and outer Hahn support criteria for the regular core,
sharp at $\alpha=2/3$; the formal Einstein-tensor reduction over
$C^{\infty}(U,\R)[[h]]$; the global-phase section deriving $\ExpEK$ with kernel
$2\pi i\Oz$ together with its proved incompatibility with a normalized
Berarducci--Mantova derivation, and the real-time corollary
$\ExpEK(i\omega t)=1$; a units-and-falsifiability section; the explicit
refutation-by-counterexample of the published surreal-singularity proposal; and
audit, notation and claim-status appendices.

\paragraph{Member 13 --- no-go obstructions.}
A deliberately negative-leaning assessment with a $50$-check suite. It supplies
the three-way taxonomy of replacement proposals kept rigorously separate
throughout; the algebraic endpoint obstruction for $r^{6}\Kret=48m^{2}$ in any
characteristic-zero field; anisotropy, weak singularities and scale invariants
beyond one pole, with the Kasner coefficients; the double-null identities and
the Borel-residue example; the non-multiplicativity and reparametrization
dependence of finite-part extraction; the infinite-period proposition for a
global phase homomorphism; the bounded-frame valuation invariance; the
dominant-balance principle; the argument that the ultraviolet question is
dynamical and not arithmetic; a measurement, recoverability and
empirical-novelty section; and a five-stage development program with explicit
failure conditions the system should \emph{report} rather than paper over.

\paragraph{Member 14 --- black-hole singularity limits.}
A $31$-page assessment with a $37$-check suite. It supplies the
small-angular-momentum crossover with its exact dimensionless quadrature and
uniform-on-bounded-intervals limit; the order-theoretic proof that a leading
pole stays unlimited in \emph{any} ordered extension fixing $\R$; the section
arguing that the classical theory loses control at a positive real curvature
scale before any literal infinitesimal radius; the noncanonical-subtraction
example and the reading of standard part as a measurement rule; the remark that
the regular-core $r^{5}$ term becomes $\abs{x}^{5}$ along a Cartesian line;
Kerr, near-extremal Reissner--Nordstr\"om and oscillatory regimes; the survey of
published surreal-physics proposals; the beyond-all-orders factorial-series
ambiguity; and a research program with explicit success criteria.

\subsection{The merge rule}
\label{phys:sub:mergerule}

A merge is a union, not a selection. Concretely, in this report: a result every
member proves is printed once; a result proved by genuinely different routes
keeps both proofs, marked; anything one member reached that the others did not
survives; every honest limitation survives; and nothing is strengthened beyond
what a source proves. The last clause is the operative one for a physics
report, and it is what \cref{phys:sec:firewall} exists to enforce.

\section{The claim-status firewall}
\label{phys:sec:firewall}

\subsection{Why this section comes before any physics}
\label{phys:sub:whyfirewall}

This is the only report in its collection whose subject is not mathematics.
Its entire value is the firewall between what is proved and what is proposed.
The characteristic failure mode of surreal-physics writing is not a wrong
calculation; it is a correct calculation presented at the wrong epistemic
grade, so that a formal identity is read as a statement about general
relativity, or an evaluation at an infinitesimal radius is read as a physical
regime. Both misreadings are available from perfectly true sentences. The only
defence is to tag the grade of every statement and never let a tag be inferred
from context.

\begin{convention}[The three claim tiers, and the provenance marker]
\label{phys:conv:tiers}
Every substantive statement in this report carries exactly one of the
following tags, in its theorem header or at the head of its paragraph.
\begin{description}[style=nextline,leftmargin=0pt,labelsep=0pt]
\item[\TE\ Exact symbolic identity.]
A finite algebraic or differential identity in ordinary real or complex
symbols, checked exactly --- not numerically, not to a tolerance --- by at
least one of the three accompanying programs, or derived here from such an
identity by finite algebra. \emph{This tier is the entire content of what the
verification programs establish.} It includes the curvature components, the
Schwarzschild and Kasner invariants, the infall and Jacobi exponents, the
crossover integrand factor, the regular-core central limits, the finite-part
shifts, the regulator integrals, and $r_{\mathrm Q}$ as an algebraic
consequence of setting $\lP^{4}\Kret$ to order one.
\item[\TC\ Conditional mathematical theorem.]
A theorem with an ordinary mathematical proof in the text, under hypotheses
stated in full. The hypotheses are strong, and \emph{they are established for
no physical model}. This tier includes the formal and coefficientwise Einstein
reductions, the quantum-shadow results, the finite Hermitian spectral theorem,
the fine-topology collapse, the pole-persistence and endpoint obstructions, the
Hahn support criteria, and the phase and period propositions. Nothing in this
tier is machine-checked. A \TC\ statement about a formal coefficient ring is
not a statement about general relativity, and this report never lets one be
read as the other.
\item[\TA\ Assessment.]
An argument rather than a result: a judgement about what a proposal does or
does not accomplish, a modelling choice declared as such, a power-counting
estimate, a research proposal, or a negative claim scoped as a \emph{missing
implication} rather than an impossibility theorem. Assessments are the reason
this report exists; they are also the tier on which nothing may be built.
\end{description}
Separately, \TI\ marks a result \emph{imported} from the published literature,
used with its own hypotheses intact and neither reproved nor strengthened here.
\TI\ is a provenance marker, not a fourth claim tier: an imported result is
evidence about the literature, not a claim of this report.
\end{convention}

\begin{warning}[The two inferences this report refuses]
\label{phys:warn:refuse}
First: the Einstein reduction theorems of \cref{phys:sec:reduction} are
theorems about \emph{formal coefficient equations under stated conditions}.
They are not theorems about general relativity, and in particular they say
nothing whatever about a metric whose leading term degenerates --- which is
exactly the situation at a singularity. Second: ``Schwarzschild at an
infinitesimal radius'' is a \emph{mathematical extension} of a classical
formula, not a physical regime; \cref{phys:sec:control} is placed immediately
after the computation, rather than a dozen sections later, precisely so that
this cannot be missed.
\end{warning}

\subsection{The ledger}
\label{phys:sub:ledger}

\Cref{phys:tab:ledger} is the report's claim-status ledger. It is duplicated in
the accompanying \texttt{README.md} so that the contract is stated before
anyone opens the article.

\begin{longtable}{@{}L{0.28\textwidth}L{0.09\textwidth}L{0.55\textwidth}@{}}
\caption{Claim-status ledger. \TE\ exact symbolic identity; \TC\ conditional
theorem proved in the text; \TA\ assessment; \TI\ imported.}
\label{phys:tab:ledger}\\
\toprule
Claim & Tier & Essential scope, and what is \emph{not} claimed \\
\midrule
\endfirsthead
\toprule
Claim & Tier & Essential scope, and what is \emph{not} claimed \\
\midrule
\endhead
\bottomrule
\endfoot
General static-spherical $\Kret[f]$, $\Scal$, $G^{t}{}_{t}$ &
\TE &
Derived on a nondegenerate patch; independently rebuilt from the metric by all
three programs. A rational identity in $f$ and $r$, nothing more. \\
\addlinespace
$\Kret=48m^{2}/r^{6}$, $\Kret(2m)=3/(4m^{4})$, vacuum $\Ric=0$ &
\TE &
Classical general relativity, recomputed and symbolically checked here. Not new
physics. \\
\addlinespace
$m^{4}\Kret=48\eps^{-6}$ at $r=m\eps$ &
\TE &
An exact Hahn-field value. \emph{Not} a statement that a curvature observable
has become a limited physical magnitude. \\
\addlinespace
Infall $r^{3}=\tfrac92 m u^{2}$, $\Kret=64/(27u^{4})$, tidal $4/(9u^{2})$,
$-2/(9u^{2})$, Jacobi exponents $4/3,-\tfrac13,\tfrac23,\tfrac13$ &
\TE &
Exact for the $E=1$ radial timelike geodesic in the stated frame convention.
Generic rates only; special initial data change them. \\
\addlinespace
Proper time from rest at $R$ to $r=0$ equals $(\pi/2)\sqrt{R^{3}/2M}$ &
\TE &
Finite for finite $R,M$; infinitesimal, hence still finite, for infinitesimal
$R$. Infall terminates. \\
\addlinespace
Crossover quadrature and its integrand factor; $\Phi$ endpoint asymptotics &
\TE &
Exact identity plus a limit uniform only on bounded positive intervals. \\
\addlinespace
Hayward $\Kret(0)=24/\ell^{4}$, $\rho(0)=3/(8\pi\ell^{2})$,
$r_{\mathrm c}=(2m\ell^{2})^{1/3}$ &
\TE &
For fixed positive \emph{real} $\ell$. $C^{2}$-level curvature regularity for
this ansatz only; not $C^{\infty}$, not geodesic completeness, not stability,
not a microphysical Lagrangian. \\
\addlinespace
Planck-curvature radius $r_{\mathrm Q}$; two-scale
exponent $4\beta-6\alpha$ &
\TE/\TA &
The algebra is exact \TE; its reading as a threshold is \TA\ dimensional power
counting, not a derived universal radius and not a prediction of any
quantum-gravity theory. \\
\addlinespace
Finite-part shifts $\FP_{\eps}$ versus $\FP_{\theta}$; $a\mapsto a-c$ &
\TE &
Exact. Shows a prescription dependence; proves no renormalization theorem. \\
\addlinespace
Pole persistence under scalar extension and ramification &
\TC &
Hahn-field valuation statement; and, by a second independent route, an
order-field statement invariant under any ordered extension fixing $\R$.
Deliberately \emph{not} about real limits. \\
\addlinespace
No field-valued Schwarzschild endpoint &
\TC &
Conditional on preserving the same invariant identity at the endpoint; not a
no-go for a modified metric, a different equation, nonclassical observables, or
a geometry with no zero-areal-radius event. \\
\addlinespace
Set-sized collapse of the full fine topology &
\TC &
Set-sized subsets only. Explicitly \emph{not} a claim that the proper class is
discrete, and \emph{not} applicable to the intrinsic valuation topology of a
set-sized Hahn field. \\
\addlinespace
Formal Einstein reduction over $C^{\infty}(U)[[h]]$ &
\TC &
Requires a nondegenerate smooth leading metric on one common domain. Singular
perturbations are outside it; a counterexample shows the hypotheses cannot be
dropped. Says nothing about general relativity at a singularity. \\
\addlinespace
Coefficientwise Hahn reduction and the shadow corollary &
\TC &
Strictly positive global support, common smooth domain, Neumann support lemma.
Five named escape routes are left open. \\
\addlinespace
Inner/outer strong-summability criteria, sharp at $\alpha=2/3$ &
\TC &
A criterion on Hahn-field valuation data, not on real convergence. It detects
misuse of an expansion; it invents no solution. \\
\addlinespace
Finite Hermitian spectral theorem; finite quantum shadow &
\TC &
Ordinary finite dimensions and step counts; standard-part readout adopted as a
modelling rule. Conditioning covered only when the outcome probability has
nonzero standard part. Not an infinite-dimensional spectral theorem. \\
\addlinespace
$\ExpEK$ with kernel $2\pi i\Oz$; phase/derivation incompatibility;
$\ExpEK(i\omega t)=1$; infinite periods &
\TC &
Statements about one specific commutative construction and one specific
normalized derivation. They do not assert that no global surcomplex exponential
exists --- one is constructed here. \\
\addlinespace
Surreal derivations, transseries embeddings, resurgent integration &
\TI &
Imported with their stated function-class restrictions. No general Einstein PDE
transfer is claimed or available. \\
\addlinespace
Penrose incompleteness, Raychaudhuri focusing, $C^{0}$ Cauchy-horizon
stability, effective-field-theory corrections, quantum probes &
\TI &
Imported with hypotheses intact. Used to say what a singularity theorem does
and does not conclude. \\
\addlinespace
The conservative architecture: ordinary spacetime $+$ controlled Hahn or
transseries coefficients $+$ explicit real-shadow map &
\TA &
Called the most defensible starting point. \emph{Not} a derived necessity. \\
\addlinespace
Standard part as the apparatus readout rule &
\TA &
An intelligible correspondence rule and an additional modelling assumption.
Not forced by the field axioms; undefined on an unlimited curvature. \\
\addlinespace
The staged program, benchmarks, projects and proposed module names &
\TA &
Proposals, not descriptions of existing modules and not results. \\
\addlinespace
A physical resolution of black-hole singularities &
--- &
\textbf{Not supplied by scalar extension, and not claimed anywhere in this
report.} \\
\end{longtable}

\section{The question that must actually be answered}
\label{phys:sec:question}

\subsection{Four different projects, three different replacement proposals}
\label{phys:sub:projects}

\TA\ Replacing the real numbers can mean four importantly different things.
An \emph{algebraic extension} lets a metric, a matrix or a curvature expression
take values in a larger field while keeping the same formulas. An
\emph{asymptotic extension} interprets a small parameter as an infinitesimal
and encodes a family of approximations in a Hahn series or a transseries. A
\emph{nonstandard reformulation} enlarges a structured mathematical universe,
retaining specified transfer and internality principles and recovering ordinary
quantities by a standard-part construction. A \emph{new physical theory}
changes the allowed states, dynamics, geometry or observational rules.

These projects are not equivalent. A useful reformulation may make proofs or
computations easier without changing a single prediction. An algebraic
extension can make an expression meaningful at a new nonzero argument without
making it meaningful at its old singular argument. A genuinely new theory may
use surreal numbers, but its physical content comes from the whole theory, not
from the name of its field.

\TA\ Cutting the same question the other way gives three concrete replacement
proposals, which this report keeps rigorously separate throughout.
\begin{description}[style=nextline,leftmargin=0pt,labelsep=0pt]
\item[(P1) Algebraic scalar extension.]
Replace real or complex coefficients in finite algebraic expressions by
elements of an ordered extension field or its complexification. Matrix
identities, polynomial equations and rational curvature formulas remain
meaningful wherever their denominators are nonzero. Mathematically
straightforward; not yet a theory of time evolution.
\item[(P2) Fields with non-Archimedean coefficients on ordinary spacetime.]
Keep an ordinary real spacetime domain and let fields have Hahn-series or
transseries coefficients, with coefficientwise spacetime derivatives and
explicit support conditions. This is a workable formal deformation theory and
is the principal constructive framework developed below.
\item[(P3) A genuinely non-Archimedean spacetime.]
Replace coordinate domains, curves, causal structure, integration, observables
and possibly probability by new objects over a surreal field. A much larger
proposal. A choice of field alone determines none of its topology, differential
geometry, global existence theorems or empirical interpretation.
\end{description}
These proposals interact, but a result in one does not silently become a
theorem in another. In particular, an expression evaluated at an infinitesimal
radius is not automatically a solution on a larger spacetime containing the
missing endpoint. Section numbers below state which proposal is in force.

\subsection{An operational hierarchy of success}
\label{phys:sub:hierarchy}

\TA\ It is useful to distinguish five increasingly demanding achievements.
\begin{enumerate}[label=\textup{(\arabic*)},leftmargin=2.4em]
\item A divergent quantity has a well-defined generalized numerical value.
\item An asymptotic solution has controlled algebraic and differential meaning.
\item A specified class of observers has well-defined observables through a
previously singular regime.
\item The spacetime or its replacement has a suitable global evolution and a
completeness or continuation theorem.
\item The model recovers tested physics and makes a discriminating, stable,
empirically meaningful prediction.
\end{enumerate}
Surreal arithmetic can directly assist the first two. The remaining steps need
additional mathematics and physics. Failure to keep these levels separate is
the main source of exaggerated claims about singularity resolution.

The same distinction can be phrased as a standard of evidence. A useful
\emph{representation} faithfully encodes specified equations and scale
relations, with computable support or remainder certificates. An \emph{analytic
model} additionally explains which formal expressions correspond to actual
solutions, with which domains and boundary data. A \emph{physical singularity
resolution} requires still more: an explicit dynamical law, a well-defined
continuation or replacement of the singular regime, a stated regularity and
measurement criterion, and physically interpretable predictions.

\subsection{What a proposed resolution must deliver: a four-question test}
\label{phys:sub:fourquestions}

\TA\ Concretely, require an explicit answer to four questions.
\begin{enumerate}[label=\textup{(Q\arabic*)},leftmargin=3.1em]
\item What states or events replace the classical endpoint?
\item What equations govern propagation there?
\item Which observables remain meaningful, and in what sense?
\item Does the replacement determine a continuation uniquely from admissible
data, or does it introduce new boundary conditions?
\end{enumerate}
The answers may involve nonclassical geometry or quantum states rather than a
smooth metric; that is allowed, and no hypothesis of this report excludes it.
What is not enough is to retain every classical formula, replace the word
``divergent'' by ``surreal-valued'', and regard the missing choices as settled.

\subsection{How strong is the affirmative case?}
\label{phys:sub:affirmative}

\TI\ There is substantial mathematics behind the affirmative case. Berarducci
and Mantova constructed compatible differential structures on the surreal
numbers and related them to transseries \cite{phys:BM,phys:BMgerms};
Aschenbrenner, van den Dries and van der Hoeven placed the resulting
differential field in a universal $H$-field framework with transseries and
Hardy-field embeddings \cite{phys:ADH}; van den Dries and Ehrlich established
surreal exponentiation and its elementary-extension properties
\cite{phys:vdDE}; Ehrlich and Kaplan constructed surreal ordered exponential
fields together with trigonometric and surcomplex exponential structures
\cite{phys:EK,phys:EKErratum}; and Costin and Ehrlich developed extension and
integration operators for a substantial, explicitly delimited class of
resurgent functions, along with obstructions to substantially more
unrestricted constructive extensions \cite{phys:CE}. These are serious tools
for asymptotic analysis, not informal manipulations of infinity. Their
availability makes it reasonable to investigate reduced gravitational equations
and multiscale models by surreal methods.

\TA\ The more ambitious statement --- that physical spacetime or quantum
amplitudes fundamentally use $\No$ or $\SC$ rather than $\R$ or $\C$ --- does
not follow. No predictive black-hole completion is established by the
repository or by any surreal-physics proposal examined here. That is an
assessment of the sources inspected, not a theorem excluding all future
surreal-valued theories.

\subsection{What the repository provides, and what it does not}
\label{phys:sub:repo}

All three manuscripts inspected \Repo\ at the pinned snapshot
\begin{center}\small\Commit.\end{center}
\TA\ Between them they read the root \texttt{README.md}, the documentation
catalogue \texttt{docs/README.md}, the surcomplex analysis report's README and
the trigonometry report's README. The catalogue describes fifteen research
packages --- six on surreal numbers, seven on surcomplex mathematics, two on
foundations and computation --- and labels them AI-assisted drafts rather than
refereed or machine-checked foundations
\cite{phys:RepoRoot,phys:RepoDocs,phys:RepoAnalysis,phys:RepoTrig}.

The distinction between the reports and the formalized prerequisites matters.
The root README describes implemented Lean prerequisites covering generic
finite algebra and geometry, polynomial operations and jets, complexification,
and a Hahn summability and standard-part layer with compatibility between
complexification and Hahn coefficients. It states explicitly that these
prerequisites do not yet construct the surreal field or establish its
normal-form bridge to Hahn series, and do not establish all claims in the
reports. It would therefore be wrong both to call the repository unformalized
and to treat a successful Lean build as verification of its analytic reports.
No manuscript ran that build, and neither did this merge.

Three documented features are used below. The analysis report separates strong
Hahn summation from ordinary topological convergence and distinguishes
coefficient rings and function classes --- fixed-domain coefficient functions,
common-domain germs, individually convergent germs with no uniform domain, and
purely formal coefficients. The trigonometry report separates canonical
finite-angle constructions from the additional choices involved in a global
phase, and explicitly declines to select a scalar derivation. The documentation
recommends set-sized coefficient workspaces over unrestricted proper-class
constructions.

A concrete instance of the common-domain warning, taken from the catalogue,
is the series
\begin{equation}\label{phys:eq:nocommondisk}
F(z)=\sum_{n\ge1}\frac{t^{n\eta}}{1-nz},\qquad \eta>0.
\end{equation}
Every coefficient is a holomorphic germ at zero, but the pole at $1/n$ prevents
the family from sharing an ordinary disk of holomorphy. Treating this as a
field on one fixed neighbourhood introduces a false hypothesis. The physics
analogue is to claim a uniform spacetime domain from coefficientwise local
solutions; \cref{phys:sub:commondomain} makes the requirement explicit.

\TA\ None of these ingredients is already a relativistic spacetime, an Einstein
initial-value theory, a quantum state space with infinite-dimensional spectral
theory, or a rule for measuring surreal-valued observables. A suitable
extension of the repository would therefore have to add a \emph{physical
semantics and a differential-geometric layer}, not merely more arithmetic
identities. The missing pieces are coefficient-function algebras, tensor
differentiation, physical interpretation, analytic realization, and an account
of evolution. \Cref{phys:sec:program} turns these into concrete development
targets. Proposed physics work should build on the verified algebraic
interfaces where applicable, and should not cite the repository as an already
formalized theory of surreal general relativity.
\section{The scalar systems and their observational shadows}
\label{phys:sec:scalars}

\subsection{Surreal numbers are a field, not a collection of error symbols}
\label{phys:sub:field}

\TI\ Conway's surreal numbers form a proper-class ordered real-closed field
containing $\R$ and the ordinals. Every individual surreal number has a
set-sized normal form, although the collection of all of them is not a set.
Adjoining $i$ with $i^{2}=-1$ gives the algebraically closed
\begin{equation}\label{phys:eq:SC}
\SC=\{a+ib:a,b\in\No\},\qquad
\overline{a+ib}=a-ib,\qquad
\abs{a+ib}=\sqrt{a^{2}+b^{2}}\in\No_{\ge0}.
\end{equation}
The order lives in $\No$, not in $\SC$. Algebraic closure holds for ordinary
finite-degree polynomials and is not analytic completeness; the modulus is
surreal-valued and is not an ordinary real-valued Banach norm. The
multiplicative and triangle identities are algebraic consequences of real
closedness. These facts, the normal-form description and the exponential
structure are in \cite{phys:Conway,phys:Gonshor,phys:BM,phys:EK,phys:Ehrlich}.

An infinite positive element $\omega$ satisfies $\omega>n$ for every ordinary
natural number $n$, and $\eps=\omega^{-1}$ is positive and smaller than every
positive real. Nevertheless
\begin{equation}\label{phys:eq:nodivzero}
\eps\ne0,\qquad \eps^{-1}=\omega,\qquad 0^{-1}\text{ does not exist}.
\end{equation}
\TC\ No field extension can change the last statement while retaining the
field axioms: if $0x=1$ then $0=1$. Likewise there is no largest positive
surreal and no smallest positive infinitesimal, since $x+1>x$ and
$0<\eps/2<\eps$.

\TA\ An infinitesimal is therefore not automatically a minimal length, a
shortest time step, or an ultraviolet cutoff. Its presence adds finer scales
rather than removing them. Choosing a surreal infinitesimal does not create a
minimum length. Surreal numbers are also commutative and have no nonzero
nilpotents: anticommuting Grassmann variables, noncommuting observables and
operator algebras still require their own constructions. Nothing in this report
is available over a noncommutative base without reproof, and the
surquaternionic material in this repository's \texttt{docs/surquaternions/}
family --- where conjugation is an anti-automorphism, the modulus need not be
multiplicative, and the value group need not be central --- must be treated as
a separate subject. Every result below uses commutativity somewhere: valuation
additivity, the Neumann series for a metric inverse, $\st$ as a ring
homomorphism, the Rayleigh-quotient step in the spectral theorem, Hahn
convolution.

\begin{definition}[Limited, infinitesimal, unlimited]
\label{phys:def:limited}
In an ordered field $F\supseteq\R$, an element $x$ is \emph{limited} if
$\abs x<M$ for some real $M>0$; \emph{infinitesimal} if $\abs x<r$ for every
real $r>0$; and \emph{unlimited} if it is not limited.
\end{definition}

This vocabulary is used deliberately in place of ``finite'' and ``infinite''
to block a specific slide: from ``curvature is now a field element'' to
``curvature is a finite physical magnitude relative to ordinary units.'' An
infinite surreal is an individual field element; it is not thereby a limited
observable. The two phrasings ``finite element'' and ``limited element'' denote
the same set $\Ofin_{F}$ below and are used interchangeably only where no
physical reading is in play.

\subsection{Set-sized Hahn workspaces}
\label{phys:sub:workspace}

For concrete calculations, fix a divisible ordered abelian group $\Gamma$ that
is a \emph{set} and work in
\begin{equation}\label{phys:eq:hahn}
F_{\Gamma}=\R((t^{\Gamma})),\qquad
K_{\Gamma}=\C((t^{\Gamma}))=F_{\Gamma}[i].
\end{equation}
An element is a formal sum $a=\sum_{\gamma\in S}a_{\gamma}t^{\gamma}$ with
$S\subseteq\Gamma$ well ordered. The convention is that $t^{\gamma}$ is
\emph{small} for $\gamma>0$, so the valuation of a nonzero series is
$\val(a)=\min\supp(a)$, with $\val(0)=+\infty$ as a bookkeeping symbol and not
a newly adjoined field element. Then
\begin{equation}\label{phys:eq:valrules}
\val(ab)=\val(a)+\val(b),\qquad
\val(a+b)\ge\min\{\val(a),\val(b)\}.
\end{equation}
\TI\ For divisible $\Gamma$, $F_{\Gamma}$ is real closed and $K_{\Gamma}$ is
algebraically closed. If $\Gamma$ is realized as an additive subgroup of
$\No$, the normal-form realization sends $t^{\gamma}$ to $\omega^{-\gamma}$,
consistently with this convention
\cite{phys:BM,phys:EK,phys:Conway,phys:Gonshor}. The notation
$\omega^{-\gamma}$ denotes Conway's monomial map; for a general surreal
exponent it must not be silently identified with a Gonshor exponential
expression, and exponential or logarithmic closure is extra structure, not a
property of every selected workspace \cite{phys:vdDE,phys:EK}.

\TA\ Every set of surreal coefficients uses only a set of normal-form
exponents, so the divisible additive subgroup they generate is a set-sized
workspace containing that data. This removes the proper-class obstacle from
individual calculations. It proves \emph{no} closure under a chosen
exponential, a scalar derivation, composition, integration or infinite
summation; those obligations must be discharged separately. It also does not
turn a proper-class-indexed field theory into a set-sized one by changing
notation. The precise statement of the localization principle, its three
clauses and the size obstructions that any candidate foundation must respect
belong to the companion merged report \emph{Foundations for Surreal and
Surcomplex Mathematics}, held in
	exttt{docs/foundations-and-computation/foundations/}
of this repository, which proves that every \emph{set} of surcomplex numbers
lies inside one divisible set-sized algebraically closed workspace
$K_{\Gamma}$, that any set of such workspaces lies inside a larger one, and
that a polynomial splitting in one acquires no new roots in a larger
surcomplex extension. This report imports that discipline and does not
re-argue set-theoretic legality; readers weighing proposal (P3) of
\cref{phys:sub:projects} should consult its size obstructions rather than this
report.

For the calculations below, $\Gamma=\Q$ suffices for the Schwarzschild
Puiseux exponents, $\Gamma=\R$ for arbitrary Kasner powers, and
beyond-all-orders exponential scales require an enlarged group or a named
transseries subfield --- see \cref{phys:sub:borel}, where
$\ee^{-1/\eps}\notin\R((\eps^{\Q}))$ is proved by an exponent argument.

\subsection{Strong Hahn summability is a support condition, not convergence}
\label{phys:sub:summability}

\begin{definition}[Strong Hahn summability]
\label{phys:def:strongsum}
A family $(a_{j})_{j\in J}$ of Hahn series is \emph{strongly summable} when the
union of its supports is well ordered \emph{and} each exponent lies in the
support of only finitely many members. Addition is then coefficientwise.
\end{definition}

The second clause is a finitary-multiplicity condition and is not implied by
the first. It does the work in both directions.

\begin{example}[\TE\ What is and is not strongly summable]
\label{phys:ex:summable}
For an infinitesimal $\eps$ of positive valuation,
\begin{equation}\label{phys:eq:geom}
\sum_{n\ge0}\eps^{n}=(1-\eps)^{-1},
\qquad\text{with exact finite remainder}\quad
(1-\eps)^{-1}-\sum_{n=0}^{N}\eps^{n}=\frac{\eps^{N+1}}{1-\eps},
\end{equation}
an exact algebraic identity, \emph{not} a limit of partial sums in any
topology. The family $\sum_{n\ge0}n!\,\eps^{n+1}$ is likewise strongly summable
--- its support lies in the positive integers and each exponent occurs once ---
although the corresponding real power series has zero radius of convergence.
By contrast the ordinary convergent real series $\sum_{n\ge0}2^{-n}$, and the
family of real constants $1/n!$, are \emph{not} strongly summable when viewed
as families of constant Hahn series: exponent $0$ occurs infinitely often. So
$e$ is first defined by ordinary real analysis in the coefficient field and
then embedded; it is not obtained by ignoring the support rule. Finally, a
countable family each of whose members equals the same positive infinitesimal
$\eps$ is not strongly summable, which is why countably additive
non-Archimedean probability is not free (\cref{phys:sub:probability}).
\end{example}

\subsection{Finite elements, infinitesimals, and standard part}
\label{phys:sub:st}

For an ordered field $F\supseteq\R$ write
\[
\Ofin_{F}=\{x\in F:\abs x<n\text{ for some }n\in\N\},\qquad
\mm_{F}=\{x\in F:\abs x<1/n\text{ for every }n\ge1\}.
\]

\begin{proposition}[\TC\ Standard part]
\label{phys:prop:st}
Each $x\in\Ofin_{F}$ has a unique decomposition $x=r+\delta$ with $r\in\R$ and
$\delta\in\mm_{F}$. The map $\st(x)=r$ is an order-preserving ring
homomorphism with kernel $\mm_{F}$, inducing $\Ofin_{F}/\mm_{F}\cong\R$. For
$K=F[i]$ it extends componentwise to a homomorphism
$\Ofin_{F}+i\Ofin_{F}\to\C$ commuting with conjugation.
\end{proposition}
\begin{proof}
The set $A=\{a\in\R:a\le x\}$ is nonempty and bounded above; let
$r=\sup_{\R}A$. For every positive real $h$ the supremum property gives
$r-h<x<r+h$, with harmless weak inequalities at an endpoint, so $x-r$ is
infinitesimal. Two distinct reals cannot differ by an infinitesimal, giving
uniqueness. Sums of infinitesimals and products of an infinitesimal with a
limited element are infinitesimal; expanding $(r+\delta)(s+\eta)$ proves the
homomorphism assertion. Order preservation and the complex assertion are
immediate. The proof uses completeness of $\R$ and nothing about $F$ beyond
its order.
\end{proof}

For the Hahn field \eqref{phys:eq:hahn}, $\Ofin$ consists of the series of
nonnegative valuation, $\mm$ of those of strictly positive valuation, and
$\st$ extracts the coefficient of $t^{0}$. Both descriptions are used below;
they agree on $\Ofin$.

\begin{proposition}[\TC\ Two elementary obstructions]
\label{phys:prop:elementary}
No field extension makes division by zero valid. Moreover $\st$ cannot be
extended to a unital ring homomorphism from an entire non-Archimedean field to
$\R$ or $\C$ annihilating the nonzero infinitesimals.
\end{proposition}
\begin{proof}
In a field $0b=0$ for every $b$, so $0b=1$ is impossible. For the second
assertion pick a nonzero infinitesimal $\eps$; a multiplicative extension
would give $1=\st(\eps\eps^{-1})=\st(\eps)\st(\eps^{-1})=0$.
\end{proof}

So $1/\eps$ is meaningful and unlimited while $1/0$ remains undefined; and
\TA\ ``read the observable by taking the standard part'' is a rule for
\emph{limited} values only, unless another measurement prescription is
supplied. Infinite curvature has no finite real standard part merely because it
has acquired a surreal name. An observational rule based on standard part is a
possible modelling choice, not a consequence of the field axioms; it discards
infinitesimal corrections by design, and a theory that instead measures them
must specify what apparatus, outcomes and probability rule do so.
\Cref{phys:sec:measurement} returns to this.

\subsection{Dimensionful constants are not infinitesimals by being small}
\label{phys:sub:units}

\TA\ A physical length is a numerical value \emph{and} a unit. The statements
$r/m=\eps$ and $\lP/m=\eps^{\beta}$ are well-defined scale comparisons; the
statement ``the radius is $\omega^{-1}$'', with no unit or reference length, is
not. Multiplication by a fixed positive real conversion factor preserves
infinitesimality, but a complete model must also say which dimensional
constants and dimensionless ratios are held fixed in its asymptotic family.
Throughout this report a \emph{dimensionless ratio}, never a dimensionful
quantity, is what may be declared infinitesimal.

The Planck length is a positive real length in conventional physics; its small
numerical value in metres does not make it infinitesimal. If $m$ and $\lP$ are
fixed nonzero real lengths and $\eps$ is a genuine infinitesimal then
$m\eps\ll\lP$ --- not merely ``somewhat below an astrophysical radius''. A
formal parameter representing a real-family limit $\lP/m\to0$ is useful, but
identifying that parameter with a physical constant is a separate interpretive
step. This distinction becomes load-bearing in \cref{phys:sec:control}.

\TA\ There is also no physical reason, from the field alone, to prefer
$\omega^{-1}$ to $\omega^{-2}$ or to $\exp(-\omega)$. Their radically
different orders become useful when dictated by an equation or a matching
condition. Choosing one as a fundamental cutoff without such a principle adds
unexplained structure.

\subsection{Transfer: what one may and may not infer}
\label{phys:sub:transfer}

\TI\ Real-closed-field algebra has first-order transfer in the language of
ordered fields, and richer elementary embeddings preserving restricted analytic
and exponential structure exist in the settings studied by van den Dries,
Ehrlich and Kaplan \cite{phys:vdDE,phys:EK,phys:Ehrlich}. It would be incorrect
to say that surreals admit no transfer.

\TA\ What does not follow from a bare field embedding is a transfer principle
for all real functions, all subsets, Lebesgue measure, Hilbert spaces, or
theorems about solutions of nonlinear partial differential equations. The
language, the expanded structure and the admissible quantifiers matter. A
nonstandard superstructure carrying internal sets and transferred integration
supplies additional data that an ordered field alone does not specify; Loeb's
construction of an ordinary measure from nonstandard measure data is a useful
illustration of how much surrounding structure such a result needs
\cite{phys:Loeb}. Thus ``the formula still makes sense'' is a much weaker
statement than ``the existence, uniqueness and stability theorem still holds'',
and it would be equally misleading to say that everything transfers or that
nothing does. The question is always which language, structure and hypotheses
have actually been transported.

\section{Topology and calculus are part of the physical proposal}
\label{phys:sec:calculus}

\subsection{The full fine topology is not ordinary spacetime topology}
\label{phys:sub:fine}

Use the positive surreal-valued modulus to define \emph{fine balls}
$B(a,r)=\{z\in\SC:\abs{z-a}<r\}$, allowing \emph{every} positive surreal
radius $r$; equivalently, take order balls on $\No$ and modulus balls on $\SC$
as the default neighbourhoods on the full proper class. All three manuscripts
prove the following, by the same Conway-cut argument; it is printed once, with
the union of their statements and of their caveats.

\begin{proposition}[\TC\ Set-sized collapse in the full fine topology]
\label{phys:prop:discrete}
Every \emph{set} $A\subseteq\SC$ is closed and discrete in the full fine
topology. Every convergent set-indexed net is eventually equal to its limit.
Every continuous map from an ordinary connected real interval into this
topology is constant. In particular the usual topology of $\R$ is \emph{not}
the subspace topology inherited from the full fine space.
\end{proposition}
\begin{proof}
Fix $x\in\SC$. The nonzero distances $\abs{x-a}$ for $a\in A\setminus\{x\}$
form a \emph{set} of positive surreals, so the Conway cut
$\{0\mid \abs{x-a}\}$ produces a positive $\delta$ smaller than every member of
that set. The ball $B(x,\delta)$ therefore meets $A$ only at $x$ if $x\in A$,
and misses $A$ otherwise; an empty distance set causes no difficulty, since any
positive radius serves. This gives discreteness and, applied to points outside
$A$, closedness. The range of a set-indexed net together with its proposed
limit is a set, so a ball isolating the limit forces eventual equality.
Finally the image of a connected real interval is a set, hence discrete, and
the continuous image of a connected space is connected, so the image is a
single point.
\end{proof}

\begin{warning}[The three things \cref{phys:prop:discrete} does not say]
\label{phys:warn:discrete}
\begin{enumerate}[label=\textup{(\alph*)},leftmargin=2.4em]
\item It does \emph{not} say the proper class $\SC$ is a discrete space. A
singleton is not an open ball, and the proof cannot choose one radius below a
\emph{proper class} of positive distances. The restriction to set-sized subsets
is essential and is stated as a hypothesis, not a convenience.
\item It does \emph{not} apply to the \emph{intrinsic} valuation or order
topology of a set-sized Hahn field. In $\R((t^{\Q}))$ one has $t^{n}\to0$ in
its own valuation topology; after embedding into $\No$ a smaller scale such as
$\omega^{-\omega}$ blocks that convergence. The two ambient topologies must
never be identified. Even in a non-Archimedean ordered field, the embedded
real sequence $1/n$ does not approach zero in the full order topology, because
one fixed positive infinitesimal is below all of its terms.
\item It does \emph{not} rule out coarser topologies or coefficientwise smooth
structures. Four alternatives change the model in four different ways: the
valuation topology on a set-sized field; a coefficientwise topology; the
standard-part topology on $\Ofin$, defined by inverse images of ordinary open
sets under $\st$, which cannot separate two points with the same standard part
and whose natural separated quotient is ordinary space; and a specialized
notion of surreal analytic function. None may be silently substituted for
another in a proof about causal curves or limits.
\end{enumerate}
\end{warning}

\TA\ The constructive consequence adopted below is the fourth option in
\cref{phys:sub:projects}: retain ordinary real spacetime as the parameter
domain and attach formal coefficients to its fields. That is proposal (P2),
and it is chosen precisely so that no claim depends on ordinary real histories
being continuous in the full fine topology.

\subsection{Three derivatives, and a fourth object, that must not be conflated}
\label{phys:sub:derivatives}

This is the report's central hygiene rule. All three manuscripts state a
version of it; their union is the following table, which is referred to
throughout and never restated.

\begin{longtable}{@{}L{0.20\textwidth}L{0.31\textwidth}L{0.41\textwidth}@{}}
\caption{The derivative and exponential ledger: domain, law, and what fails.
No row may be used for another row's object.}
\label{phys:tab:derivatives}\\
\toprule
Object & Domain and normalization & Law it obeys, and what fails \\
\midrule
\endfirsthead
\toprule
Object & Domain and normalization & Law it obeys, and what fails \\
\midrule
\endhead
\bottomrule
\endfoot
Coefficientwise spacetime derivative $\partial_{\mu}$ &
Series $\sum_{\gamma}a_{\gamma}(x)t^{\gamma}$ with all $a_{\gamma}\in
C^{\infty}(U)$ on \emph{one common} ordinary domain $U\subseteq\R^{4}$;
$\partial_{\mu}t^{\gamma}=0$ &
Commutes with the other $\partial_{\nu}$, Leibniz, commutes with $\st$ and with
$\pi_{0}$ by construction. This is the derivative used for every metric
perturbation below. It is \emph{not} a derivative in the asymptotic parameter.
\\
\addlinespace
Parameter derivative $\dd/\dd\eps$ &
$\dd(\eps^{q})/\dd\eps=q\eps^{q-1}$ &
Ordinary calculus in the parameter. Distinct from both neighbours; holding
$t$ constant while claiming to differentiate in $r$ after setting $t=r/m$ is
an error, since then $\partial_{r}=m^{-1}\partial_{t}$. \\
\addlinespace
Scalar derivation $D$ on $\No$ &
Berarducci--Mantova derivation, normalized $D\omega=1$, $D|_{\R}=0$
\cite{phys:BM,phys:ADH} &
Leibniz and strong additivity on $\No$; Liouville closed. Then
$D(\omega^{-1})=-\omega^{-2}=-\eps^{2}$, neither $1$ nor $0$. It
differentiates \emph{elements read as asymptotic functions}; it is not a
physical time derivative, and one scalar derivation does not supply four
commuting coordinate derivatives or a connection. Choosing it by analogy with
time evolution can make a supposedly constant regulator time-dependent. \\
\addlinespace
Derivative of a surreal \emph{function} &
Costin--Ehrlich extensions and integration on specified resurgent classes
\cite{phys:CE} &
Genuine calculus on its stated class, with proved obstructions to broader
constructive extension. A fourth object, not the scalar derivation. \\
\addlinespace
Fine difference-quotient derivative &
A chosen topology on a non-Archimedean field &
Definable, but does not automatically preserve mean-value, compactness or
existence theorems. \\
\addlinespace
Real surreal exponential on $\No$ &
$\exp:\No\to\No_{>0}$ \cite{phys:vdDE} &
Order isomorphism onto the positives; growth and decay only. Defines no
oscillatory phase. \\
\addlinespace
Infinitesimal Hahn exponential $\Emm$ &
$\Emm(h)=\sum_{n\ge0}h^{n}/n!$ for $h\in\mm_{\C}$ &
Strongly summable; a homomorphism $(\mm_{\C},+)\to$ units, injective, and
satisfies the chain rule for a compatible derivation. Defined only on
infinitesimals. \\
\addlinespace
Finite-angle phase $\cis$ &
$\cis(\theta)=\ee^{i\st\theta}\Emm(i(\theta-\st\theta))$ for
$\theta\in\Ofin$ &
Group law with kernel $2\pi\Z$; the ordinary factor comes from complex
analysis, the infinitesimal tail from strong summation. Says nothing about
infinite angles. \\
\addlinespace
Global phase $\ExpEK$ &
Ehrlich--Kaplan integer-part construction with $\Oz$ as the canonical initial
integer part \cite{phys:EK} &
A group homomorphism with kernel $2\pi i\Oz$, extending the ordinary complex
exponential --- but \emph{no} chain rule for any derivation with $D\omega=1$
(\cref{phys:prop:phase}), and $\ExpEK(i\omega t)=1$ for every real $t$
(\cref{phys:cor:constantphase}). \\
\addlinespace
Quaternionic radial exponential &
Surquaternionic manuscripts, \texttt{docs/surquaternions/} &
\emph{Not} an additive-to-multiplicative homomorphism for noncommuting
arguments, and not an everywhere Hahn-summed power series. Named here only so
that $\ExpEK$'s homomorphism property is not carried across; no result of this
report may be lifted to a noncommutative base without reproof. \\
\end{longtable}

\begin{warning}[The rows are different objects, not competing accounts]
\label{phys:warn:ledger}
$\ExpEK$ \emph{is} a group homomorphism; the quaternionic radial exponential is
\emph{not}; and a componentwise derivation can admit \emph{no} global
commuting exponential chain rule. All three statements are true, of three
different objects. Flattening them into one symbol ``the surreal exponential''
manufactures a contradiction between correct theorems. Likewise
$D(\omega^{-1})=-\eps^{2}$ contradicts nothing: $\partial_{\mu}$, $\dd/\dd\eps$
and $D$ answer different questions, and $D(t)=0$ for every real scalar $t$
does not mean that a real-coordinate physical field has zero derivative.
\end{warning}

\subsection{Common domains and support conditions are essential}
\label{phys:sub:commondomain}

An expression $A(x)=\sum_{\gamma}a_{\gamma}(x)t^{\gamma}$ is used below only
when the coefficient functions live on a specified common ordinary domain and
the support rules justify the operations being performed. Having an analytic
germ for each coefficient is strictly weaker than having a common
neighbourhood: \eqref{phys:eq:nocommondisk} is the witness. A coefficient
domain that shrinks to nothing cannot support an advertised common spacetime
neighbourhood.

\TA\ Likewise, a Taylor expansion does not determine an arbitrary smooth
function: $\ee^{-1/x^{2}}$ for $x>0$, continued by zero for $x\le0$, has zero
Taylor series at the origin and is not identically zero. A rule that extends
all smooth functions by formal Taylor coefficients necessarily loses such
information. A physical proposal must choose a function class or a richer
extension rule rather than assuming that every classical field extends
canonically.

\subsection{Integration: substantial results, specified scope}
\label{phys:sub:integration}

\TI\ Costin and Ehrlich construct extensions of a specified class of resurgent
real functions to surreal arguments, together with compatible
antidifferentiation and integration; their class includes important
Borel-summable examples and solutions of nonresonant meromorphic systems, and
changes of variables allow finite singular locations to be treated. They also
prove obstructions to substantially more unrestricted constructive extensions
\cite{phys:CE}.

\TA\ This is a positive starting point for physics, and a reduced
gravitational ODE might fall inside such a class --- in which case a surreal
function extension would be far more informative than assigning a formal value
to a divergent term. Membership, the relevant sector or branch and the chosen
integration operator would all have to be checked. These results are not a
general Einstein PDE existence theorem, and they concern derivations on
\emph{functions}, a different object from the scalar derivation of
\cref{phys:tab:derivatives}.

In the coefficientwise framework, integration over a fixed ordinary domain is
instead defined by integrating each coefficient, provided all coefficient
integrals exist and the resulting family remains admissible:
\begin{equation}\label{phys:eq:coeffint}
\int_{U}\Bigl(\sum_{\gamma}L_{\gamma}(x)t^{\gamma}\Bigr)\dd^{4}x
:=\sum_{\gamma}\Bigl(\int_{U}L_{\gamma}(x)\dd^{4}x\Bigr)t^{\gamma},
\end{equation}
whose support is contained in the original well-ordered support. This is not
the full fine-topological limit of Riemann sums, and it is the only integration
used below.

\subsection{Coordinate and frame discipline}
\label{phys:sub:frames}

\begin{proposition}[\TC\ Bounded-frame valuation invariance]
\label{phys:prop:frame}
For a finite tuple $V$ over a Hahn field put
$\val_{*}(V)=\min_{j}\val(V_{j})$. If a matrix $A$ and its inverse $A^{-1}$
both have limited entries, then $\val_{*}(AV)=\val_{*}(V)$.
\end{proposition}
\begin{proof}
Limited entries give $\val_{*}(AV)\ge\val_{*}(V)$ by \eqref{phys:eq:valrules};
applying the same inequality to $A^{-1}(AV)$ gives the reverse.
\end{proof}

\TA\ An infinite boost, or a transformation of infinitesimal determinant, need
not preserve this criterion, so component valuations are \emph{not}
unrestricted coordinate invariants. Curvature scalars, areal radius, proper
time and tidal eigenvalues in a specified observer frame are the better primary
diagnostics, and a general theory must say which formal coordinate changes
preserve its coefficient algebra and what counts as a physically comparable
observer.

\subsection{Valuations classify balances, not solutions}
\label{phys:sub:balance}

\begin{lemma}[\TC\ Dominant balance: the minimum is attained twice]
\label{phys:lem:balance}
Suppose a dimensionless equation is a finite sum
$\sum_{j=1}^{N}c_{j}\prod_{k=1}^{d}X_{k}^{a_{jk}}=0$ with nonzero coefficients
and integer exponents where the denominators exist. If the valuations of the
terms had a unique smallest value, the corresponding leading coefficient could
not cancel. Hence any solution requires the minimum valuation to be attained
at least twice.
\end{lemma}

\TA\ This gives a concrete dominant-balance algorithm: enumerate competing
leading orders, solve the resulting leading-coefficient equations, then
determine corrections. It can guide asymptotic Einstein or matter equations
once a gauge and ansatz are fixed. It is a \emph{necessary} condition and
explicitly not sufficient: leading coefficients may fail to solve, constraints
may conflict, logarithmic terms may be forced, and an apparently valid
recursion may have no controlled real realization. Valuation arithmetic is
attractive when several corrections compete --- a candidate term proportional
to $\eps^{a}x^{-b}$ changes order according to $a-b\alpha$ when
$x=\eps^{\alpha}$, and a collection of such inequalities can locate balances
and propose rescalings before a costly expansion is attempted --- but the
physical coefficients and admissible operators must still come from an
effective action or another specified dynamical model.
\part{Black holes: the worked examples, and what they settle}
\label{phys:part:bh}

\section{What a black-hole singularity is}
\label{phys:sec:singularity}

\subsection{Four notions that must remain separate}
\label{phys:sub:fournotions}

\TA\ A coordinate formula can diverge even when the underlying geometry is
regular. A curvature tensor can diverge in a freely falling frame even if some
scalar invariants remain bounded, and the absence of a particular divergent
invariant is not a general regularity theorem. A causal geodesic can be
incomplete with no established curvature blow-up. And a spacetime can admit a
continuous metric extension but no extension with the differentiability needed
for the classical field equations. These are four different defects requiring
four different notions of resolution
\cite{phys:Witten,phys:Landsman,phys:Senovilla,phys:Dafermos,phys:DafermosLuk}.

Artificially excising a regular point produces incompleteness with no physical
curvature singularity; demanding maximality or appropriate inextendibility
addresses that separate issue \cite{phys:Senovilla}. So the questions
``does a bad chart merely need replacement'', ``is the spacetime extendible,
and with which differentiability'', ``do curvature scalars or freely falling
tidal components diverge'', and ``can a chosen matter field or quantum
observable continue even when classical geometry does not'' are logically
independent. Answering one settles none of the others.

\TA\ These distinctions survive any serious change of scalars. A metric must be
nondegenerate to define its inverse. Curvature needs derivatives in a specified
sense. A geodesic needs an affine parameter and a solution theory. An extension
needs a larger space, compatible charts or their replacement, and agreement
with the original geometry. Adding $\omega$ to the scalar vocabulary supplies
none of these structures. Replacing coefficients by surreal numbers cannot
supply the missing geometric definitions; without them, ``the geodesic
continues through an infinite value'' has no precise geometric content.

\TA\ The phrase ``where usual real analysis fails'' therefore needs refinement.
Real differential geometry successfully identifies incomplete geodesics,
distinguishes coordinate artefacts from invariant pathology, and estimates
diverging tidal fields. Real analysis often succeeds precisely at proving that
a particular classical solution cannot be extended with a specified
regularity. That is a diagnosis of the solution or the model, not a
contradiction in the real numbers, and describing the problem as a failure of
real arithmetic misidentifies what needs repair.

\subsection{What the singularity theorems actually constrain}
\label{phys:sub:theorems}

\TI\ One standard form of Penrose's theorem assumes global hyperbolicity with
a noncompact Cauchy surface, an appropriate compact trapped surface and the
null convergence condition, and concludes future null geodesic incompleteness
\cite{phys:Penrose,phys:Witten,phys:Landsman}. The conclusion is not that
every possible singularity is a point of infinite density, nor that every
curvature scalar must diverge, nor that a computer has encountered a large
number. Its proof combines local focusing with global geometry.

\TI\ The local null Raychaudhuri equation for a twist-free congruence in four
dimensions gives, under the null convergence condition,
\begin{equation}\label{phys:eq:ray}
\frac{\dd\theta}{\dd\lambda}
=-\tfrac12\theta^{2}-\sigma_{ab}\sigma^{ab}-R_{ab}k^{a}k^{b}
\le-\tfrac12\theta^{2},
\end{equation}
with the shear contraction on the positive screen space
\cite{phys:Raychaudhuri,phys:Witten}. An initially negative expansion is driven
to focusing within affine parameter at most $2/\abs{\theta_{0}}$, as long as
the regular evolution in question exists. The equality model is
\begin{equation}\label{phys:eq:raymodel}
\theta(\lambda)=\frac{\theta_{0}}{1+\tfrac12\theta_{0}\lambda},
\qquad\theta_{0}<0,
\end{equation}
with a pole at $\lambda_{*}=-2/\theta_{0}$.

\begin{example}[\TE\ Evaluating the focusing model at an infinitesimal]
\label{phys:ex:ray}
Evaluating \eqref{phys:eq:raymodel} at $\lambda_{*}-\eps$ replaces a large
negative real expansion by an unlimited negative surreal expansion; at
$\lambda_{*}$ the denominator is still zero. Infinite numbers alone do not
prevent focusing. The algebraic inequality \eqref{phys:eq:ray} can be
represented in an ordered extension field; the global causal and completeness
assumptions cannot be replaced by field arithmetic.
\end{example}

\TA\ A caustic of a congruence is not itself necessarily a spacetime
singularity: the global argument is what turns focusing, under further
assumptions, into incompleteness. A new theory may evade a hypothesis by
modifying the effective energy condition, causal structure, topology,
differentiability, or the spacetime concept. Adopting an ordered extension of
$\R$ does not identify which hypothesis changes and proves no replacement
theorem. Conversely, showing that the classical proof cannot be copied into an
unfamiliar topology --- where, by \cref{phys:prop:discrete}, global
compactness and causal-topological arguments are not preserved --- does not
exhibit a nonsingular physical spacetime.

\subsection{Changing variables does not automatically change the physics}
\label{phys:sub:changevars}

\begin{example}[\TE\ A pole, a projective coordinate, and what each means]
\label{phys:ex:projective}
For $y'=y^{2}$ the solution is $y=(t_{*}-t)^{-1}$. At $t=t_{*}-\eps$ the value
is $\eps^{-1}$; at $t=t_{*}$ the identity $(t_{*}-t)y=1$ has no field-valued
solution at all. The projective coordinate $z=1/y$ obeys $\dot z=-1$ and
crosses zero smoothly, so a compactified state space continues the
\emph{projective} trajectory. But the inverse map fails at $z=0$: if $y$ is the
measured stress, replacing it by $z$ has not made the measured stress
bounded. Which of $y$ and $z$ is the physical observable is additional
information.
\end{example}

\TA\ This guards against two opposite errors. Not every divergence is a
fundamental obstruction, because some are removable by changing the geometric
description. But not every compactification is a physical resolution, because
the regularity of the relevant observables must be checked.

\TE\ Reparametrizing time is likewise not a cure. If a timelike geodesic
reaches the end of its domain at finite proper time $\tau_{*}$, then
$u=-\log(\tau_{*}-\tau)$ moves that endpoint to $u=+\infty$ without changing
the elapsed proper time; $u$ is not the proper-time parameter and in general is
not affine. Incompleteness is not cured by relabelling the clock. A surreal
value for $u$ may encode an extraordinarily late stage of that
reparametrization; the missing proper-time continuation is a separate problem.
A compactification, a blow-up, a reciprocal field and a geodesic completion are
four different operations.

\section{Schwarzschild at an infinitesimal radius}
\label{phys:sec:schwarzschild}

\subsection{Units and conventions}
\label{phys:sub:units2}

Signature $(-,+,+,+)$, geometrized units $G=c=1$. The mass parameter $m=GM/c^{2}$
has dimensions of length; proper time is expressed in length units, so
$\tau=c\,\tau_{\text{physical}}$ and $[\Kret]=\text{length}^{-4}$. Per
\cref{phys:sub:units}, a dimensionless ratio is what may be declared
infinitesimal.

\subsection{The horizon is already regular over the reals}
\label{phys:sub:horizon}

The Schwarzschild metric in its familiar static chart is
\begin{equation}\label{phys:eq:schwarz}
\dd s^{2}=-f(r)\dd t^{2}+f(r)^{-1}\dd r^{2}+r^{2}\dd\Omega^{2},
\qquad f(r)=1-\frac{2m}{r},
\end{equation}
with $\dd\Omega^{2}=\dd\theta^{2}+\sin^{2}\theta\,\dd\phi^{2}$. Introduce the
advanced coordinate $v=t+r_{*}$ with $\dd r_{*}/\dd r=f^{-1}$; substitution
gives
\begin{equation}\label{phys:eq:ef}
\dd s^{2}=-f(r)\dd v^{2}+2\,\dd v\,\dd r+r^{2}\dd\Omega^{2}.
\end{equation}

\begin{computation}[\TE\ The horizon needs no new scalars]
\label{phys:comp:horizon}
The $(v,r)$ block of \eqref{phys:eq:ef} has determinant $-1$ even at $r=2m$,
and away from the ordinary angular-coordinate poles the full determinant is
$-r^{4}\sin^{2}\theta$, nonzero there. Hence the divergence of $g_{rr}$ at the
horizon in \eqref{phys:eq:schwarz} is removed by an ordinary \emph{real}
coordinate change. No generalized number is needed. All three companion
programs check this determinant.
\end{computation}

\TE\ Inside the horizon $g^{ab}(\partial_{a}r)(\partial_{b}r)=f(r)<0$, so the
gradient of the areal radius is timelike. \TA\ The future singular boundary is
therefore not analogous to an obstacle at the centre of a room that a traveller
could circle around, and allowing a complex-valued radius does not preserve the
real Lorentzian problem.

\subsection{The invariant calculation}
\label{phys:sub:invariant}

\begin{computation}[\TE\ General static-spherical curvature]
\label{phys:comp:Kgeneral}
For a metric of the form \eqref{phys:eq:schwarz} with general $f$,
\begin{equation}\label{phys:eq:Kgeneral}
\Kret[f]:=R_{abcd}R^{abcd}
=(f'')^{2}+4\Bigl(\frac{f'}{r}\Bigr)^{2}
+4\Bigl(\frac{1-f}{r^{2}}\Bigr)^{2},
\end{equation}
and, contracting before squaring,
\begin{equation}\label{phys:eq:Rgeneral}
\Scal=-f''-\frac{4f'}{r}+\frac{2(1-f)}{r^{2}},
\qquad
G^{t}{}_{t}=G^{r}{}_{r}=\frac{rf'+f-1}{r^{2}};
\end{equation}
for $f=1-2M(r)/r$ the last becomes $G^{t}{}_{t}=-2M'(r)/r^{2}$.
\Cref{phys:app:curvature} derives these by an orthonormal coframe.
\emph{Each of the three companion programs rebuilds the Christoffel symbols,
the full four-index Riemann tensor, the Ricci tensor, the scalar curvature and
the mixed Einstein component from the metric alone, with $f$ an unevaluated
symbolic function, and contracts them to confirm
\eqref{phys:eq:Kgeneral}--\eqref{phys:eq:Rgeneral}.} None of them substitutes
into an assumed answer.
\end{computation}

\begin{computation}[\TE\ Schwarzschild]
\label{phys:comp:Kschwarz}
Substituting $f=1-2m/r$ into \eqref{phys:eq:Kgeneral} gives
\begin{equation}\label{phys:eq:Kschwarz}
\Kret=\frac{48m^{2}}{r^{6}},\qquad
\Kret(2m)=\frac{3}{4m^{4}},
\end{equation}
with the three contributions of \eqref{phys:eq:Kgeneral} equal to
$16m^{2}/r^{6}$ each. For $r>0$ Schwarzschild is vacuum: all sixteen Ricci
components and the Ricci scalar vanish, and $G^{t}{}_{t}=0$. The divergent
invariant is built from the full Riemann tensor, equivalently the Weyl tensor
in this vacuum region.
\end{computation}

\begin{warning}[\TA\ The divergence is not an infinite matter density]
\label{phys:warn:weyl}
Because $\Ric=0$ for $r>0$, it is misleading to explain this example as an
infinite ordinary matter density that could be replaced by a surreal density.
The vanishing of the Ricci tensor in this vacuum solution does not contradict
the nonzero, divergent Weyl curvature. Whatever a replacement theory does, it
is not merely renaming a density.
\end{warning}

\begin{computation}[\TE\ Evaluation at an infinitesimal radius]
\label{phys:comp:Keps}
Let $x=r/m$, so the dimensionless invariant is $m^{4}\Kret=48x^{-6}$; with
$x=\eps>0$ a positive infinitesimal,
\begin{equation}\label{phys:eq:Keps}
m^{4}\Kret=48\eps^{-6}.
\end{equation}
Equivalently, with an ordinary reference length $L$, $m=\mu L$ and $r=L\eps$,
\begin{equation}\label{phys:eq:KepsL}
L^{4}\Kret=48\mu^{2}\eps^{-6}.
\end{equation}
This is an exact and genuinely useful non-Archimedean description: it records
both the power and the coefficient of the divergence, and the compared quantity
is dimensionless. The value can be added, multiplied, compared and inverted; it
dominates $m^{-4}\eps^{-5}$ and is dominated by $m^{-4}\eps^{-7}$, information
that a bare symbol $+\infty$ destroys. The valuation of $m^{4}\Kret$ is exactly
the blow-up exponent.
\end{computation}

\begin{warning}[\TA\ What \cref{phys:comp:Keps} does and does not say]
\label{phys:warn:Keps}
The value \eqref{phys:eq:Keps} is unlimited: larger than every positive real
and with no finite standard part. It exceeds every real curvature bound in
these units, and the physics of tidal exposure is not made finite by
reclassifying it as a legitimate scalar. The radius here is the \emph{areal}
radius, defined by the area $4\pi r^{2}$ of the symmetry spheres, so no change
of coordinates alters the scalar value at a fixed event; one may alter the
notion of event or replace the spacetime model, but that is more than an
algebraic relabelling. At $r=0$ the defining expression still has a pole.
\end{warning}

\subsection{Poles persist: two independent routes}
\label{phys:sub:pole}

Two of the three manuscripts prove that the pole survives, by arguments with
different hypotheses and different strengths. Both are printed.

\begin{proposition}[\TC\ Route 1, valuation-theoretic: poles survive scalar
extension and ramification]
\label{phys:prop:pole}
Let $p\ge1$ be an integer, let $u(r)\in\R[[r]]$ have nonzero constant term $c$,
and put $f(r)=r^{-p}u(r)$. Evaluate the series at a positive infinitesimal
$\eps$ in a Hahn field, using strong summation. Then
\[
\val f(\eps)=-p\,\val\eps,
\qquad
f(\eps)=c\,\eps^{-p}(1+\delta),\quad \delta\in\mm_{F},
\]
so $f(\eps)$ is unlimited. Moreover, under a formal change of radial variable
$r=s^{q}w(s)$ with $q\ge1$ an integer and $w(0)\ne0$, the pole order becomes
$pq$ and does not disappear.
\end{proposition}
\begin{proof}
Strong summation of $u(\eps)$ is valid by the positive-valuation power-series
rule, and its constant term is $c$, so $u(\eps)=c(1+\delta)$ with
$\val\delta>0$ unless $\delta=0$. Multiplying by $\eps^{-p}$ gives the stated
leading term and valuation. For the change of variable,
\[
f(s^{q}w(s))=s^{-pq}\,w(s)^{-p}\,u(s^{q}w(s)),
\]
and the factor following $s^{-pq}$ is a unit with nonzero constant term
$w(0)^{-p}c$, which cannot cancel the pole.
\end{proof}

\begin{proposition}[\TC\ Route 2, order-theoretic: a leading pole stays
unlimited in every ordered extension]
\label{phys:prop:unlimited}
Let $F\supseteq\R$ be an ordered field, $\eps>0$ infinitesimal, and $p>0$ such
that $\eps^{p}$ is defined and infinitesimal in $F$. If
$a\in\R\setminus\{0\}$ and $u$ is infinitesimal, then
$I=a\eps^{-p}(1+u)$ is unlimited; and \emph{no} ordered extension of $F$ fixing
$\R$ makes that same element limited.
\end{proposition}
\begin{proof}
Since $\abs u<1/2$ we get $\abs I>\abs a\eps^{-p}/2$, which exceeds every
positive real bound because $\eps^{p}$ is infinitesimal. Every inequality with
a real bound is preserved by an ordered extension.
\end{proof}

\begin{remark}[Why both routes are kept]
\label{phys:rem:tworoutes}
They are not the same theorem. \Cref{phys:prop:pole} is a statement about
Hahn-field valuation data and is what yields \emph{ramification invariance}:
no radial reparametrization of the stated form removes the blow-up.
\Cref{phys:prop:unlimited} is a statement about the order alone and is what
yields \emph{extension invariance}: passing from a small ordered asymptotic
field to all of $\No$ cannot turn the same curvature observable into a limited
one. Neither implies the other, and the two together are the precise content of
``a larger field does not help''. Subtracting terms, rescaling the measured
observable by an infinitesimal, or changing the metric can change the
conclusion --- but each introduces additional structure or a new physical
claim.
\end{remark}

\begin{computation}[\TE\ Ramification invariance, checked]
\label{phys:comp:ramification}
\Cref{phys:prop:pole} applies to \eqref{phys:eq:Kschwarz} with $p=6$ and
$c=48m^{2}$. Substituting $r=s^{q}(1+s)$ into $48m^{2}/r^{6}$ and expanding
gives a pole of order exactly $6q$ with leading coefficient exactly $48m^{2}$,
unchanged, for $q=1,2,3$. This was verified directly from the metric, by
building the Christoffel symbols and the Riemann tensor and contracting: the
Kretschmann scalar is exactly $48M^{2}/r^{6}$ and the substituted pole order
and leading coefficient are as stated. \emph{No radial reparametrization of
this form removes the blow-up.}
\end{computation}

\begin{corollary}[\TC\ The real-analytic geometric consequence]
\label{phys:cor:noext}
If a dimensionless scalar curvature invariant tends to infinity along a curve,
there is no extension along that curve to an ordinary nondegenerate $C^{2}$
metric at an endpoint at which that invariant is continuous and finite:
continuity would give a finite limit.
\end{corollary}
\status{A restricted and useful obstruction. \TA\ It is \emph{not} a claim that
all conceivable generalized geometries have been ruled out, and it says nothing
about a degenerate limiting metric, a weaker regularity class, or a geometry
with no such endpoint event.}

\subsection{The algebraic endpoint obstruction}
\label{phys:sub:endpoint}

Two of the three manuscripts prove the following by the same one-line argument;
it is printed once.

\begin{proposition}[\TC\ No field-valued Schwarzschild endpoint]
\label{phys:prop:noendpoint}
On the Schwarzschild region,
\begin{equation}\label{phys:eq:invariantidentity}
r^{6}\Kret=48m^{2}.
\end{equation}
Let $E$ be any characteristic-zero field containing a nonzero $m$. There is no
assignment of a value $\Kret_{0}\in E$ at $r=0$ extending
\eqref{phys:eq:invariantidentity} to that endpoint. In particular, adjoining
surreal or surcomplex values does not extend this invariant identity through
$r=0$: a regular extension \emph{preserving this relation} at its endpoint is
impossible in any field, not merely in $\R$.
\end{proposition}
\begin{proof}
At $r=0$ the left side of \eqref{phys:eq:invariantidentity} is zero for every
field element $\Kret_{0}$, including one unlimited relative to $\R$. The right
side is nonzero.
\end{proof}
\status{Deliberately a precise, limited statement, and explicitly a
\emph{conditional} obstruction. \TA\ It does not forbid a modified metric, a
different field equation, a nonclassical notion of observable, a quantum
replacement of the endpoint, or a geometry that has no zero-areal-radius event
at all. Identifying an endpoint by an adjoined compactification symbol would
change the algebra and would still require a metric and an evolution law there.
Setting $r=\eps$ merely avoids $r=0$; it does not provide the missing
assignment.}

\subsection{An unlimited number is not an added regular event}
\label{phys:sub:notanevent}

\TA\ Assigning a nonzero infinitesimal radius describes another point in a
punctured extension of the coordinate domain. It does not add the missing
endpoint. There is also no ``last positive radius'' before zero, since
$\eps/2$ is smaller. A proposed rule that replaces zero by one chosen
infinitesimal has introduced a \emph{regulator}, and must explain the choice,
its physical units, the resulting equations, and the invariance of observations
under changing that choice.

\TA\ Allowing a metric to have unlimited coefficients might be part of a
broader theory. In that case the theory must define regularity for those
coefficients and the behaviour of clocks, rods and tidal observables there. It
cannot infer regularity from the fact that a coefficient belongs to $\No$.

\section{The classical theory loses control before a literal infinitesimal scale}
\label{phys:sec:control}

\begin{warning}[Placement is deliberate]
\label{phys:warn:placement}
This section follows \cref{phys:sec:schwarzschild} immediately, and not a
dozen sections later, because the preceding computation is the one a reader is
most likely to misread as a physical regime. One of the three manuscripts makes
the present argument its headline; another makes it inside an
effective-field-theory section; the third defers it. The buried version invites
exactly the misreading this report exists to prevent, so the argument is
promoted here to the status the first manuscript gives it.
\end{warning}

\subsection{A positive real high-curvature scale is encountered first}
\label{phys:sub:rq}

\TI\ Low-energy quantum gravity can be treated as an effective field theory;
its controlled use depends on suitable separation from the ultraviolet scale,
not on an assumption that the classical equations hold at arbitrarily high
curvature \cite{phys:Donoghue,phys:DonoghueReview}. The coefficients, field
content and state are physical inputs, not consequences of choosing $\No$.

\begin{computation}[\TE\ The Planck-curvature radius]
\label{phys:comp:rq}
With $\lP=\sqrt{\hbar G/c^{3}}$, the dimensionless indicator
$\mathcal Q=\lP^{4}\Kret=48m^{2}\lP^{4}/r^{6}$ on Schwarzschild. Setting
$\mathcal Q\sim1$ gives
\begin{equation}\label{phys:eq:rq}
r_{\mathrm Q}=48^{1/6}\bigl(m\lP^{2}\bigr)^{1/3}.
\end{equation}
For $m\gg\lP$,
\begin{equation}\label{phys:eq:rqratios}
\frac{r_{\mathrm Q}}{\lP}=48^{1/6}\Bigl(\frac{m}{\lP}\Bigr)^{1/3}\gg1,
\qquad
\frac{r_{\mathrm Q}}{2m}=\frac{48^{1/6}}{2}\Bigl(\frac{\lP}{m}\Bigr)^{2/3}\ll1.
\end{equation}
\end{computation}
\status{The algebra of \eqref{phys:eq:rq}--\eqref{phys:eq:rqratios} is exact
and is checked by two of the three programs. \TA\ Its \emph{reading} is
dimensional power counting, not a derived universal threshold, not a sharply
defined transition radius, and not an exact prediction of any quantum-gravity
theory. Derivatives, matter scales, the quantum state and the coefficients of
the effective action can introduce other relevant criteria and can invalidate a
simple curvature criterion sooner.}

\TA\ Two consequences follow, and they are the point of this section.

First, a macroscopic areal radius deep inside a large black hole can coexist
with Planck-scale curvature: by \eqref{phys:eq:rqratios} the radius at which
$\mathcal Q\sim1$ is far above the Planck length while far below the horizon.
``Small radius'' and ``high curvature'' cannot be identified without
specifying the mass. A black-hole horizon, a Planck-curvature regime and a
classical singular boundary are three distinct things, and low horizon
curvature for a large black hole is entirely compatible with high interior
curvature.

Second --- and this is the load-bearing statement --- fix ordinary positive
real $L$ and $\lP/L$. A literal positive infinitesimal $\eps$ satisfies
$\eps<\lP/L$, and indeed lies below \emph{every} positive real cutoff ratio.
Evaluating the unchanged classical solution at $r=L\eps$ therefore explores a
regime beyond any such fixed cutoff. Equivalently, if $r=m\eps$ with fixed real
$m,\lP>0$ then
\begin{equation}\label{phys:eq:Qunlimited}
\mathcal Q=48\Bigl(\frac{\lP}{m}\Bigr)^{4}\eps^{-6},
\end{equation}
which is unlimited. A surreal-valued computation can faithfully report that the
expansion parameter is large; it cannot thereby justify neglecting the
higher-order terms that made the classical approximation unreliable. The
computation of \cref{phys:sec:schwarzschild} is therefore a \emph{mathematical
extension} of a classical formula, not a physical regime.

\subsection{Two-parameter power counting becomes exact valuation arithmetic}
\label{phys:sub:twoscale}

The three manuscripts present two versions of the same device, with different
parametrizations. Both are printed because their independent variables differ.

\begin{computation}[\TE\ Version A: radius and Planck length scaled against
the mass]
\label{phys:comp:twoscaleA}
Let $r/m=\eps^{\alpha}$ and $\lP/m=\eps^{\beta}$ with $\alpha,\beta>0$ real.
Then
\begin{equation}\label{phys:eq:EFTvaluation}
\lP^{4}\Kret=48\,\eps^{4\beta-6\alpha},
\end{equation}
which is infinitesimal if $2\beta>3\alpha$, of order one if $2\beta=3\alpha$,
and unlimited if $2\beta<3\alpha$.
\end{computation}

\begin{computation}[\TE\ Version B: mass scaled against the Planck length]
\label{phys:comp:twoscaleB}
Let $m=\lP\eps^{-a}$ and $r=m\eps^{b}$ with $a,b>0$. Then
\begin{equation}\label{phys:eq:twoscale}
\mathcal Q=\lP^{4}\Kret=48\,\eps^{4a-6b},
\end{equation}
sub-Planckian for $b<2a/3$, of order one at $b=2a/3$, super-Planckian for
$b>2a/3$ --- even though $r/m$ is infinitesimal in \emph{every} case.
\end{computation}

\TA\ Version B does not contradict \eqref{phys:eq:Qunlimited}: it changes an
explicit assumption, namely that the mass is held at a fixed real value. That
is precisely the merit of the device. ``Infinitesimally far inside relative to
the mass scale'' does not determine the curvature regime when the mass is
simultaneously scaled, and a valuation representation makes the competing
assumptions visible and mechanically checkable. \TA\ The ratio assignments
describe a \emph{family of models}; they do \emph{not} assert that a fixed
measured Planck length is literally an infinitesimal. Physical application
still requires choosing a finite large mass and estimating how accurately the
formal family describes it.

\subsection{Hawking modes are a different question from the central singularity}
\label{phys:sub:hawkingmodes}

\TA\ A local mode frequency is measured relative to an observer, for example
through $-k_{a}u^{a}$. Large backward-traced frequencies in horizon
calculations concern redshift, states, propagation and short-distance dynamics;
they are not the same invariant issue as $\Kret\to\infty$ at $r=0$.
\TI\ Unruh and Sch\"utzhold derive conditions under which modified
high-frequency dispersion still yields the Hawking effect to leading order, and
also exhibit ways those conditions can fail \cite{phys:Unruh}. \TA\ A
surcomplex symbol for an unlimited phase or frequency selects none of the
dispersion law, the state or the boundary condition used in such a calculation.
It can be a bookkeeping variable inside a specified model; it is not a
substitute for those inputs.

\section{Proper-time infall and the physical tidal field}
\label{phys:sec:infall}

\subsection{An exact radial geodesic}
\label{phys:sub:radial}

\begin{computation}[\TE\ The $E=1$ radial infall]
\label{phys:comp:infall}
For a radial timelike Schwarzschild geodesic with conserved specific energy
$E=f\,\dd t/\dd\tau$, four-velocity normalization gives
$(\dd r/\dd\tau)^{2}=E^{2}-f(r)$. Choose infall from rest at infinity, $E=1$:
\begin{equation}\label{phys:eq:infall}
\Bigl(\frac{\dd r}{\dd\tau}\Bigr)^{2}=\frac{2m}{r},
\qquad
\frac{\dd r}{\dd\tau}=-\sqrt{\frac{2m}{r}}.
\end{equation}
With $u=\tau_{*}-\tau>0$ the remaining proper time,
\begin{equation}\label{phys:eq:properr}
u=\frac{2r^{3/2}}{3\sqrt{2m}},\qquad
r^{3}=\frac92 m u^{2},\qquad
r=\Bigl(\frac{9m}{2}\Bigr)^{1/3}u^{2/3},\qquad
\frac{m}{r^{3}}=\frac{2}{9u^{2}},\qquad
\Kret=\frac{64}{27u^{4}}.
\end{equation}
The cancellation of the mass in the proper-time curvature is exact for this
geodesic.
\end{computation}

\TA\ The fractional power $r\propto u^{2/3}$ is exactly what a ramified Hahn or
Puiseux workspace represents well, and taking $u=L\eps$ produces exact
infinitesimal distances and unlimited curvature while retaining the relation to
the infaller's clock. That is an improvement in \emph{representation}. It is
not an improvement in the boundedness of \eqref{phys:eq:properr}, and it does
not make the remaining proper time to the classical endpoint infinite.

\begin{computation}[\TE\ The infall terminates at finite proper time ---
and at infinitesimal proper time for infinitesimal radius]
\label{phys:comp:finiteproper}
For radial infall from rest at areal radius $R$, the proper time to reach
$r=0$ is
\begin{equation}\label{phys:eq:taufinite}
\tau=\int_{0}^{R}\frac{\dd r}{\sqrt{2M\bigl(r^{-1}-R^{-1}\bigr)}}
=\frac{\sqrt2\,\pi R^{3/2}}{4\sqrt M}
=\frac{\pi}{2}\sqrt{\frac{R^{3}}{2M}},
\end{equation}
finite for finite $R,M$. This was verified independently by evaluating the
quadrature in closed form. Moreover, for \emph{infinitesimal} $R$ the value
\eqref{phys:eq:taufinite} has $\val\tau=\tfrac32\val R>0$, so it is
infinitesimal --- and hence still limited, indeed still \emph{finite}. The
infall terminates either way.
\end{computation}

\begin{warning}[\TA\ The one place a reader may hope for relief, and why there
is none]
\label{phys:warn:tau}
A tempting thought is that at an infinitesimal radius the remaining proper time
becomes ``infinitely short'' in a way that makes the endpoint unreachable.
\Cref{phys:comp:finiteproper} refutes it exactly: the proper time is
infinitesimal, and an infinitesimal is limited. It is not $+\infty$, it is not
unreachable, and no ordered extension makes it so. The clock runs out.
\end{warning}

\subsection{Geodesic deviation resolves the scale hierarchy}
\label{phys:sub:tides}

\begin{computation}[\TE\ Tidal equations and Jacobi exponents]
\label{phys:comp:jacobi}
In a parallel-propagated radial freely falling orthonormal frame, with the
convention in which the radial relative acceleration is stretching,
\begin{equation}\label{phys:eq:jacobi}
\frac{\dd^{2}\xi_{r}}{\dd\tau^{2}}=\frac{2m}{r^{3}}\xi_{r}
=\frac{4}{9u^{2}}\xi_{r},
\qquad
\frac{\dd^{2}\xi_{\perp}}{\dd\tau^{2}}=-\frac{m}{r^{3}}\xi_{\perp}
=-\frac{2}{9u^{2}}\xi_{\perp}.
\end{equation}
Since $\dd^{2}/\dd\tau^{2}=\dd^{2}/\dd u^{2}$, the ansatz $\xi=u^{a}$ gives the
indicial equations $a(a-1)=\tfrac49$ and $a(a-1)=-\tfrac29$, hence
\begin{align}
\xi_{r}&=A\,u^{4/3}+B\,u^{-1/3},\label{phys:eq:radmodes}\\
\xi_{\perp}&=C\,u^{2/3}+D\,u^{1/3}.\label{phys:eq:transmodes}
\end{align}
For a generic diagonal family of three independent Jacobi separations the
radial direction stretches like $u^{-1/3}$ while the two transverse directions
shrink like $u^{1/3}$, so the associated volume scales like $u^{1/3}$ and tends
to zero. Independently, the tidal component
$R^{r}{}_{trt}=2M(2M-r)/r^{4}$ diverges as $r\to0$ and is unlimited at
infinitesimal $r$.
\end{computation}
\status{Exact for this geodesic family in the stated frame convention. \TA\ The
displayed rates are \emph{generic}: special initial data can remove a leading
term and change them, though that only relocates the dominant power and never
removes the divergence of the curvature itself. A radial boost preserves the
principal tidal coefficients, which is why \eqref{phys:eq:jacobi} holds in the
freely falling frame and not only in the static one.}

\TA\ The surreal description retains all these scales simultaneously: for
$u/m=\eps$ it records an unlimited tidal coefficient, an unlimited radial
stretching mode and infinitesimal transverse separations, all at once, and it
connects the singular regime to an observer's clock and to neighbouring freely
falling trajectories. That is more informative than ``infinite curvature is now
allowed''. Calling these legitimate numbers is mathematically coherent. Calling
them bounded physical forces, or a nonsingular passage, is not. A surreal
evaluation supplies no surviving detector, restores no nonzero volume and
determines no propagation through the endpoint.

\subsection{Why a formal bounce is not a geodesic extension}
\label{phys:sub:bounce}

\TA\ One can write $r\propto\abs u^{2/3}$ on both sides of $u=0$, and the
algebraic relation $r^{3}=(9m/2)u^{2}$ admits a formal other branch. This does
not produce a smooth Schwarzschild continuation through the singularity: the
first derivative of $r$ is unbounded there, the curvature still behaves as
$u^{-4}$, and the original metric and geodesic equation are undefined at the
joining event. The relation alone omits the differentiability, curvature,
orientation and matching information a physical extension needs. Matching two
formulas across an omitted point supplies neither the missing geometry nor a
unique law for outgoing matter.

\begin{computation}[\TE\ A smoother parameter is not a regular metric]
\label{phys:comp:smoother}
Write $u=L\lambda^{3}$ with a fixed real length $L$, so $r\propto\lambda^{2}$
and the fractional power in the curve disappears. But
$\dd\tau/\dd\lambda=-3L\lambda^{2}$ vanishes at zero, and the proper-time
metric on the curve pulls back to
\begin{equation}\label{phys:eq:degenerate}
-\dd\tau^{2}=-9L^{2}\lambda^{4}\dd\lambda^{2},
\end{equation}
which degenerates there. This is not a nonsingular coordinate transformation at
the endpoint.
\end{computation}
\status{\TA\ Resolving a ramified formula, or blowing up an algebraic variety,
does not resolve a nondegenerate Lorentzian metric. Algebraic desingularization
remains a useful organizational tool provided the distinction stays explicit.}

\TA\ There is no assertion here that a quantum-gravitational bounce is
impossible. Such a bounce would need a different dynamical or geometric
account of the transition. Surreal methods might describe its asymptotic
incoming and outgoing regions; the transition law is precisely the physical
information the classical formulas do not contain.

\section{A constructive multiscale application: small angular momentum}
\label{phys:sec:crossover}

The $u^{-4}$ curvature law of \cref{phys:comp:infall} must not be universalized
to every Schwarzschild geodesic. Its dependence on angular momentum is the
clearest example in this report of something an ordered asymptotic field
genuinely organizes.

\begin{computation}[\TE\ Nonradial infall changes the proper-time exponent]
\label{phys:comp:nonradial}
Restrict by spherical symmetry to the equatorial plane and let
$j=r^{2}\dd\phi/\dd\tau$ be the conserved specific angular momentum. For
finite $E$ and $j\ne0$,
\begin{equation}\label{phys:eq:radialgeneral}
\dot r^{\,2}=E^{2}-\Bigl(1-\frac{2m}{r}\Bigr)\Bigl(1+\frac{j^{2}}{r^{2}}\Bigr)
=\frac{2mj^{2}}{r^{3}}+\frac{2m}{r}-\frac{j^{2}}{r^{2}}+E^{2}-1.
\end{equation}
The first term dominates as $r\downarrow0$ with $j$ fixed and nonzero, and
integrating the leading balance gives
\begin{equation}\label{phys:eq:nonradial}
r\sim\Bigl[\tfrac52\sqrt{2m}\,\abs j\,s\Bigr]^{2/5},
\qquad
\Kret\sim 48m^{2}\Bigl[\tfrac52\sqrt{2m}\,\abs j\Bigr]^{-12/5}s^{-12/5},
\end{equation}
where $s$ is the remaining proper time. The divergence remains; its power in
remaining proper time changes from $-4$ to $-12/5$.
\end{computation}

\TA\ This is not a change in the curvature as a function of areal radius. It is
a change in how the worldline approaches the same invariantly characterized
region.

\begin{computation}[\TE\ An exact dimensionless crossover quadrature]
\label{phys:comp:crossover}
Let $J=\abs j>0$, set $r=JR$, and define the dimensionless remaining clock
$S=\sqrt{2m}\,J^{-3/2}s$. For the ingoing branch reaching $r=0$, integrating
the exact first integral \eqref{phys:eq:radialgeneral} gives
\begin{equation}\label{phys:eq:crossoverexact}
S=\int_{0}^{R}
\frac{u^{3/2}\,\dd u}
{\sqrt{\,1+u^{2}+\dfrac{J}{2m}\bigl[(E^{2}-1)u^{3}-u\bigr]}},
\end{equation}
the radicand being positive in the sufficiently small interior region under
consideration. \emph{The exact integrand denominator factor in
\eqref{phys:eq:crossoverexact} is symbolically checked.}
\end{computation}

\begin{proposition}[\TC\ The small-$J$ limit, uniform on bounded intervals]
\label{phys:prop:crossover}
As $J/m\to0$ with $E$ fixed and $R$ ranging over any bounded positive interval,
the reciprocal square root in \eqref{phys:eq:crossoverexact} converges
uniformly to that of $1+u^{2}$. Consequently
\begin{equation}\label{phys:eq:Phi}
S\longrightarrow\Phi(R):=\int_{0}^{R}\frac{u^{3/2}}{\sqrt{1+u^{2}}}\dd u,
\end{equation}
with the elementary endpoint asymptotics
\begin{equation}\label{phys:eq:Philimits}
\Phi(R)\sim\tfrac25R^{5/2}\ \ (R\downarrow0),
\qquad
\Phi(R)\sim\tfrac23R^{3/2}\ \ (R\to\infty).
\end{equation}
\end{proposition}
\status{The first asymptotic in \eqref{phys:eq:Philimits} yields the
angular-momentum-dominated law \eqref{phys:eq:nonradial}; the second recovers
the radial law \eqref{phys:eq:properr} in the overlap region $J\ll r\ll m$.
\TA\ The convergence is uniform \emph{only} on bounded positive
$R$-intervals; the large-$R$ behaviour of the limiting function is a
matched-asymptotic statement and is \emph{not} a claim of uniform convergence
on an unbounded interval.}

\begin{computation}[\TE\ Exponent sector tests]
\label{phys:comp:sectors}
With $m=\mu L$, $J=L\eps^{\beta}$ and $s=L\eps^{\alpha}$ for real
$\alpha,\beta>0$, the consistency conditions are
\begin{align}
\alpha<\tfrac32\beta&:\quad
r/L\ \text{has leading exponent}\ \tfrac{2\alpha}{3},
\qquad r\gg J,\label{phys:eq:radialsector}\\
\alpha>\tfrac32\beta&:\quad
r/L\ \text{has leading exponent}\ \tfrac{2(\alpha+\beta)}{5},
\qquad r\ll J.\label{phys:eq:angularsector}
\end{align}
At equality neither monomial may be discarded, and $\Phi$ describes the
transition. Real coefficients of order one can be restored without changing
these exponent tests.
\end{computation}

\TA\ This is what useful surreal-compatible physics mathematics looks like:
retain several small scales, determine which terms dominate, avoid taking
incompatible limits. Note carefully what supplies what. Ordinary real analysis
supplies the actual crossover \emph{function} $\Phi$; Hahn valuations supply a
concise way to \emph{select} its asymptotic regions. The construction offers no
new gravitational law, and it is equally expressible in a suitable transseries
or generalized-series field. Its value must be judged by clarity, support
control and computational reliability --- not by whether the ambient field
contains a proper class of still larger numbers.

\section{Beyond one pole: anisotropy, weak singularities, and other interiors}
\label{phys:sec:beyond}

\subsection{Kasner geometry as an anisotropic benchmark}
\label{phys:sub:kasner}

\begin{computation}[\TE\ Kasner curvature]
\label{phys:comp:kasner}
For the vacuum Kasner metric, with reference scales absorbed into the
coordinates,
\begin{equation}\label{phys:eq:kasner}
\dd s^{2}=-\dd\tau^{2}+\sum_{j=1}^{3}\tau^{2p_{j}}\dd x_{j}^{2},
\qquad \sum_{j}p_{j}=1,\quad\sum_{j}p_{j}^{2}=1,
\end{equation}
a diagonal orthonormal frame gives curvature components with coefficients
proportional to $p_{j}(p_{j}-1)/\tau^{2}$ in the time--space planes and
$p_{j}p_{k}/\tau^{2}$ in the spatial planes, whence
\begin{equation}\label{phys:eq:kasnerK}
\Kret=\frac{4}{\tau^{4}}
\Bigl[\sum_{j}p_{j}^{2}(p_{j}-1)^{2}
+\sum_{j<k}p_{j}^{2}p_{k}^{2}\Bigr].
\end{equation}
For $(p_{1},p_{2},p_{3})=(-\tfrac13,\tfrac23,\tfrac23)$ the coefficient is
$64/27$ --- the same coefficient met in the Schwarzschild proper-time infall
\eqref{phys:eq:properr}. For each of the three permutations of $(1,0,0)$ it
vanishes, corresponding to a flat coordinate example rather than a genuine
curvature divergence. Both constraints, the $64/27$ coefficient and all three
flat exceptional cases are symbolically checked.
\end{computation}

\TA\ The value of a Hahn description here is not ``curvature is infinite''. It
simultaneously retains three anisotropic scales, their products and the orders
of the tidal terms, and it makes \emph{cancellation} visible: identical
exponents alone do not prove a nonzero leading curvature coefficient, as the
flat exceptional case demonstrates. \TI\ Rigorous singular asymptotics already
exist in special Einstein--matter settings; Andersson and Rendall's quiescent
Einstein--scalar and stiff-fluid analysis is a concrete benchmark
\cite{phys:AR}. \TA\ A useful surreal project would reproduce a stated solution
class and its error estimates, not replace those estimates by an appeal to
infinitesimal notation.

\begin{warning}[\TA\ One Kasner expansion is not a theorem about interiors]
\label{phys:warn:kasner}
Generic oscillatory sequences of changing Kasner epochs raise additional
issues, and a single well-ordered expansion must not be assumed to encode every
possible interior evolution. \TI\ BKL-type analyses describe changing
anisotropic regimes, often represented by billiard behaviour, with nontrivial
scope and assumptions \cite{phys:DHN}. \TA\ Hardy fields and most asymptotic
transseries constructions are designed for eventually ordered, nonoscillatory
behaviour; oscillatory phases, repeated changes of dominant balance, or
spatially localized features may require additional variables, sectors, or a
different function category. A Schwarzschild Puiseux law licenses \emph{no}
inference about generic interiors. The real surreal exponential does not by
itself define an unrestricted periodic function at every infinite argument.
\end{warning}

\subsection{Divergent curvature need not have the same integrated effect}
\label{phys:sub:weaktide}

\begin{computation}[\TE\ An integrated tidal criterion]
\label{phys:comp:weaktide}
Consider the schematic tidal magnitude --- \emph{not} a universal formula for
black-hole interiors ---
\begin{equation}\label{phys:eq:weaktide}
\abs{T(s)}\asymp\frac{A}{s^{2}[\log(L_{0}/s)]^{q}},
\qquad 0<s<s_{0}<L_{0}/\ee .
\end{equation}
A natural absolute-integrability diagnostic for the displacement kernel of a
Jacobi equation is the twice-integrated, $s$-weighted quantity
\begin{equation}\label{phys:eq:weakintegral}
\int_{0}^{s_{0}}s\abs{T(s)}\dd s
\asymp A\int_{\log(L_{0}/s_{0})}^{\infty}v^{-q}\dd v,
\end{equation}
by $v=\log(L_{0}/s)$, which is finite exactly when $q>1$. The weight $s$
arises by integrating a positive forcing twice up to the endpoint and reversing
the order of integration. Both the $q\ne1$ primitive and the logarithmic
boundary primitive are symbolically checked.
\end{computation}
\status{Curvature can diverge while this weighted tidal integral converges, so
two profiles both with leading power $s^{-2}$ can differ in integrated effect
because of logarithmic factors. \TA\ This is a diagnostic of an integral
equation and a \emph{model} criterion; it is explicitly not a theorem giving an
extension of a spacetime or its matter fields, and not a substitute for all
hypotheses in a geometric definition of singularity strength.}

\TA\ A surreal or transseries representation preserves that distinction, which
the undifferentiated statement ``the tidal field is an infinite number''
cannot. Integrability of the distortion kernel is information supplied by the
calculus and the equations; it does not follow from making $T(\eps)$ a named
field element. \TI\ Weak mass-inflation-type singularities motivate exactly
this examination; representative charged black-hole analyses separate
diverging curvature from finite tidal distortion
\cite{phys:Bhattacharjee}.

\subsection{Charged and rotating interiors are not Schwarzschild copies}
\label{phys:sub:kerr}

\TI\ Inner-horizon instability and mass inflation in charged or rotating
models introduce blueshift, late-time tails and regularity questions different
from the spacelike Schwarzschild singularity. Dafermos's spherically symmetric
Einstein--Maxwell--scalar analysis shows mass-inflation behaviour coexisting
with a continuous metric extension while stronger regularity is obstructed
\cite{phys:Dafermos}; Dafermos and Luk prove $C^{0}$ stability of the Kerr
Cauchy horizon under prescribed interior-data hypotheses, obtaining continuous
metric extensions across a nontrivial portion of that horizon
\cite{phys:DafermosLuk}.

\TA\ Neither is a theorem that generic black holes have a smooth, unique,
physically harmless interior continuation. What they establish for present
purposes is that even inside real mathematics, asking whether ``the singularity
is removed'' without naming the regularity class is an incomplete question.
Finite or continuous metric components do not imply bounded curvature, a
$C^{2}$ extension, or uniqueness of the continued Einstein evolution. A $C^{0}$
extension does not by itself provide finite curvature or a classical $C^{2}$
solution beyond the boundary, and replacing those components by surreal numbers
settles none of these distinctions.

\TA\ Exponential growth factors together with inverse powers or logarithms of
an advanced-time parameter are natural candidates for transseries descriptions,
and surreal exponentiation and logarithms can retain separated scales that a
truncated power series misses. But a schematic expression involving
$\ee^{\kappa v}v^{-p}$ is \emph{not} a universal theorem about all Kerr
interiors; the model, tail assumptions, gauge and quantities being estimated
must all be specified. Computing a large generalized scalar is not a substitute
for a nonlinear stability argument.

\begin{computation}[\TE\ Near-extremal Reissner--Nordstr\"om bookkeeping]
\label{phys:comp:rn}
For $f(r)=1-2m/r+q^{2}/r^{2}$ the horizon polynomial is $r^{2}-2mr+q^{2}$. In
geometrized units, with $\delta=\sqrt{m^{2}-q^{2}}/m$ for $0<\abs q<m$,
\begin{equation}\label{phys:eq:rn}
r_{\pm}=m(1\pm\delta),\qquad
\kappa_{+}=\frac{r_{+}-r_{-}}{2r_{+}^{2}}=\frac{\delta}{m(1+\delta)^{2}}.
\end{equation}
A symbolic infinitesimal $\delta$ records the splitting of a double root and
the small surface-gravity scale cleanly, and several parameters can keep this
limit separate from a wave-frequency limit or a small perturbation amplitude.
\end{computation}
\status{Algebraic identities and asymptotic bookkeeping. \TA\ Explicitly
\emph{not} a claim that extremality or a Cauchy horizon is physically
harmless.}
\part{The positive theorems, all of them conditional}
\label{phys:part:positive}

\section{Formal and coefficientwise Einstein reduction}
\label{phys:sec:reduction}

\begin{warning}[Read \cref{phys:warn:refuse} again before this section]
\label{phys:warn:reduction}
Everything in this section is tier \TC. The theorems below are statements about
\emph{formal coefficient equations under stated hypotheses}. They are not
theorems about general relativity, they assert nothing about a metric whose
leading term degenerates, and they must never be quoted as showing that a
non-Archimedean deformation of a singular spacetime is well behaved.
\Cref{phys:sec:hypotheses} exhibits the counterexamples that prove their
hypotheses cannot be dropped.
\end{warning}

\subsection{Two conservative geometric frameworks}
\label{phys:sub:frameworks}

The three manuscripts use two genuinely different coefficient algebras, and
both are kept. In each, the set of spacetime events is \emph{not} replaced:
$U$ is an ordinary coordinate domain in a smooth four-manifold, or in $\R^{4}$,
and fields are formal families over it. This is proposal (P2) of
\cref{phys:sub:projects}.

\paragraph{Framework F1: formal power series in one parameter.}
Put
\begin{equation}\label{phys:eq:AU}
\Alg(U)=C^{\infty}(U,\R)[[h]],
\end{equation}
with $h$ a formal parameter killed by the coordinate derivatives. Products are
formal Cauchy products, with only finitely many contributions at each integer
order. Reduction
\[
\pi_{0}:\Alg(U)\to C^{\infty}(U),
\qquad \sum_{n\ge0}f_{n}h^{n}\mapsto f_{0},
\]
is a ring homomorphism and satisfies $\pi_{0}\partial_{\mu}=\partial_{\mu}\pi_{0}$
\emph{by definition}. Formal series can later be evaluated at a positive
infinitesimal in a suitable Hahn workspace, coefficient by coefficient. That
evaluation does not make the map $x\mapsto f(x;h)$ smooth in the full fine
topology; the specified derivative, not an unstated topological
identification, carries the construction.

\paragraph{Framework F2: a Hahn coefficient-function algebra.}
Put
\begin{equation}\label{phys:eq:AGamma}
\Alg_{\Gamma}(U)=\Bigl\{A(x)=\sum_{\gamma\in S}A_{\gamma}(x)t^{\gamma}
: A_{\gamma}\in C^{\infty}(U),\ S\subseteq\Gamma\ \text{well ordered}\Bigr\},
\end{equation}
where $S$ is the \emph{global} coefficient support. Products use Hahn
convolution. This is a ring, not a field: smooth coefficients have zeros.
Write $\Alg^{\ge0}_{\Gamma}(U)$ and $\Alg^{>0}_{\Gamma}(U)$ for the
nonnegative- and positive-support parts; extraction of the exponent-zero
coefficient gives
\[
\st:\Alg^{\ge0}_{\Gamma}(U)\to C^{\infty}(U),
\qquad \st\partial_{\mu}=\partial_{\mu}\st.
\]
Restriction to smaller ordinary domains is allowed and does not enlarge
support. \emph{The distinction between \eqref{phys:eq:AGamma} and a collection
of coefficient germs with shrinking domains is essential}
(\cref{phys:sub:commondomain}); so is the distinction between $\Alg(U)$, where
each order receives finitely many contributions automatically, and
$C^{\infty}(U)((t^{\Gamma}))$, which needs controlled supports and summability
under every operation employed.

\TA\ The notation deliberately separates two questions: whether a formal
equation is \emph{well defined}, and whether its formal solution represents an
actual family of real solutions with controlled remainders. The first is
algebraic. The second is analytic and generally much harder.

\subsection{The reduction theorem in Framework F1}
\label{phys:sub:reduceF1}

\begin{theorem}[{\TC\ Reduction of the Einstein tensor over
$C^{\infty}(U,\R)[[h]]$}]
\label{phys:thm:einstein}
Let $g(h)=g_{0}+hg_{1}+h^{2}g_{2}+\cdots$ be a symmetric $4\times4$ matrix with
entries in $\Alg(U)$, and assume $g_{0}$ is an ordinary smooth nondegenerate
Lorentzian metric on $U$. Define the connection and curvature by the usual
coordinate formulas, differentiating coefficients and holding $h$ fixed. Then
the inverse metric exists in $M_{4}(\Alg(U))$ and reduction modulo $h$
satisfies
\begin{align}
\pi_{0}(g^{-1})&=g_{0}^{-1},\label{phys:eq:ginvreduce}\\
\pi_{0}\bigl(\Ric[g(h)]\bigr)&=\Ric[g_{0}],\label{phys:eq:ricreduce}\\
\pi_{0}\bigl(G[g(h)]\bigr)&=G[g_{0}].\label{phys:eq:einreduce}
\end{align}
If $\Lambda(h)$ and $T(h)$ have the same coefficientwise interpretation and
$G[g(h)]+\Lambda(h)g(h)=8\pi T(h)$ holds coefficientwise on $U$, then
$G[g_{0}]+\Lambda_{0}g_{0}=8\pi T_{0}$. In particular, for real constants
$\Lambda,\kappa$, if $G[g]+\Lambda g=\kappa T$ then
$G[g_{0}]+\Lambda g_{0}=\kappa T_{0}$.
\end{theorem}
\begin{proof}
Write $g=g_{0}(I+H)$ with $H\in hM_{4}(\Alg(U))$. Then
\begin{equation}\label{phys:eq:neumann}
g^{-1}=\Bigl(\sum_{n\ge0}(-H)^{n}\Bigr)g_{0}^{-1},
\end{equation}
and at any specified power of $h$ only finitely many terms contribute; this is
a formal inverse, using no topological convergence, and it proves
\eqref{phys:eq:ginvreduce}. The map $\pi_{0}$ preserves finite sums and
products and commutes with $\partial_{\mu}$, since the latter differentiates
only coefficient functions. Apply it to
\[
\Gamma^{a}{}_{bc}=\tfrac12 g^{ad}
\bigl(\partial_{b}g_{dc}+\partial_{c}g_{db}-\partial_{d}g_{bc}\bigr),
\qquad
R^{a}{}_{bcd}=\partial_{c}\Gamma^{a}{}_{db}-\partial_{d}\Gamma^{a}{}_{cb}
+\Gamma^{a}{}_{ce}\Gamma^{e}{}_{db}-\Gamma^{a}{}_{de}\Gamma^{e}{}_{cb}.
\]
Reduction commutes with the contractions defining the Ricci curvature and the
scalar curvature, hence with
$G_{ab}=R_{ab}-\tfrac12\Scal(g)g_{ab}$. Applying reduction to the formal
Einstein equation gives the last assertions.
\end{proof}
\status{\TC. Requires a nondegenerate smooth leading metric on \emph{one
common} domain. The same argument works for common-domain Hahn coefficient
families when their supports justify multiplication and the inverse expansion,
but those hypotheses must be \emph{stated}, never inferred from pointwise
values. \TA\ Self-contained; not a claim that the repository has formalized
Einstein geometry.}

\subsection{The reduction theorem in Framework F2, and its shadow corollary}
\label{phys:sub:reduceF2}

This is a genuinely different theorem: its coefficient algebra is a Hahn ring
rather than a power-series ring, its perturbation has strictly positive
\emph{global} support rather than positive integer order, and its inverse
needs a support lemma rather than an order count. Both are printed.

\begin{theorem}[\TC\ Standard-part reduction over a Hahn coefficient algebra]
\label{phys:thm:reduction}
Assume a formal metric deformation
\begin{equation}\label{phys:eq:metricdeformation}
g=g_{0}+h,\qquad h\in M_{4}\bigl(\Alg^{>0}_{\Gamma}(U)\bigr),
\end{equation}
on one common smooth domain $U$, with $g_{0}$ an ordinary smooth Lorentzian
metric with $\det g_{0}$ nowhere zero, $h$ globally well supported of strictly
positive order, and $T\in M_{4}(\Alg^{\ge0}_{\Gamma}(U))$. Then
\begin{equation}\label{phys:eq:inversemetric}
g^{-1}=\sum_{n\ge0}\bigl(-g_{0}^{-1}h\bigr)^{n}g_{0}^{-1}
\end{equation}
is a valid matrix Hahn sum, and
\begin{equation}\label{phys:eq:Einsteinreduction}
\st(g^{-1})=g_{0}^{-1},\qquad \st\bigl(G(g)\bigr)=G(g_{0}).
\end{equation}
If $G_{\mu\nu}(g)+\Lambda g_{\mu\nu}=8\pi T_{\mu\nu}$ holds coefficientwise
with $\Lambda\in\R$, its ordinary shadow satisfies
\begin{equation}\label{phys:eq:shadow}
G_{\mu\nu}(g_{0})+\Lambda(g_{0})_{\mu\nu}=8\pi\st(T_{\mu\nu}).
\end{equation}
The same coefficient-zero compatibility holds for the coordinate Riemann and
Ricci tensors and their finite scalar contractions.
\end{theorem}
\begin{proof}
Positive global support and the Neumann support lemma guarantee both a
well-ordered resulting support and finitely many coefficient contributions at
each exponent, so \eqref{phys:eq:inversemetric} is a legitimate element of
$M_{4}(\Alg_{\Gamma}(U))$; only the $n=0$ term has exponent zero, giving the
first identity of \eqref{phys:eq:Einsteinreduction}. Standard part commutes
with addition, with products, with real constants, and with the coefficient
derivatives; applying it to the Christoffel formula yields the Christoffel
formula for $g_{0}$, and applying it to the finite curvature and contraction
formulas gives the second identity. Reducing the formal Einstein equation
coefficientwise gives \eqref{phys:eq:shadow}.
\end{proof}
\status{\TC. At each ordinary point the positive-order perturbation preserves
the signature over the ordered field --- put $g_{0}$ in a real Lorentz-diagonal
frame and use finite elimination; the pivots remain infinitesimal perturbations
of nonzero real pivots and keep their signs. \emph{That statement is pointwise
only}; no compactness and no uniform positive lower bound on $U$ is inferred.}

\begin{corollary}[\TC\ Positive-order corrections cannot repair a singular
shadow]
\label{phys:cor:shadow}
A metric satisfying the hypotheses of \cref{phys:thm:reduction} cannot have a
smooth, nondegenerate ordinary shadow through a point of zero areal radius
while that shadow agrees with the Schwarzschild geometry on a punctured
neighbourhood with $m\ne0$.
\end{corollary}
\begin{proof}
A smooth nondegenerate ordinary metric has finite continuous curvature scalars
locally. On the punctured Schwarzschild region the scalar is $48m^{2}/r^{6}$ by
\cref{phys:comp:Kschwarz}, unbounded as the areal radius tends to zero. That
contradicts such an extension.
\end{proof}
\status{A \emph{conditional} obstruction, not a universal theorem against
non-Archimedean physics. \TA\ It leaves open, explicitly, five escape routes:
a degenerate or nonexistent standard-part metric; negative-order fields;
singular boundary layers; new field equations; and nonclassical observables or
a different notion of spacetime. Any proposed evasion should identify precisely
which hypothesis is being replaced, and how the missing physics is supplied.}

\subsection{Connection, Bianchi identities and gauge}
\label{phys:sub:bianchi}

\TC\ In either framework, with commuting coefficient derivatives, define
\begin{equation}\label{phys:eq:christoffel}
\Gamma^{\alpha}_{\mu\nu}=\tfrac12 g^{\alpha\beta}
\bigl(\partial_{\mu}g_{\beta\nu}+\partial_{\nu}g_{\beta\mu}
-\partial_{\beta}g_{\mu\nu}\bigr),
\end{equation}
and form the curvature, Ricci curvature and Einstein tensor by their usual
finite contractions. The algebraic derivations of torsion-freeness, metric
compatibility and the Bianchi identities --- including the contracted Bianchi
identity --- use only these operations and their identities, so they hold
coefficientwise. At first order one obtains the linearized Einstein operator
together with the corresponding matter perturbation.

\TA\ Ordinary coordinate transformations of $U$ act on the coefficient tensors
in the usual way, and formal diffeomorphisms near the identity can be defined
order by order using Lie derivatives. Neither operation changes the invariant
content merely because the coefficients lie in a larger algebra. Conservation
must either follow from the chosen matter equations or be imposed consistently
with the Bianchi identity; arbitrarily choosing coefficients of $T$ does not
guarantee an admissible matter model. And this is a route to formalization, not
an assertion that any of it is already present in the repository.

\TA\ A practical implementation should treat support conditions as part of
every operation's input contract: inversion requires an invertible leading
term, composition requires the relevant positive-order or domain hypotheses,
and integration requires both a coefficientwise integral and an admissible
combined support. It is unsafe to derive an identity in a formal ring,
interpret it as an ordinary convergent sum, and then apply ordinary
interchange-of-limit theorems without proving a bridge.

\subsection{What the reduction theorems say about new physics}
\label{phys:sub:newphysics}

\TA\ A \emph{regular} infinitesimal deformation of a nondegenerate classical
metric obeys the classical Einstein equation at leading order. So a model
cannot both keep this regular common-domain reduction and obtain an arbitrary
different macroscopic leading equation without changing another assumption ---
the source, the action, or the scaling regime. Read positively, that is a
precise consistency theorem: ordinary general relativity is recovered at
leading order in a regular formal deformation.

\TA\ Read negatively --- and this is the direction that matters here --- the
same theorem limits what the construction can claim. If the leading metric has
a classical singular boundary, the theorem does not extend it through that
boundary. Its hypotheses require a common smooth domain and an invertible
leading metric. The theorems do \emph{not} exclude singular perturbations: a
leading metric may degenerate at a boundary, inverse coefficients may acquire
negative valuation, resummed terms may become dominant, a new leading geometry
may appear, or derivatives may promote a seemingly small correction to leading
order, so that a rescaled inner problem differs substantially from the outer
one. Those are exactly the situations a controlled asymptotic language is for.
They are singular perturbation problems, not counterexamples to the reduction
theorems, whose hypotheses have simply ceased to apply.

\section{Why the hypotheses cannot be dropped}
\label{phys:sec:hypotheses}

\subsection{A boundary layer with uniformly small values and order-one
curvature}
\label{phys:sub:boundarylayer}

\begin{example}[\TE\ The family $b_{\delta}$]
\label{phys:ex:boundarylayer}
For a positive ordinary real parameter $\delta$ put
\begin{equation}\label{phys:eq:boundarylayer}
b_{\delta}(x)=\frac{\delta^{2}}{1+(x/\delta)^{2}}
=\frac{\delta^{4}}{\delta^{2}+x^{2}}.
\end{equation}
For every real $x$ one has $0\le b_{\delta}(x)\le\delta^{2}$, so the family
converges uniformly to zero, and its first derivatives are uniformly of order
$\delta$. Nevertheless
\begin{equation}\label{phys:eq:boundarylayerder}
b_{\delta}''(0)=-2
\end{equation}
for every $\delta$, while the limit function has second derivative zero. Both
\eqref{phys:eq:boundarylayerder} and the derivative scaling are symbolically
checked.
\end{example}

\TA\ Uniform smallness of a metric perturbation therefore does not by itself
imply small curvature: second-derivative control is the missing ingredient, and
this family has first-derivative control without it. Evaluating the same
rational formula at an infinitesimal $\delta$, every pointwise value has zero
standard part, but a coefficientwise or appropriately extended second
derivative at zero still yields $-2$. There is no common formal expansion in
nonnegative powers of $\delta$ with smooth coefficients valid across $x=0$:
away from zero, expansion in $\delta^{2}/x^{2}$ produces coefficients singular
at zero. So the example does not contradict \cref{phys:thm:einstein}; it
demonstrates exactly why its hypotheses cannot be omitted, and why a
\emph{pointwise} standard part is an insufficient basis for differentiating a
limiting metric.

\subsection{A purely real counterexample to interchanging reduction and
differentiation}
\label{phys:sub:oscillatory}

\begin{example}[\TE\ The family $\delta\sin(x/\delta)$]
\label{phys:ex:oscillatory}
The real functions $f_{\delta}(x)=\delta\sin(x/\delta)$, $\delta>0$, converge
uniformly to zero as $\delta\downarrow0$, yet $f_{\delta}'(0)=1$ for every
$\delta$, while the derivative of the limit function is zero.
\end{example}

\TA\ This is a real-family counterexample: no value for $\sin\omega$ is needed
to state it, and no non-Archimedean object appears in it. Its role is to
identify precisely what \cref{phys:thm:einstein} and \cref{phys:thm:reduction}
are buying. In those theorems reduction commutes with differentiation
\emph{by definition of the regular formal algebra}. That commutation does not
hold for arbitrary families followed by a limiting or standard-part operation.
Highly oscillatory and boundary-layer families require additional estimates
before reduction and differentiation may be interchanged.

\subsection{Position-dependent exponents break the coefficient algebra}
\label{phys:sub:varyingexponent}

\TC\ A substantial issue arises when an exponent varies with position:
differentiating $t^{p(x)}$ introduces $p'(x)\log t$, so the constant-exponent
coefficient algebra \eqref{phys:eq:AGamma} is \emph{not} closed under such
expressions. \TA\ Three responses are available, and one of them must be
chosen: restrict the ansatz; enlarge the differential transseries class with a
proof of closure; or use separate local expansions with a justified matching
rule. This is exactly the kind of concrete obstruction a research prototype
should expose rather than paper over.

\subsection{Formal certificates versus analytic certificates}
\label{phys:sub:certificates}

\TA\ A coefficient solution of a formal Einstein equation is not automatically
an actual family of real solutions with that asymptotic expansion. For that one
needs a realization or summation theorem, error bounds, a gauge choice,
constraint propagation, and a well-posed problem in an appropriate function
space. A recursion with nonzero denominators at every formal order need not
have controlled real coefficients or a convergent expansion.

Two different outputs must therefore be distinguished by name.
\begin{description}[style=nextline,leftmargin=0pt,labelsep=0pt]
\item[A formal certificate]
proves support admissibility and vanishing of the residual through a valuation
threshold.
\item[An analytic certificate]
additionally proves approximation or summability and quantifies a real error.
\end{description}
\TA\ The first is valuable without the second, but it must not be advertised as
an actual spacetime solution. A remainder of valuation greater than $N$ bounds
scale in the formal ordering only; a numerical bound at a finite real parameter
needs control of the coefficients and a realization map. Conversely two real
solutions differing by $\ee^{-A/\eps}$ can have identical power-series data
while remaining physically distinct in a process that amplifies that sector.
A representation is most useful when it prevents such information from being
discarded prematurely, and when it \emph{rejects} a product or substitution
whose support condition fails rather than assigning a symbolic infinity and
proceeding.

\subsection{Actions, conservation laws, and measures, with an honest scope}
\label{phys:sub:actions}

\TC\ On a fixed ordinary compact integration domain, a common-support family of
smooth densities integrates coefficientwise by \eqref{phys:eq:coeffint}, and
the support of the result is contained in the original well-ordered support.
Ordinary integration by parts applies separately to every coefficient under its
usual boundary assumptions. For the Einstein--Hilbert density, the inverse
metric exists formally by \eqref{phys:eq:neumann} or
\eqref{phys:eq:inversemetric}, and since $\det g$ has a nonzero ordinary
leading term, the branch of $\sqrt{-\det g}$ around a nondegenerate Lorentzian
$g_{0}$ can be expanded by an admissible positive-order binomial series.
Consequently a formal local action principle, Euler--Lagrange equations and
coefficientwise conservation identities can be developed, and their local
Euler--Lagrange expressions recover the formal equations of
\cref{phys:sec:reduction}.

\TA\ That is a conservative extension of classical variational algebra, and
nothing more. It does not define a nonperturbative path-integral measure,
select boundary conditions at a singularity, or prove convergence of the
expansion, and it does not pretend that a fine-topological Riemann integral
over the whole surreal class has been constructed. A quantum functional
integral is a different task: the field space, the measure or its replacement,
the regularization, the gauge treatment and the limiting procedure all need
definitions, and these are already substantial over $\R$ and $\C$. A
proper-class sum over all surreal-valued histories is \emph{not} licensed by
Hahn summability. Formal sums, finite-dimensional integrals and rigorously
constructed functional measures must remain explicitly different objects.

\section{A regular-core comparison: where new physics actually enters}
\label{phys:sec:core}

\subsection{A real modification already removes the central pole}
\label{phys:sub:hayward}

To separate a change of dynamics from a change of scalars, all three
manuscripts use the static metric function of Hayward's regular-core model
\cite{phys:Hayward}:
\begin{equation}\label{phys:eq:hayward}
f_{\ell}(r)=1-\frac{2mr^{2}}{r^{3}+2m\ell^{2}},\qquad m>0,\ \ell>0,
\end{equation}
inserted into the static spherical form \eqref{phys:eq:schwarz}.

\begin{computation}[\TE\ The central limits, for fixed real $\ell$]
\label{phys:comp:hayward}
At large radius the correction to Schwarzschild in $f$ begins at order
$r^{-4}$. At small radius
\begin{equation}\label{phys:eq:haywardcenter}
f_{\ell}(r)=1-\frac{r^{2}}{\ell^{2}}+\frac{r^{5}}{2m\ell^{4}}+O(r^{8}),
\end{equation}
a de Sitter-like core, and substituting into \eqref{phys:eq:Kgeneral} gives
\begin{equation}\label{phys:eq:haywardK}
\Kret(0)=\frac{24}{\ell^{4}},
\qquad \Scal(0)=\frac{12}{\ell^{2}},
\end{equation}
with the three contributions of \eqref{phys:eq:Kgeneral} tending to
$4/\ell^{4}$, $16/\ell^{4}$ and $4/\ell^{4}$ respectively. The effective mass
function and, via $G^{t}{}_{t}=-2M'_{\ell}(r)/r^{2}$ and $G_{ab}=8\pi T_{ab}$
with $T^{t}{}_{t}=-\rho$, the effective energy density are
\begin{equation}\label{phys:eq:haywardrho}
M_{\ell}(r)=\frac{mr^{3}}{r^{3}+2m\ell^{2}},
\qquad
\rho(r)=\frac{M_{\ell}'(r)}{4\pi r^{2}}
=\frac{3m^{2}\ell^{2}}{2\pi(r^{3}+2m\ell^{2})^{2}},
\qquad \rho(0)=\frac{3}{8\pi\ell^{2}},
\end{equation}
with $p_{r}=-\rho$, so at the centre the stress is of cosmological-constant
form. The crossover radius, from balancing the denominator terms, is
\begin{equation}\label{phys:eq:haywardcrossover}
r_{\mathrm c}=\bigl(2m\ell^{2}\bigr)^{1/3},
\end{equation}
and for fixed $r>0$ letting the real $\ell\to0$ recovers Schwarzschild
pointwise. All of \eqref{phys:eq:haywardcenter}--\eqref{phys:eq:haywardrho},
the expansion through $r^{5}$, the leading outer correction, the mass-function
identity, the Einstein source identity and the $\ell\to0$ limit are
symbolically checked.
\end{computation}

\begin{warning}[\TA\ What does the regularizing, and what ``regular'' means
here]
\label{phys:warn:hayward}
The local de Sitter-type core and the finite limiting curvature are
consequences of the \emph{new metric function} and its effective source, not of
any non-Archimedean arithmetic. Unlike Schwarzschild, the model is not vacuum
in the core. Four qualifications are part of the statement.
\begin{enumerate}[label=\textup{(\alph*)},leftmargin=2.4em]
\item ``Regular core'' here means finite curvature plus the corresponding
classical $C^{2}$-level metric regularity --- \emph{not} analytic or
all-orders smoothness. In fact the $r^{5}$ term in
\eqref{phys:eq:haywardcenter} becomes $\abs x^{5}$ along a Cartesian line
through the centre, so this ansatz must not be advertised as a $C^{\infty}$
Cartesian metric there. The vanishing spherical-coordinate determinant at the
centre must be assessed in Cartesian coordinates; the expansion yields a
nondegenerate $C^{2}$ local metric, and no claim about higher derivatives
follows from the curvature limit alone.
\item Writing $T_{ab}=G_{ab}/8\pi$ is a consistent \emph{effective}
description. It supplies no microphysical matter Lagrangian and establishes no
stability.
\item A finite-curvature core is not a proof of global geodesic completeness,
unique evolution, or physically stable evaporation. The regularity class,
inner-horizon behaviour, global causal assumptions and stress-energy
interpretation must all be examined. \TI\ Classifications of geodesically
complete black-hole scenarios make exactly this larger geometric issue explicit
\cite{phys:CR}.
\item Keeping a finite real $\ell$, possibly motivated by a physical
ultraviolet scale, can give finite central invariants --- but choosing the
scale and the modified source remains a physical hypothesis. Surreal arithmetic
does not determine that $\ell$ is the Planck length, or that a particular
regular-core ansatz is the correct quantum-gravitational dynamics.
\end{enumerate}
\end{warning}

\subsection{Making the core radius infinitesimal reverses the conclusion}
\label{phys:sub:infinitesimalcore}

\begin{computation}[\TE\ Regulator nonuniformity]
\label{phys:comp:regnonuniform}
Set $\ell=m\eps$ with $\eps$ a positive infinitesimal. For every fixed positive
ordinary regulator, \eqref{phys:eq:haywardK} is finite. As a non-Archimedean
value, however,
\begin{equation}\label{phys:eq:coreeps}
m^{4}\Kret(0)=24\eps^{-4},
\end{equation}
which is unlimited; equivalently, with $\ell=L\delta$ and $\delta$
infinitesimal, $L^{4}\Kret(0)=24\delta^{-4}$. The denominator at the centre is
algebraically nonzero, and the central curvature is bounded by no real number.
\end{computation}

\begin{warning}[\TA\ Two statements that are not the same]
\label{phys:warn:twostatements}
\emph{Regular for each fixed real regulator} and \emph{uniformly finite in the
regulator's infinitesimal limit} are different assertions, and the first does
not imply the second. Equally, \emph{a field-valued smooth expression exists}
and \emph{the curvature has a finite ordinary observable shadow} are different
regularity criteria. Taking the regulator infinitesimal makes the
nonuniformity explicit; it does not supply a uniformly finite core, and the
scalar field has not supplied a finite observational curvature bound.
\end{warning}

\TA\ The two real-parameter limits illustrate the same nonuniformity from the
other side. At every fixed $r>0$, letting real $\ell\to0$ recovers
Schwarzschild; near $r^{3}\sim2m\ell^{2}$ that outer approximation ceases to be
uniform and the core must be resolved. A Hahn representation is well suited to
organizing the outer and inner scales. It is the altered metric or action, not
surreal arithmetic, that changes the central behaviour.

\subsection{Hahn summability detects the correct boundary layer}
\label{phys:sub:matching}

This is the sharpest genuinely positive application in the report, and it is
unique to one member. Write
\begin{equation}\label{phys:eq:dimensionlesscore}
x=\frac rm,\qquad \eps=\frac{\ell}{m},
\qquad f(x,\eps)=1-\frac{2x^{2}}{x^{3}+2\eps^{2}},
\end{equation}
treat $\eps$ as a Hahn monomial of positive valuation, and take
$x=\eps^{\alpha}$ with $\alpha>0$ real; the requisite real exponents lie in a
divisible set-sized value group.

\begin{proposition}[\TC\ Outer and inner strong-summability criteria]
\label{phys:prop:matching}
For \eqref{phys:eq:dimensionlesscore}, the outer geometric expansion
\begin{equation}\label{phys:eq:outer}
f=1-\frac2x\sum_{k\ge0}\Bigl(-\frac{2\eps^{2}}{x^{3}}\Bigr)^{k}
\end{equation}
is strongly Hahn summable exactly when $\alpha<2/3$, whereas the inner
expansion
\begin{equation}\label{phys:eq:inner}
f=1-\frac{x^{2}}{\eps^{2}}\sum_{k\ge0}\Bigl(-\frac{x^{3}}{2\eps^{2}}\Bigr)^{k}
\end{equation}
is strongly Hahn summable exactly when $\alpha>2/3$. At $\alpha=2/3$
\emph{neither} displayed geometric family is strongly summable. The exact
transition-scale expression, with $x=\eps^{2/3}X$, is
\begin{equation}\label{phys:eq:transition}
f=1-2\eps^{-2/3}\frac{X^{2}}{X^{3}+2},
\end{equation}
whose denominator is a unit, and which is therefore well defined, for positive
limited $X$ with positive real standard part.
\end{proposition}
\begin{proof}
The ratio in \eqref{phys:eq:outer} has valuation $(2-3\alpha)\val(\eps)$,
positive exactly for $\alpha<2/3$; its powers then have increasing well-ordered
support with finite multiplicity at each exponent, so the geometric identity
applies. For $\alpha>2/3$ the exponents strictly decrease and the family is not
summable. The ratio in \eqref{phys:eq:inner} has valuation
$(3\alpha-2)\val(\eps)$, giving the complementary criterion. At $\alpha=2/3$
all geometric terms contribute at the same exponent, so infinitely many nonzero
contributions land on that exponent and strong summability fails.
Equation~\eqref{phys:eq:transition} follows by direct substitution in the
original rational function, without summing either invalid family.
\end{proof}
\status{\TC. A criterion on Hahn-field data --- valuations of the expansion
ratio --- and \emph{not} on real convergence. Its outer, inner and
transition-scale rational rewrites are symbolically checked.}

\begin{computation}[\TE\ The exact finite remainder underlying the criterion]
\label{phys:comp:remainder}
For every finite $N$,
\begin{equation}\label{phys:eq:remainder}
\frac{1}{1+A}=\sum_{k=0}^{N-1}(-A)^{k}+\frac{(-A)^{N}}{1+A},
\end{equation}
so if $\val A>0$ the remainder has valuation $N\val A$. This identity was
checked exactly for $N=1,\dots,8$. It provides a certificate of asymptotic
order and is a natural target for existing Neumann-series and rational-jet
infrastructure.
\end{computation}

\TA\ The crossover is at $r^{3}\sim2m\ell^{2}$, matching
\eqref{phys:eq:haywardcrossover}; a formally small correction in the outer
region becomes leading in the inner region, and the same sector test can be
phrased in exponents: with $r/L=\eps^{a}$, $\ell/L=\eps^{b}$ and $m/L$ ordinary
and nonzero, comparing $3a$ with $2b$ identifies the outer and inner sectors,
with both terms retained at equality. The gain is real and it is precisely
delimited: \emph{support conditions detect when an approximation has been used
outside its asymptotic regime.} At the transition the failure is a failure of
that expansion, not a singularity of the rational model, and no new physical
solution has been invented by the series algebra.
\part{Quantum structures: what extends, and what does not}
\label{phys:part:quantum}

\section{Finite-dimensional surcomplex quantum algebra}
\label{phys:sec:quantum}

\subsection{The algebraic part works well}
\label{phys:sub:algebra}

Let $F\supseteq\R$ be a real-closed ordered field and $K=F[i]$. Define
\begin{equation}\label{phys:eq:hermform}
\overline{a+ib}=a-ib,\qquad \abs{a+ib}=\sqrt{a^{2}+b^{2}},
\qquad \langle\psi,\phi\rangle=\sum_{j=1}^{n}\overline{\psi_{j}}\phi_{j},
\qquad \lVert\psi\rVert^{2}=\sum_{j=1}^{n}\abs{\psi_{j}}^{2}.
\end{equation}
\TC\ Then $\langle v,v\rangle>0$ for every $v\ne0$, since it is a finite sum of
squares of real and imaginary parts, so \eqref{phys:eq:hermform} is a
positive-definite Hermitian form on $K^{n}$. Positivity and Cauchy--Schwarz
hold by ordinary finite algebra, unitary matrices preserve the norm, and a
normalized vector has nonnegative candidate Born weights
$p_{j}=\abs{\psi_{j}}^{2}$ summing to one in $F$. Orthogonality, tensor
products of finitely many systems and projective-state algebra all make sense.
The same finite algebra supports Lorentzian quadratic forms, their linear
isometry groups, Clifford algebras and finite spinor constructions over
appropriate extensions.

\TA\ What has been preserved is an \emph{algebraic} structure. Compactness of a
group, spectral measures, the topology of a state space and continuum dynamics
are separate issues.

\begin{theorem}[\TC\ Finite Hermitian spectral theorem]
\label{phys:thm:spectral}
Let $F$ be a real-closed ordered field containing $\R$, let $K=F[i]$, and let
$n$ be an ordinary finite integer. Every Hermitian $A=A^{\dagger}\in M_{n}(K)$
has an orthonormal eigenbasis, and all its eigenvalues belong to $F$.
\end{theorem}
\begin{proof}
The characteristic polynomial splits over the algebraically closed $K$, so
there is an eigenvector $v\ne0$ with eigenvalue $\lambda$. Since
$v^{\dagger}v>0$,
\[
\lambda v^{\dagger}v=v^{\dagger}Av=\overline{v^{\dagger}Av}
=\overline{\lambda}\,v^{\dagger}v,
\]
whence $\lambda=\overline{\lambda}\in F$. Normalize $v$ using the positive
square root of $v^{\dagger}v$, which exists because $F$ is real closed. The
orthogonal complement of $v$ is invariant under $A$, because
$\langle v,Aw\rangle=\langle Av,w\rangle=\lambda\langle v,w\rangle$, and the
restriction is Hermitian. Induction on the ordinary finite $n$ completes the
proof.
\end{proof}
\status{All three manuscripts prove this, by the same argument; it is printed
once. The proof requires \emph{no} completeness of $F$. \TA\ It is a finite
algebraic theorem and not an infinite-dimensional spectral theorem for
unbounded operators. Finite quantum kinematics, finite spin matrices and
finite unitary algebra therefore do not force an Archimedean scalar field ---
but that is a scalar-extension result, not a new quantum prediction. An
ordered extension also permits infinitesimal probabilities, whose physical
interpretation is an additional rule; and unlimited eigenvalues can be
represented, but whether they represent physical energies requires an
interpretation and a dynamics.}

\subsection{Finite observations have an ordinary shadow}
\label{phys:sub:shadow}

The three manuscripts state three overlapping versions. Their union is the
following, with the sharp conditioning boundary that only one of them makes
explicit.

\begin{theorem}[\TC\ Reduction of finite quantum protocols]
\label{phys:thm:quantumshadow}
Let all dimensions and step counts be ordinary finite integers.
\begin{enumerate}[label=\textup{(\alph*)},leftmargin=2.4em]
\item If $\psi\in K^{n}$ satisfies $\psi^{\dagger}\psi=1$ then every component
is limited, and $\psi_{0}=\st\psi\in\C^{n}$ satisfies
$\psi_{0}^{\dagger}\psi_{0}=1$; moreover each $p_{a}=\abs{\psi_{a}}^{2}$ is
limited with
\begin{equation}\label{phys:eq:bornshadow}
\sum_{a}\abs{\psi_{0,a}}^{2}=1,\qquad \st(p_{a})=\abs{\psi_{0,a}}^{2}.
\end{equation}
\item Every unitary $U\in M_{n}(K)$ has limited entries and $U_{0}=\st U$ is an
ordinary complex unitary; every orthogonal projection $P=P^{\dagger}=P^{2}$ has
limited entries and $P_{0}=\st P$ is an ordinary orthogonal projection.
\item $\st(\psi^{\dagger}P\psi)=\psi_{0}^{\dagger}P_{0}\psi_{0}$ and
$\st(U\psi)=U_{0}\psi_{0}$; finite observables with limited entries have their
expectation values $\psi^{\dagger}A\psi$ reduced correctly. Consequently the
real standard parts of joint probabilities in a fixed finite sequence of
unitary evolutions and projective measurements agree with the corresponding
ordinary complex calculation.
\item Conditioning on an outcome reduces correctly \emph{provided} the outcome
probability has nonzero standard part.
\end{enumerate}
\end{theorem}
\begin{proof}
Each $\abs{\psi_{j}}^{2}$ is nonnegative and at most the total sum $1$, so each
component is limited; apply \cref{phys:prop:st} to the norm identity. The same
argument applied to the columns of $U$ bounds every entry by $1$ and gives
$U_{0}^{\dagger}U_{0}=I$. For a projection and a coordinate vector $e_{j}$,
$\lVert Pe_{j}\rVert^{2}=P_{jj}$ and $\lVert(I-P)e_{j}\rVert^{2}=1-P_{jj}$ are
both nonnegative, so $0\le P_{jj}\le1$ and every entry of the $j$th column has
modulus at most $1$. Reduction preserves the self-adjointness and idempotence
identities, and it commutes with conjugation, with finite sums and with finite
products, hence with the expressions in \eqref{phys:eq:bornshadow} and with the
finite products defining a protocol; joint probabilities may be computed as
squared norms of products of the relevant projectors and unitaries applied to
$\psi$, with no normalization by a possibly infinitesimal intermediate
probability. For (d), conditioning divides by a probability and takes a
positive square root for normalization, and both commute with reduction when
the denominator's standard part is positive.
\end{proof}
\status{\TC, and conditional on adopting real standard parts as measured
probabilities (\cref{phys:sub:readout}). \TA\ It is \emph{not} a universal
theorem about every non-Archimedean probability interpretation, and not a
theorem that non-Archimedean models are experimentally indistinguishable. Two
exclusions are part of the statement: an outcome of infinitesimal probability
has zero ordinary shadow, so reduction cannot determine the state obtained by
dividing by that infinitesimal; and infinite or hyperfinite experiments lie
outside the theorem entirely.}

\begin{example}[\TE\ A rare outcome stays rare for finitely many trials]
\label{phys:ex:rare}
For the two-level state
\[
\psi=\frac{(1,\eps)}{\sqrt{1+\eps^{2}}},
\qquad p_{2}=\frac{\eps^{2}}{1+\eps^{2}},\qquad \st(p_{2})=0,
\]
any fixed ordinary finite number of independent repetitions preserves that
zero standard-part probability for observing the rare outcome. The
normalization and the rare-outcome limit are symbolically checked.
\end{example}
\status{\TA\ Introducing an infinite repetition count, or an amplifier with
infinite resources, changes the operational assumptions; it is not evidence
that the original finite experiment had a different real shadow. Postselection
on an infinitesimal-probability event lies outside
\cref{phys:thm:quantumshadow}.}

\subsection{Nonuniform amplification: the interchange genuinely fails}
\label{phys:sub:amplification}

\begin{example}[\TE\ $H=\eps\sigma_{z}$ with $T=\eps^{-1}$]
\label{phys:ex:amplify}
In units with $\hbar=1$, take $H=\eps\sigma_{z}$ and evolution time
$T=\eps^{-1}$. Then $HT=\sigma_{z}$, so the finite-phase evolution is
nontrivial although $\st H=0$. Reducing $H$ before an unbounded duration loses
the information.
\end{example}
\status{This is why \cref{phys:thm:quantumshadow} does not imply equivalence
under arbitrary nonuniform limits: reduction and unbounded duration do not
commute. \TA\ It demonstrates a \emph{mathematical} failure of interchange. It
expressly does not establish a feasible way to observe a literal infinitesimal
probability or energy shift, and a proposed observable effect must state how
such a duration, coupling, postselection or amplification is physically
realized. More generally $x\mapsto x/\eps$ promotes an infinitesimal deviation
to an order-one result; the operation is defined algebraically, but an
experimental proposal must identify the physical gain, its energy or stability
cost, and what fixes $\eps$ --- without silently inserting an infinite
resource.}

\subsection{Probability requires a summation theory, not just positive scalars}
\label{phys:sub:probability}

\TC\ For infinitely many outcomes, ``the probabilities sum to one'' needs a
definition. A countable family assigning the same positive infinitesimal $\eps$
to each outcome is \emph{not} strongly Hahn summable, since infinitely many
summands contribute to the same exponent (\cref{phys:ex:summable}). Ordinary
finite additivity therefore cannot be promoted to a countably additive
probability theory by declaring each term infinitesimal.

\TI\ Non-Archimedean probability theories exist and introduce explicit
structures for infinite addition and conditioning \cite{phys:BHW}. \TA\ They
are examples of additional mathematics a physical proposal would have to
choose, not consequences of $F$ being an ordered field. An exact infinitesimal
probability is nonzero in the ordered field and has zero standard part; for a
fixed ordinary finite number of trials the probability of seeing such an event
remains infinitesimal in a finite independent-trial model. A new probability
interpretation is possible, but it must specify how records relate to the
field-valued weights. \TI\ Ordinary measure-theoretic or nonstandard
constructions --- Loeb's among them --- can be incorporated only after their
assumptions and interpretation are supplied \cite{phys:Loeb}.

\subsection{What this does and does not say about dynamics}
\label{phys:sub:dynamics}

\TA\ \Cref{phys:thm:quantumshadow} associates an ordinary unitary matrix to
every chosen finite surcomplex unitary matrix. It does not prove that the
resulting family $U_{0}(t)$ satisfies a specified Schr\"odinger equation, is
strongly continuous, or has a particular Hamiltonian: those are differential
and topological assertions. An observable can have unlimited eigenvalues even
though all normalized-state components are limited, and a rule for reading such
an observable cannot be obtained by taking the standard part of an unlimited
number.

\TA\ The expression $U(t)=\exp(-iHt/\hbar)$ is likewise not defined by the
existence of $K=F[i]$. A real exponential on $F$ defines growth and decay; it
does not choose values for oscillatory phases with unlimited arguments. A
general complex exponential requires a compatible function theory, periodic
structure, and branch or extension choices; parametrizing the unit circle
algebraically or using finite angles does not by itself define a canonical
value of $\sin\omega$. Three tractable subclasses exist: keep the unperturbed
real-time propagator over $\C$ and solve a formal perturbation problem
coefficientwise; exponentiate a positive-order formal matrix, for which each
coefficient receives finitely many relevant terms; or work with an explicitly
constructed analytic extension on a restricted domain. What is not legitimate
is to infer a global unitary evolution for an unlimited Hamiltonian from
\cref{phys:thm:spectral} alone.

\TA\ The immediate lesson is favourable but modest: finite quantum algebra is
not obstructed by real-closed non-Archimedean coefficients. To obtain new
physics one must explain which dynamics or observational rule cannot already be
represented by the ordinary shadow.

\subsection{Infinite-dimensional quantum theory is not supplied by algebraic
closure}
\label{phys:sub:infdim}

\TA\ Ordinary Hilbert-space completeness, unbounded self-adjoint operators,
spectral measures, strong continuity of time evolution and countably additive
probability require more structure than the field operations. A space $K^{d}$
with an $F$-valued norm is not automatically a Hilbert space in the usual
real-normed sense, and a surreal-valued inner product is not an ordinary
Hilbert norm. Increasing $d$ to a surreal number is not a definition of a
vector space of that dimension, let alone a measure-theoretic construction. The
full fine topology is especially unsuitable for ordinary countable
approximation processes, by \cref{phys:prop:discrete}. Finite Fourier
identities and finite Hermitian spectral algebra do not establish infinite
Fourier completeness. A generalized quantum theory must explicitly choose its
state spaces, topology, operator domains and observables; a hyperfinite
nonstandard model, a coefficientwise deformation of an ordinary Hilbert space,
and a module over a Hahn field are three different options, and none inherits
the useful features of the others merely because its scalars admit
infinitesimals.

\TI\ There is also a distinct notion of \emph{quantum} singularity. Horowitz
and Marolf study scalar quantum probes of static singular spacetimes and
relate unique quantum evolution to essential self-adjointness of the relevant
spatial operator \cite{phys:HM}. \TA\ Their setting establishes no universal
resolution of the spacelike Schwarzschild singularity. It does illustrate what
a meaningful question looks like: does the dynamics require extra boundary
conditions? Changing scalars alone selects no operator domain and removes no
ambiguity in a self-adjoint extension. For semiclassical gravity the demand is
greater still --- one must specify the quantum state, the renormalized stress
tensor and the backreaction prescription, none of which is selected by the
availability of unlimited surreal energies or surcomplex amplitudes.

\section{Global phases exist, but are not automatically time evolution}
\label{phys:sec:phases}

\subsection{Finite phases, and a genuine global extension}
\label{phys:sub:finitephase}

\TA\ It would be wrong to dismiss surcomplex quantum theory by saying that no
global surcomplex exponential exists. One does.

\TC\ For a limited surreal angle $\theta=a+\delta$ with $a\in\R$ and $\delta$
infinitesimal, the ordinary real phase together with an infinitesimal Taylor
expansion gives
\begin{equation}\label{phys:eq:finitephase}
\cis(\theta)=\ee^{ia}\sum_{n\ge0}\frac{(i\delta)^{n}}{n!}
=\ee^{ia}\,\Emm(i\delta),
\end{equation}
where the ordinary factor $\ee^{ia}$ is supplied by complex analysis and the
infinitesimal tail is a strong Hahn sum. Every point of the surcomplex unit
circle has a \emph{limited} angle --- for instance by the rational half-angle
chart together with the ordinary choice of branch for its finite part. \TA\
This finite-angle fact assigns no physical temporal frequency to an arbitrary
unlimited angle.

\TI\ Ehrlich and Kaplan construct a global exponential using an integer part
of the ordered field; for initial surreal structures there is a canonical
choice, the omnific integers $\Oz$, and the repository's trigonometry report
includes that construction \cite{phys:EK,phys:RepoTrig}. Writing a general
surreal as $b=b_{\mathrm{PI}}+b_{\mathrm{fin}}$, separating its purely infinite
normal-form terms from its finite terms, the global normalization takes the
form
\begin{equation}\label{phys:eq:EKfinite}
\sinEK b=\sin(b_{\mathrm{fin}}),\qquad \cosEK b=\cos(b_{\mathrm{fin}}),
\end{equation}
because $b_{\mathrm{PI}}/(2\pi)$ belongs to $\Oz$, so the purely infinite part
is a period in the integer-part construction. The corresponding surcomplex
exponential is
\begin{equation}\label{phys:eq:EKexp}
\ExpEK(a+ib)=\exp(a)\bigl(\cosEK b+i\sinEK b\bigr),
\end{equation}
with kernel $2\pi i\Oz$. \TC\ It extends the ordinary complex exponential and
\emph{has} the group-homomorphism property.

\begin{remark}[The construction uses more than the bare field]
\label{phys:rem:EKinputs}
The normalization \eqref{phys:eq:EKfinite} uses the chosen integer part and,
for canonicity, the surreal initial-subtree structure --- strictly more than
the ordered field plus local Taylor axioms. This report adopts the usual
$\cos b+i\sin b$ convention in \eqref{phys:eq:EKexp}, not the reversed pair
appearing in one displayed preprint formula; the convention is fixed by the
value $1$ at $b=0$ together with the addition law. The published erratum to
\cite{phys:EK} concerns in-text bibliography reference numbering and not the
withdrawal of the construction \cite{phys:EKErratum}.
\end{remark}

\subsection{An explicit derivative incompatibility}
\label{phys:sub:phaseincompat}

The global group law is not enough to ensure compatibility with a selected
scalar derivation.

\begin{proposition}[\TC\ Canonical phase versus scalar differentiation]
\label{phys:prop:phase}
Let $D$ be a derivation on $\No$ that kills $\R$ and satisfies $D\omega=1$ ---
in particular the normalized Berarducci--Mantova derivation. Then the
functions in \eqref{phys:eq:EKfinite} do \emph{not} satisfy the global scalar
chain rule $D(\sinEK x)=\cosEK(x)\,D(x)$ for every $x\in\No$.
\end{proposition}
\begin{proof}
Since $\omega$ is purely infinite, $\sinEK\omega=0$ and $\cosEK\omega=1$.
Therefore $D(\sinEK\omega)=D(0)=0$ while $\cosEK(\omega)D\omega=1$. The
asserted global chain rule fails already at $x=\omega$.
\end{proof}
\status{\TC, and it is an \emph{incompatibility of two structures}, not a
contradiction in either established construction. A derivation of numbers and a
chosen global periodic extension are different objects, and neither
construction claims that they combine with all familiar identities without
qualification.}

\begin{corollary}[\TC\ The canonical global phase is constant in ordinary time]
\label{phys:cor:constantphase}
For every ordinary real $t$, the product $\omega t$ is purely infinite or zero,
so
\begin{equation}\label{phys:eq:constantphase}
\ExpEK(i\omega t)=1\qquad (t\in\R).
\end{equation}
As a function of an ordinary real time parameter the right-hand side is
constant. It is therefore \emph{not} a coefficientwise real-time solution of
$\psi'=i\omega\psi$ with $\psi(0)=1$.
\end{corollary}
\status{\TA\ A derivative defined using arbitrarily small \emph{surreal}
increments is a different limit operation and does not repair this real-time
substitution. The lesson is to choose a dynamical function calculus first,
rather than infer one from the existence of a phase homomorphism.
Equation~\eqref{phys:eq:constantphase} should be included as an explicit test
in any proposed phase interface: a construction that silently turns
$\ExpEK(i\omega t)$ into a real-time infinite-frequency oscillator is mixing
incompatible conventions.}

\begin{proposition}[\TC\ Infinite periods of a global phase homomorphism]
\label{phys:prop:periods}
Suppose $E:(\No,+)\to\{u\in\SC:u\bar u=1\}$ is a group homomorphism extending
the finite-angle rule \eqref{phys:eq:finitephase}. Then $E$ has an infinite
period.
\end{proposition}
\begin{proof}
Every unit-modulus $u$ is limited, so its complex standard part is $\ee^{ia}$
for some ordinary real angle $a$, and $u\ee^{-ia}=1+\eta$ with $\eta$
infinitesimal. The formal logarithm of $1+\eta$ is defined, and
$(1+\eta)\overline{(1+\eta)}=1$ shows that logarithm is purely imaginary; hence
$u=\cis(a+\delta)$ for a limited angle. Given any infinite $X$, choose such a
limited $\theta$ representing $E(X)$; then $E(X-\theta)=1$ and $X-\theta$ is
infinite.
\end{proof}

\begin{remark}[\Cref{phys:prop:periods} and $\ExpEK$ agree]
\label{phys:rem:periodsagree}
These two results are often mistaken for a conflict, so the merge records the
resolution explicitly. \Cref{phys:prop:periods} says any global phase
homomorphism extending the finite-angle rule \emph{must} have an infinite
period; $\ExpEK$ has kernel $2\pi i\Oz$, which contains infinite omnific
integers. The construction is therefore an \emph{instance} of the proposition,
not a counterexample to it. What fails is only the ordinary assertion that all
periods are ordinary integer multiples of $2\pi$: finite-angle agreement plus
the homomorphism law do not reproduce it. Neither result invalidates a chosen
global normalization, and whether such a construction represents a physical
high-frequency limit requires a link to the phase \emph{function}, its state
and its derivative, not merely to its scalar value.
\end{remark}

\subsection{Why the naive Taylor series also fails}
\label{phys:sub:naivetaylor}

\begin{remark}[\TC\ Non-summability of an infinite-angle Taylor family]
\label{phys:rem:naivetaylor}
For an infinite monomial $A$ with $\val(A)<0$, the terms of
$\sum_{n\ge0}(iA)^{n}/n!$ have exponents $n\val(A)$ descending without a least
element, so the family is not strongly Hahn summable. In particular the formal
series for $\ee^{i/\eps}$ has exponents unbounded below and is an inadmissible
Hahn Taylor substitution. Even for an ordinary nonzero real angle the direct
Taylor family fails at exponent $0$, all terms landing there; its ordinary
complex sum must be taken in $\C$ first, as in \eqref{phys:eq:finitephase}.
\end{remark}

\TA\ An appropriate oscillatory differential algebra, a transferred internal
function space, or a sectorial analytic construction might handle an
infinite-frequency limit. Such a theory may use surcomplex coefficients, but
it must specify which phases are functions, what derivatives act on them, and
how measured finite-time motion is recovered. A global extension of sine alone
is not that theory.

\subsection{The phase obstruction belongs to a companion report}
\label{phys:sub:diffeqcite}

The sharp differential-algebraic statement about oscillation over $\SC$ is not
reproved here. It is proved, in three independent ways, in the companion merged
report \emph{Differential Algebra and Differential Equations over the Surreal
and Surcomplex Numbers}, held in
	exttt{docs/surcomplex/differential-equations/}
of this repository, whose oscillator corollary establishes that for
$\lambda\in\C$ the equation $\partial y=\lambda y$ has a nonzero solution in
$\SC$ if and only if $\lambda\in\R$ --- so $\partial y=iy$ and
$\partial^{2}y+y=0$ have only the zero solution. \Cref{phys:prop:phase} and
\cref{phys:cor:constantphase} are the \emph{physical} reading of that
obstruction, and they land squarely on that report's phase conventions. This
report therefore adopts its vocabulary without exception.

\begin{convention}[The three phase symbols, imported]
\label{phys:conv:phases}
Three different objects in this literature are called ``phase'', and one of
them is not an angle at all. The companion report fixes three names and three
symbols; this report uses them and no others.
\begin{description}[style=nextline,leftmargin=0pt,labelsep=0pt]
\item[The finite angle $\theta$, living in $\Ofin$.]
The argument of $\cis$; the \emph{whole} angle of a unit, limited by theorem.
No value of a circular function at an infinite angle is ever used.
\item[The infinitesimal phase residue $\eta$, living in $\mm$.]
Every unit factors uniquely as $u=c\,\Emm(i\eta)$ with $c$ an ordinary unit
complex number; $c$ carries all the ordinary phase and $\eta$ all the variable
intrinsic phase. It is the infinitesimal \emph{part} of $\theta$, not a synonym
for it: $\theta=\st(\theta)+\eta$ and $c=\ee^{i\st(\theta)}$.
\item[The accumulated phase $B$, an arbitrary surreal.]
A primitive of the imaginary coefficient. It is \emph{not} the angle of any
surcomplex number; it need not be limited, and when it is infinite it lies
outside $\Ofin$ entirely. Its purely infinite part is the obstruction, and the
solvability criterion is exactly the finiteness of $B$.
\end{description}
\end{convention}

\begin{warning}[The sentence not to write]
\label{phys:warn:phase}
A merged sentence such as ``the phase of the solution is $\log\omega$'' is
meaningless, and the companion report deliberately prints it beside the true
one. The correct statement is: \emph{the accumulated phase of the coefficient
$i/\omega$ is $\log\omega$, which is infinite, so there is no solution.}
Likewise ``every surcomplex number has a limited phase'' is a theorem about
$\theta$; ``the variable phase of a unit is infinitesimal'' is a theorem about
$\eta$; and neither implies anything about $B$. This report never prints the
false form, and it does not re-derive the obstruction that the companion report
owns.
\end{warning}

\section{Quantum fields, divergent expansions, and renormalization}
\label{phys:sec:qft}

\subsection{The jump from finite matrices to fields is substantial}
\label{phys:sub:fields}

\TI\ Quantum field theory on a curved spacetime needs local field algebras,
states, distributional correlation functions, a suitable spectrum condition and
renormalized composite observables; in a Hilbert-space presentation it also
needs infinite-dimensional completion and operator domains. Hollands and Wald
give a precise formulation of these issues for ordinary globally hyperbolic
spacetimes \cite{phys:HW}.

\TA\ Replacing the coefficient field supplies none of a normed-space
completion, countable additivity, a spectral measure, or the hypotheses of a
unitary-evolution theorem. A chosen set-sized topology or a coefficientwise
construction may be useful, but it is additional input.

\subsection{Formal perturbative coefficients are a realistic application}
\label{phys:sub:formalpert}

\TI\ Formal parameters are already a powerful organizing principle in
mathematical physics; deformation quantization is a well-developed example in
which associative products are built as formal series rather than convergent
numerical sums \cite{phys:Kontsevich}.

\TA\ A Hahn or surreal realization can place several such scales in one ordered
framework and add fractional powers, logarithms or exponential sectors where a
justified differential structure is available; a scale-resolving field can
represent quantities of the schematic form
\begin{equation}\label{phys:eq:schematictrans}
\sum_{\alpha,k,n}c_{\alpha k n}\,
\eps^{\alpha}\bigl[\log(1/\eps)\bigr]^{k}\exp(-nA/\eps),
\end{equation}
\emph{provided} the chosen transseries construction admits the displayed
support and operations. The display is schematic: \emph{an arbitrary
triple-indexed family is not automatically admissible.} In suitable classes it
keeps algebraic, logarithmic and beyond-all-orders sectors separate, which can
help formulate matching conditions, compare competing terms, and retain
exponentially small information a power expansion discards. Much of this work
can nevertheless be done in a smaller Hahn or transseries field; the full
proper class is usually an ambient unifying language, not a computational
necessity.

\TA\ This improves the handling of orders of magnitude. It does not make an
arbitrary formal expansion into a unique actual state or scattering amplitude.

\subsection{A divergent series, an exact element, and hidden data}
\label{phys:sub:borel}

\begin{computation}[\TE\ The factorial series is an exact Hahn element]
\label{phys:comp:factorial}
The series
\begin{equation}\label{phys:eq:factorial}
\sum_{n\ge0}(-1)^{n}n!\,\eps^{n}
\qquad\text{and}\qquad
\widehat F(\eps)=\sum_{n\ge0}n!\,\eps^{n+1}
\end{equation}
are well-defined Hahn elements for positive-valuation $\eps$, although the
associated ordinary power series have zero radius of convergence: for nonzero
real $z$ the ratio of successive term magnitudes in $\widehat F$ is
$(n+1)\abs z$. Exact denotation at $\eps$ does not give a real-valued sum at
$z$; the formal series and an ordinary function with that asymptotic expansion
are different objects.
\end{computation}

\begin{computation}[\TE\ The beyond-all-orders ambiguity, made explicit]
\label{phys:comp:beyond}
Consider the dimensionless equation
\begin{equation}\label{phys:eq:factorialODE}
y'(x)+y(x)=\frac1x,\qquad x\to+\infty .
\end{equation}
Substituting $\sum a_{n}x^{-n-1}$ gives $a_{0}=1$ and $a_{n}=na_{n-1}$, so
$a_{n}=n!$ and the formal inverse-power solution is
$\widehat y(x)=\sum_{n\ge0}n!\,x^{-n-1}$. For the truncation $y_{N}$,
\begin{equation}\label{phys:eq:factorialresidual}
y_{N}'+y_{N}-\frac1x=-\frac{(N+1)!}{x^{N+2}},
\end{equation}
checked exactly for $N=1,\dots,8$. At $x=\omega$ the series is Hahn summable,
and under a compatible strong derivation with $D\omega=1$ its differentiated
terms cancel in the same way. But the homogeneous equation has solutions
$C\ee^{-x}$, so
\begin{equation}\label{phys:eq:ambiguity}
y(x)=y_{\mathrm{part}}(x)+C\ee^{-x}
\end{equation}
shares the same inverse-power expansion for every ordinary constant $C$; in the
enlarged exponential field $C\ee^{-\omega}$ is a distinct beyond-all-orders
scale. The exponentially small homogeneous ambiguity is itself a checked
identity.
\end{computation}
\status{\TA\ The coefficient sequence does not specify the complete solution:
$C$ must come from boundary or initial data, or from a prescribed summation and
matching condition. A richer field can \emph{retain} that information once
supplied; it cannot infer it from data that do not determine it. For positive
ordinary $x\to0$, two functions differing by $C\ee^{-1/x}$ have the same power
asymptotic expansion to every order, and arithmetic alone selects neither $C$,
nor a contour prescription, nor boundary conditions. This is a miniature of the
difficulty faced when an asymptotic black-hole solution is mistaken for
complete dynamics.}

\begin{computation}[\TE\ Borel transform and lateral summation]
\label{phys:comp:borel}
With the Borel convention sending $n!z^{n+1}$ to $\xi^{n}$, the transform of
$\widehat F$ is
\begin{equation}\label{phys:eq:borel}
\mathcal B\widehat F(\xi)=\frac{1}{1-\xi},
\end{equation}
and the lateral Laplace integrals bypassing the pole on opposite sides,
\begin{equation}\label{phys:eq:lateral}
F_{\pm}(z)=\int_{0}^{\infty \ee^{\pm i0}}
\frac{\ee^{-\xi/z}}{1-\xi}\,\dd\xi,
\end{equation}
differ by a quantity of magnitude and phase proportional to
$2\pi i\,\ee^{-1/z}$, by the residue of $\ee^{-\xi/z}/(1-\xi)$ at $\xi=1$; the
sign depends on the orientation convention. The residue is symbolically
checked. The two functions have the same formal power asymptotics away from the
Stokes direction and differ by information invisible at every algebraic order.
\end{computation}

\begin{remark}[\TC\ A genuine infinitesimal needs a larger workspace]
\label{phys:rem:notinrank1}
For a genuine positive infinitesimal $\eps$, the element $\ee^{-1/\eps}$ is
smaller than every $\eps^{N}$, and it is \emph{not} an element of the rank-one
field $\R((\eps^{\Q}))$ with that prescribed monomial scale, since a nonzero
element there has a least rational exponent. A larger workspace or a named
transseries subfield is required. Extending the arithmetic and choosing how it
represents an actual function are two separate steps.
\end{remark}

\TA\ A transseries must therefore include a suitable exponential sector
\emph{and} data selecting its coefficient. Naming the Hahn element has selected
no boundary condition, no lateral summation and no physical state. \TI\ This is
exactly why Stokes data and realization matter in resurgent physics
\cite{phys:Resurgence}. \TA\ The extension and summation choices in the
established surreal integration theories are precisely additional structure,
not consequences of an unadorned infinite sum \cite{phys:CE,phys:BMgerms}.

\subsection{A conditional asymptotic-to-surreal pipeline}
\label{phys:sub:pipeline}

\TA\ A rigorous application can proceed through a restricted pipeline: start
from a specified family of real or complex solutions and a dimensionless
parameter; derive its admissible asymptotic expansion, with uniformity on a
named ordinary domain; embed that expansion in a differential Hahn or
transseries field whose derivation represents the selected variable; and prove
that the observables extracted from the representation agree with the original
solution up to a stated error. The output is a formal or an analytic
certificate in the sense of \cref{phys:sub:certificates}. For singular ODE
reductions, existing surreal integration results suggest a feasible bridge
where their hypotheses apply \cite{phys:CE}; for the unrestricted nonlinear
Einstein system that bridge would require substantially more work. A precise
restricted model is preferable to an assertion of universal transfer.

\subsection{Unlimited values do not define products of distributions}
\label{phys:sub:distributions}

\begin{computation}[\TE\ A Gaussian regulator and its square]
\label{phys:comp:gaussian}
For an ordinary positive real regulator $\eps$ put
$\delta_{\eps}(x)=(\eps\sqrt\pi)^{-1}\exp(-x^{2}/\eps^{2})$. Gaussian
integration gives
\begin{equation}\label{phys:eq:deltareg}
\int_{\R}\delta_{\eps}(x)\dd x=1,
\qquad
\int_{\R}\delta_{\eps}(x)^{2}\dd x=\frac{1}{\eps\sqrt{2\pi}} .
\end{equation}
Both integrals are symbolically checked. The first family tends to the Dirac
distribution; the second has a divergent integral.
\end{computation}
\status{\TA\ Naming the right-hand side of the second identity an unlimited
surreal value constructs no canonical distribution $\delta^{2}$ and proves no
regulator independence. Note also what was used: the integral in this example
was the \emph{ordinary} integral for each real regulator, and no unannounced
surreal integration rule appeared.}

\begin{computation}[\TE\ Pointwise infinitesimals do not produce a
distribution]
\label{phys:comp:distrib}
In ordinary distribution theory
\begin{equation}\label{phys:eq:distribution}
\lim_{\eta\downarrow0}\frac{1}{x+i\eta}=\PV\frac1x-i\pi\delta(x),
\end{equation}
which follows from $\frac{1}{x+i\eta}=\frac{x}{x^{2}+\eta^{2}}
-i\frac{\eta}{x^{2}+\eta^{2}}$: against a smooth compactly supported test
function the imaginary kernel integrates to $\pi$ in the limit, by $x=\eta y$
and $\int_{\R}(1+y^{2})^{-1}\dd y=\pi$, while the odd real kernel tends to the
principal value after locally subtracting the test function's value at zero.
This is an \emph{integral} statement, not a pointwise identity of functions.
Setting $\eta$ instead to a fixed positive infinitesimal, pointwise standard
part gives $1/x$ at each nonzero ordinary $x$, and at $x=0$ the value is
unlimited so standard part is unavailable. \emph{The $-i\pi\delta(x)$
contribution cannot be recovered by pointwise standard part at all.}
\end{computation}
\status{\TA\ A test-function pairing, with an explicitly justified
coefficientwise, asymptotic or nonstandard integration construction, is
required. Writing a surcomplex $i\eps$ has supplied no distributional
prescription. The general lesson recurs throughout field theory: information can
concentrate in a shrinking region whose pointwise ordinary shadow vanishes away
from that region, so a global integration rule must resolve the region
\emph{before} taking the observable shadow.}

\TI\ Generalized-function approaches to nonlinear equations do exist, including
frameworks used in general relativity, and they retain carefully specified
regularization data and notions of equivalence or association
\cite{phys:SV}. \TA\ They should not be confused with scalar field extensions.
Surreal coefficients could be incorporated into such a framework, but the
mathematical work lies in defining its operations and their relation to
physical observables. Substituting an unlimited scalar for a delta distribution
does not retain its action on test functions, and a black-hole proposal
involving weak stress tensors must address that explicitly whatever field it
chooses.

\subsection{Removing negative powers is not a renormalization theory}
\label{phys:sub:renorm}

\begin{computation}[\TE\ Finite part is neither multiplicative nor
reparametrization invariant]
\label{phys:comp:FP}
On the ring of limited elements, standard part is multiplicative. On all
Laurent series, the operation $\FP$ extracting the exponent-zero coefficient is
not:
\begin{equation}\label{phys:eq:FPnonmult}
\FP(\eps)=\FP(\eps^{-1})=0,\qquad \FP(\eps\cdot\eps^{-1})=\FP(1)=1 .
\end{equation}
Nor is a bare finite-part prescription invariant under a harmless-looking
change of regulator coordinate. With $\theta=\eps/(1+a\eps)$ for $a\in\R$,
\begin{equation}\label{phys:eq:finitepart}
\eps^{-1}=\theta^{-1}-a,
\qquad \FP_{\eps}(\eps^{-1})=0,
\qquad \FP_{\theta}(\eps^{-1})=-a .
\end{equation}
Equivalently, in the other manuscripts' parametrization: $F(\eps)=\eps^{-1}+a$
has exponent-zero coefficient $a$, but under the equally legitimate formal
change $\eps=\eta+c\eta^{2}$,
\begin{equation}\label{phys:eq:noncanonical}
F(\eta+c\eta^{2})=\eta^{-1}+(a-c)+c^{2}\eta-\cdots,
\end{equation}
so the constant coefficient becomes $a-c$. Both reparametrization shifts are
symbolically checked.
\end{computation}
\status{So ``keep the finite part'' of an unlimited expression is
scheme-dependent. \TA\ This does \emph{not} contradict the uniqueness of
standard part in \cref{phys:prop:st}, because $F$ was never in $\Ofin_{F}$.
Physical renormalization must specify observables, allowed counterterms and
normalization conditions; field enlargement does not fix that freedom, and
effective-field-theory reasoning remains necessary for identifying which
interactions and parameters have physical significance \cite{phys:Donoghue}. A
coefficient-extraction rule must be shown to respect the relevant coordinate
and gauge symmetries, and its observable consequences must not depend on
arbitrary regulator choices --- otherwise a supposedly resolved curvature is a
scheme-dependent number rather than a physical prediction.}

\TA\ Surreal arithmetic can still help here, and the help is precisely
specifiable: divergent powers and logarithms can be retained symbolically,
counterterm cancellation can be tested exactly, and residual prescription
dependence can be \emph{exposed} rather than hidden. That is an improvement in
representation and verification. It is not a proof that ultraviolet divergences
have been physically removed.

\subsection{Hawking radiation and information require state evolution}
\label{phys:sub:hawking}

\TI\ The Hawking effect can be formulated using ordinary quantum fields on
suitable classical black-hole backgrounds, with the usual care about states and
renormalized observables \cite{phys:HW}. \TA\ A black-hole information proposal
requires an account of the full evolution and of the observables accessible to
relevant observers. Making an interior curvature or an intermediate amplitude
surcomplex neither supplies that evolution nor proves its unitarity. A
successful asymptotic calculation could certainly assist --- controlling a
near-horizon mode expansion, a late-time tail, or a nonperturbative correction
--- but to claim a physical information-resolution mechanism one would still
need to identify a state map, prove the relevant consistency properties, and
relate the result to finite measurement outcomes. Those are the missing
physical statements, not missing names for large numbers.

\subsection{A realistic wave-physics benchmark}
\label{phys:sub:qnm}

\TI\ Black-hole quasinormal-mode and scattering problems already use complex
analysis, asymptotic series and boundary conditions; Hatsuda relates
quasinormal frequencies to high-order WKB expansions and shows, in the cases
studied, how Borel summation recovers the numerical frequencies
\cite{phys:Hatsuda}. \TA\ That is a concrete and \emph{testable} benchmark for
a surreal-compatible implementation: preserve the perturbative and
nonperturbative sectors, enforce the physical boundary conditions, and
reproduce an established complex frequency calculation. Success would
demonstrate calculational utility. It would not resolve an interior spacetime
singularity, because the wave problem and its boundary conditions are separate
from a replacement of the central geometry. The benchmark is attractive
precisely because the claim is limited and can fail.
\part{Context, interpretation, and what to do next}
\label{phys:part:context}

\section{Existing proposals and the status of the idea}
\label{phys:sec:proposals}

\subsection{A direct surreal-singularity proposal deserves a direct test}
\label{phys:sub:nieto}

\TI\ Nieto's \emph{Some mathematical and physical remarks on surreal numbers}
explicitly discusses gravitational singularities and proposes that the
availability of infinite surreal values removes the usual singularity problem;
a later paper connects surreals to dualities, matroids, qubits and twistors
\cite{phys:Nieto2016,phys:Nieto2018}. In the Schwarzschild discussion the
argument passes from a larger scalar domain to a claim that the singularity
disappears, with no corresponding construction of a regular differential
spacetime or a geodesic continuation. That such work exists establishes that
the suggestion has precedents, and it is not a question that can be dismissed
as unasked or meaningless.

\begin{assessment}[\TA\ The missing step, located]
\label{phys:asmt:nieto}
The computations of \cref{phys:sec:schwarzschild,phys:sec:infall} identify
precisely where the inference fails, in four places.
\begin{enumerate}[label=\textup{(\arabic*)},leftmargin=2.4em]
\item At a nonzero infinitesimal radius the metric is evaluable, but by
\cref{phys:comp:Keps} and \cref{phys:prop:unlimited} the invariant curvature is
still unlimited in the real-boundedness sense, and by
\cref{phys:prop:pole} the $r=0$ pole is unchanged under any radial
reparametrization of the stated form.
\item The exact identity $r^{6}\Kret=48m^{2}$ still has \emph{no} field-valued
solution at $r=0$, by \cref{phys:prop:noendpoint}, in \emph{any}
characteristic-zero field and not merely in $\R$.
\item The classical endpoint is still reached at finite --- indeed
infinitesimal, hence limited --- proper time, by
\cref{phys:comp:finiteproper}, and the tidal equations
\eqref{phys:eq:jacobi} retain their singular coefficients.
\item No continuation law follows from any of the above, so none of (Q1)--(Q4)
of \cref{phys:sub:fourquestions} is answered.
\end{enumerate}
Therefore the proposed inference is not a singularity-resolution theorem. The
useful intuition in it is the \emph{resolution of scales}, not a demonstrated
removal of the endpoint.
\end{assessment}
\status{\TA\ This criticism concerns that \emph{inference}. It is not an
assertion that no future theory using surreal variables could make nontrivial
physical predictions, and it is not an impossibility theorem. A complete
alternative might replace the metric, the state space, the observables or the
dynamical law --- but those replacements must be constructed.}

\TA\ The literature reviewed by the three manuscripts does not establish an
empirically tested theory in which scalar replacement alone resolves a
black-hole singularity. That is a bounded assessment of the cited work. It is
not a proof that no future non-Archimedean theory can do so, and the
established mathematical results on surreal differential fields and integration
are genuine and must not be conflated with speculative gravitational
interpretations. Each manuscript records that its literature search was
targeted rather than systematic: absence of a predictive model in the inspected
sources is not proof that no other proposal has ever been written.

\subsection{Comparison with other non-Archimedean and singular methods}
\label{phys:sub:comparison}

\TA\ The choice of tool should be driven by the operation required, not by the
size of the number system, and a research program should explain why surreals
are preferable to an existing smaller tool for a particular task.
Universality alone is neither an efficiency theorem nor a prediction.
\Cref{phys:tab:comparison} merges the three manuscripts' comparison tables. Its
last column records a \emph{missing implication}, not a blanket impossibility,
and the table is not a ranking of research programs.

\begin{longtable}{@{}L{0.19\textwidth}L{0.36\textwidth}L{0.38\textwidth}@{}}
\caption{What each framework is naturally for, and what does not follow
automatically. The third column is a missing implication, never a blanket
impossibility.}
\label{phys:tab:comparison}\\
\toprule
Framework & Principal useful structure & What does not follow automatically \\
\midrule
\endfirsthead
\toprule
Framework & Principal useful structure & What does not follow automatically \\
\midrule
\endhead
\bottomrule
\endfoot
Real/complex asymptotics &
Limits, error bounds, analytic continuation, well-developed PDE and operator
methods &
Existence of a quantum-gravity completion of a classical singularity \\
\addlinespace
Formal power and Puiseux series &
Finite coefficient algorithms; ramified scaling; clear order-by-order equations
&
A convergent or uniquely selected real solution \\
\addlinespace
Hahn fields &
Exact ordered powers, multiple scales, support-certified arithmetic, finite
algebra and valuation balances &
Function domains, derivatives, integration, realization of formal series;
arbitrary smooth function theory or physical interpretation \\
\addlinespace
Transseries and resurgence &
Logarithmic and exponential asymptotics, beyond-all-orders sectors, asymptotic
differential algebra, selected reconstruction methods &
Admissible support, summation, Stokes and boundary data, error estimates;
universal applicability to nonlinear spacetime PDEs \\
\addlinespace
Full surreal/surcomplex scalars &
An exceptionally broad ordered and exponential ambient system, algebraic
closure, compatible differential structures for many scales, distinguished
normal forms &
An ordinary continuum topology, a canonical global oscillatory calculus, a
probability theory, a measurement rule, or singularity resolution;
set/class discipline and a chosen analytic framework remain to be supplied \\
\addlinespace
Hyperreal and nonstandard models &
Infinitesimal encodings with a specified transfer and internal-set framework,
with ordinary shadows &
New predictions merely from reformulating an old model; internal models,
saturation, measure and dynamics hypotheses must be supplied \\
\addlinespace
$p$-adic models &
Non-Archimedean norm geometry and explicitly constructed alternative field
theories &
New geometric and physical interpretations; no inherited real ordering or
Lorentzian causality \\
\addlinespace
Distributions and generalized-function geometry &
Weak limits, singular sources, boundary values, a specified notion of weak
field and regularized nonlinear operation &
Products or nonlinear equations may need extra structure; uniqueness or
regulator independence without additional hypotheses \\
\addlinespace
Regular-core or other modified dynamics &
New field equations or effective sources, possibly a regular core &
Derivation, empirical validity, global stability, geodesic completeness,
unitarity, or observable consistency from local regularity alone \\
\end{longtable}

\TA\ Four specific cautions belong with the table.

First, if a calculation is unchanged when performed in a small Hahn field, its
success does not demonstrate that the full proper class of surreals is
physically necessary. A positive result should identify whether the essential
contribution is the larger field, a support theorem, a new extension operator,
or simply a more transparent notation for an existing asymptotic argument. For
an individual expansion with rational powers a Puiseux field may suffice; for
several ordered scales a Hahn field may be the right working object. Surreals
become attractive as an ambient framework connecting many such structures, or
supporting a particularly useful normal form --- not because they are larger.

Second, an ordered non-Archimedean field can sit inside a surreal ambient
field without thereby identifying the accompanying structures. A hyperreal
formulation usually specifies a nonstandard enlargement, internal functions and
sets, and a transfer principle for that enlarged structure; those are more than
a field embedding \cite{phys:Ehrlich,phys:Loeb}.

Third, $p$-adic physics is not surreal physics in different notation: $\Q_{p}$
carries no compatible field order, whereas order is central to the real-surreal
modulus and the causal-sign constructions used throughout this report.
\TI\ Gubser and collaborators' $p$-adic AdS/CFT model is instructive precisely
because it chooses a $p$-adic boundary and a tree-like bulk explicitly and
computes correlation functions \cite{phys:Padic}: a concrete alternative
geometric model, neither an ordered extension of the real spacetime metric nor
evidence about surreal numbers and Schwarzschild.

Fourth, a complex contour can detour around a pole, but an analytic
continuation along that contour is not necessarily a real Lorentzian
trajectory. Ordinary $\C$ already permits such detours; changing to $\SC$
specifies no continuation, no boundary condition and no causal history. To turn
such a construction into a real physical continuation one must recover a real
Lorentzian slice, specify the reality conditions, choose the branch and
matching data, and show that the resulting evolution satisfies the intended
equations. Analytic continuation in a parameter is not the same operation as
extending an infalling observer's causal worldline.

\TA\ For this repository, a set-sized Hahn or transseries layer is the
lowest-risk first implementation: it uses the existing summability and
finite-algebra work while avoiding an unnecessary commitment to a full
surreal-valued event manifold. A later theory can explain when embedding those
calculations into $\No$ supplies a substantive additional benefit.

\section{Measurement, units, recoverability, and empirical novelty}
\label{phys:sec:measurement}

\subsection{Which structure is supposed to be physical?}
\label{phys:sub:whichstructure}

\TA\ A scalar replacement must say whether it changes the set of events, the
values of fields on an unchanged spacetime, the codomain of measurement
results, or only an approximation calculus. These lead to different theories.
Limited surreal field values on a real manifold can have an ordinary shadow
while retaining formal infinitesimal corrections; a manifold modelled directly
on a non-Archimedean field instead requires a new account of paths,
integration and initial data. \emph{There is no single object called ``the
surreal spacetime'' determined by the field axioms.}

\TA\ The simplicity relation and the birthday of a surreal number are extra
mathematical structure, not automatically a physical age. A rule identifying an
ordinal birthday with time would have to explain how ordinary time translation,
changes of units and coordinate covariance act on it. No such identification is
needed for the coefficient-based program developed here, and none is made.

\subsection{What an ordinary apparatus is supposed to report}
\label{phys:sub:readout}

\TA\ There are at least two readings of a surreal-valued observable. It can
encode the asymptotic behaviour of a family of ordinary observables, in which
case there is a parameter-to-physical-system map and an approximation problem.
Or it can be taken literally as the value in a new ontology, in which case an
apparatus rule must map it to records and probabilities.

The standard-part map is a simple apparatus rule for limited values. It is
\emph{not forced by field algebra}; it is natural when infinitesimals are meant
to represent deviations below every real resolution, and it is an additional
modelling assumption. It is undefined on an unlimited curvature such as
\eqref{phys:eq:Keps}, and extending the codomain to include $+\infty$ still
provides no ordinary finite detector output.
\Cref{phys:thm:reduction,phys:thm:quantumshadow} show how conservative this
rule is: regular limited data reduce to the usual theories in the stated
models.

\TC\ More generally, any expression built from limited inputs by addition,
multiplication, conjugation and inverses with nonzero standard part commutes
with reduction. So a finite algebraic readout protocol cannot distinguish two
inputs whose difference is infinitesimal under those conditions. \TA\ That is a
mathematical statement, not an assertion that every imaginable physical
apparatus is such a protocol.

\subsection{Ways to leave the conservative regime}
\label{phys:sub:leaving}

\TA\ A new prediction could arise through a singular rescaling, an unbounded
duration, a nonstandard observable rule, or an explicit change of dynamics.
Each is a meaningful research possibility, and each changes more than the bare
choice of scalar field --- \cref{phys:ex:amplify} is the cleanest illustration
of the second, and \cref{phys:comp:FP} of the third. If an ordinary finite
experiment with bounded couplings is claimed to see such an effect, the model
should derive that result without silently inserting an infinite resource.

\subsection{A new scalar ontology may have no new finite predictions}
\label{phys:sub:novelty}

\TA\ \Cref{phys:prop:st} and \cref{phys:thm:quantumshadow} show how a large
numerical extension can reduce to ordinary predictions for a specified class of
observations. That is not a defect: many valuable mathematical reformulations
are empirically equivalent. It does mean that an ontological claim --- ``the
true amplitude contains an infinitesimal term'' --- needs an operational
consequence before it becomes a testable physical distinction.

\TA\ Invariance of predictions is the decisive test. A physically meaningful
result should survive changes of admissible coordinates, equivalent asymptotic
parametrizations, choices of embedding into a larger Hahn workspace, and
regulator prescriptions subject to the same physical normalization conditions.
That requirement is strictly stronger than formal equality inside one chosen
representation, and \eqref{phys:eq:finitepart} demonstrates how an apparently
finite prediction can shift under a reparametrization when the observable rule
is underspecified. Conversely the invariant curvature identities and
\cref{phys:prop:frame} exhibit quantities that can be controlled
intrinsically. If a prediction changes when an arbitrary infinitesimal is
replaced by another with no accompanying change of physical inputs, the model
has not yet separated physics from its regulator; if the choice is intended to
be physical, it must be promoted to a law or a measurable parameter rather than
hidden in notation.

\TA\ No empirical preference for a literal surreal scalar ontology follows from
anything in this report. The immediate case for these methods rests on
mathematical utility. An experimentally distinct theory would have to exhibit a
definite finite prediction, show why ordinary effective descriptions do not
already contain it, and connect its parameters to a reproducible physical
procedure.

\subsection{The relevant notion of resolution must be stated}
\label{phys:sub:notionofresolution}

\TA\ A theory might allow unlimited intrinsic quantities while maintaining
finite detector responses. Another might replace a spacetime endpoint by a
transition between quantum states. A third might define a weak geometric
extension. These are logically possible strategies and nothing here rejects
them by definition. Each carries a distinct burden of proof: the first must
\emph{derive} detector responses rather than postulate a standard part of a
nonexistent limited value; the second needs a dynamical state map and an
interpretation of time; the third needs a regularity class and an evolution or
matching principle. Examples of what ``resolved'' might mean include extension
of every relevant causal geodesic, bounded detector responses, existence and
uniqueness of a weak evolution, and a specified quantum propagation criterion
--- and these conditions are \emph{not equivalent}. In all cases the physical
gain would be the successful construction and its consequences, never the
permission to write an unlimited scalar.

\section{Relation to the companion reports in this repository}
\label{phys:sec:companions}

\TA\ This report is the only member of its collection whose subject is
empirical. Nothing in the neighbouring reports contains general relativity,
quantum mechanics or field theory, so the report carries the whole burden of
its own reading contract and establishes the theorem-versus-interpretation
discipline of \cref{phys:sec:firewall} for itself. Two companion reports are
nevertheless load-bearing, and in both cases the correct relation is to
\emph{inherit and cite}, not to restate.

\paragraph{\emph{Foundations for Surreal and Surcomplex Mathematics}.}
In 	exttt{docs/foundations-and-computation/foundations/}.
That report owns the workspace-localization theorem and the size obstructions
any candidate foundation must respect. \Cref{phys:sub:workspace} cites its
localization theorem for the workspace discipline used throughout, and defers
to its size obstructions for anything bearing on proposal (P3), a genuinely
surreal-valued spacetime. This report does not re-argue set-theoretic legality,
and there is no content duplication: the physics manuscripts ask a different
question of the same facts, namely which of (P1)--(P3) a given calculation
actually needs. Readers evaluating a proposed surreal-valued event manifold
should begin there and not here.

\paragraph{\emph{Differential Algebra and Differential Equations over the
Surreal and Surcomplex Numbers}.}
In 	exttt{docs/surcomplex/differential-equations/}.
That report owns the oscillation obstruction, the finite-primitive criterion
and the three phase symbols. \Cref{phys:sub:diffeqcite} cites its oscillator
corollary rather than re-deriving the obstruction;
\cref{phys:conv:phases} adopts its three phase symbols verbatim, and
\cref{phys:warn:phase} reproduces its warning about the sentence not to write.
\Cref{phys:prop:phase} and \cref{phys:cor:constantphase} are the physical
reading of that obstruction and land on that report's phase conventions; they
strengthen nothing there and contradict nothing there.

\paragraph{The surcomplex analysis and trigonometry reports.}
\Cref{phys:sub:commondomain,phys:sub:fine} use the analysis report's
distinction between Hahn summation and ordinary convergence, its analytic
function classes, and the common-domain warning example
\eqref{phys:eq:nocommondisk}; \cref{phys:sub:finitephase} uses the
trigonometry report's inclusion of the Ehrlich--Kaplan integer-part global
normalization together with its explicit refusal to select a scalar derivation
\cite{phys:RepoAnalysis,phys:RepoTrig}. Neither is treated as a verified
physics foundation, and the trigonometric and complex-algebraic material must
not be read as a finished unrestricted phase calculus for quantum evolution.

\paragraph{The surquaternionic family (\texttt{docs/surquaternions/}).}
Named only to prevent a transfer. That material is noncommutative, so
conjugation, modulus multiplicativity and valuation need separate treatment
there, and nothing in this report --- every result of which uses commutativity
somewhere --- may be lifted into it without reproof. The last row of
\cref{phys:tab:derivatives} records the one specific hazard: the quaternionic
radial exponential is not an additive-to-multiplicative homomorphism for
noncommuting arguments, while $\ExpEK$ of \eqref{phys:eq:EKexp} is a
homomorphism on the commutative surcomplex field. Both statements are true of
different objects.

\section{A staged research program}
\label{phys:sec:program}

\TA\ Everything in this section is tier \TA: these are proposals, not
descriptions of existing modules and not results. The central research
challenge is to make a proposal precise enough that it can \emph{fail} a test.
A larger collection of named numbers is not yet such a proposal.

\subsection{Recommended architecture}
\label{phys:sub:architecture}

The most defensible starting point is
\begin{equation}\label{phys:eq:architecture}
\boxed{\begin{gathered}
\text{ordinary spacetime}\\
+\;\text{controlled differential Hahn/transseries coefficients}\\
+\;\text{explicit real-shadow map}
\end{gathered}}
\end{equation}
This preserves the classical geometric domain while enriching the asymptotic
fields. \TA\ It is called the most defensible starting point; it is \emph{not}
a derived necessity. A fully non-Archimedean spacetime can be studied
separately, but should not be required before the coefficient approach has
produced a demonstrable advantage.

A coefficient interface should record the value group, support certificates,
coefficient domains, conjugation, the real form, and the derivatives under
which the space is closed; and it should distinguish a spacetime derivative
from a scalar derivation and from a formal-parameter derivative \emph{at the
type or interface level}, per \cref{phys:tab:derivatives}. It should store the
coefficient domain, exponent group, derivative convention and any summation
prescription alongside each series; reject a product or substitution whose
support condition fails; and distinguish an exact Hahn equality from an
asymptotic equality of functions and from equality after standard part. The
repository's existing finite algebra, Neumann support, rational-jet and
standard-part layers are useful prerequisites, and their stated current scope
must be respected.

\subsection{A minimal model specification}
\label{phys:sub:minimalspec}

A candidate theory should specify: its scalar workspace and foundational size;
its spaces of states, events and fields; its topology and derivatives; its
rules for infinite sums and integration; its field equations and admissible
initial data; its causal or quantum evolution law; its gauge and unit
conventions; and its operational readout. Recovery of ordinary physics must be
a theorem or a controlled approximation \emph{in the chosen model}, never an
appeal to the inclusion $\R\subset\No$. For a singularity claim the
specification must additionally state which of the meanings in
\cref{phys:sub:notionofresolution} is intended.

\subsection{Stage A: formalize a conservative local theory}
\label{phys:sub:stageA}

Target the model of \cref{phys:sec:reduction}, not unrestricted quantum
gravity. Define a common-domain algebra of globally supported coefficient
functions; prove support preservation under products and coefficient
derivatives; construct matrix inversion around a nonsingular ordinary leading
matrix; and make standard part commute with these operations \emph{by theorem},
not by convention. Then define tensors, the Levi-Civita connection, curvature
and the Einstein tensor, and prove \cref{phys:thm:einstein},
\cref{phys:thm:reduction}, the finite Bianchi identities,
\cref{phys:prop:frame}, and compatibility with ordinary coordinate
transformations, recording the exact coefficient-ring assumptions. Proposed
module names such as \texttt{Physics/CoefficientFields},
\texttt{Physics/FormalMetric} and \texttt{Physics/StandardPartGR} are
\emph{proposed additions}, not claims about the inspected tree.

The failure conditions matter as much as the successes, and they should be
\emph{reported} rather than papered over: a zero leading determinant; a
non-well-ordered global support; a coefficient family lacking a common domain;
an unsupported substitution. A system that assigns a symbolic infinity and
proceeds has failed.

\subsection{Stage B: verify exact backgrounds and observer diagnostics}
\label{phys:sub:stageB}

Use Schwarzschild as a mandatory control case: verify the cancellation of the
horizon-coordinate singularity (\cref{phys:comp:horizon}), the vacuum Ricci
tensor away from zero, the Kretschmann invariant
(\cref{phys:comp:Kschwarz}), the radial proper-time scaling and the Jacobi
exponents (\cref{phys:comp:infall,phys:comp:jacobi}), in \emph{both} static and
ingoing coordinates, and the Kasner family (\cref{phys:comp:kasner}) to test
multiple real exponents and vanishing leading coefficients. Add the
regular-core ansatz to distinguish changed dynamics from scalar extension.

\emph{The benchmark is passed only if the method identifies the horizon as
removable and $r=0$ as obstructed under the unchanged invariant identity.} A
system that declares both regular merely because the denominators have become
surreal symbols has failed. Success means more than producing a large
expression: the tool should retain units, verify agreement of invariant
quantities between charts, certify which denominators are units, and report
when a requested evaluation meets a zero denominator or an unsupported
expansion. In particular it must distinguish $g_{rr}$ at the Schwarzschild
horizon from \eqref{phys:eq:Keps} at the central curvature pole. The supplied
programs check the concrete identities used in this article; they are not a
proof-assistant verification of the general framework.

\subsection{Stage C: certified dominant balances, matching, and one matter
problem}
\label{phys:sub:stageC}

Implement the support test and the exact finite remainder of
\cref{phys:prop:matching,phys:comp:remainder}. The desired output is a
proof-carrying classification of outer, transition and inner scales together
with the exact rational model the expansions came from; extend the test to
several small parameters and to fractional-power ansatzes. A further benchmark
is a reduced gravitational or matter ODE with a controlled singular expansion:
compute its coefficients, prove the support pattern inductively, and compare
with a real solution through a remainder theorem or a valid summation
construction. \emph{Formal recurrence alone is not an existence theorem}, and
in applying a surreal integration result one must verify its nonresonance,
function-class and branch hypotheses explicitly.

A defensible matter experiment is a scalar field on a fixed spherically
symmetric background with a specified asymptotic region and initial or
characteristic data: derive a consistent expansion for the linear equation,
compare it with ordinary asymptotic estimates, and only then attempt
backreaction in the Einstein--massless-scalar system
\begin{equation}\label{phys:eq:Einsteinscalar}
R_{ab}=8\pi\,\partial_{a}\phi\,\partial_{b}\phi,
\qquad \Box_{g}\phi=0,
\end{equation}
in geometrized units with zero cosmological constant, these being the equations
of the usual minimally coupled massless scalar action, which also fix the
normalization of $\phi$. In double-null spherical coordinates, with
\begin{equation}\label{phys:eq:doublenull}
\dd s^{2}=-\Omega^{2}(u,v)\,\dd u\,\dd v+r^{2}(u,v)\dd\Omega_{2}^{2},
\qquad g_{uv}=-\Omega^{2}/2,
\end{equation}
one null constraint and the geometric mass aspect are
\begin{equation}\label{phys:eq:nullconstraint}
\partial_{u}\Bigl(\frac{\partial_{u}r}{\Omega^{2}}\Bigr)
=-\frac{4\pi r}{\Omega^{2}}\bigl(\partial_{u}\phi\bigr)^{2},
\qquad
1-\frac{2M}{r}=-\frac{4\,\partial_{u}r\,\partial_{v}r}{\Omega^{2}},
\end{equation}
with the analogous equation in $v$. \TE\ The normalization in
\eqref{phys:eq:nullconstraint} is checked through
$R_{uu}=-2r^{-1}\bigl(r_{uu}-2\Omega_{u}r_{u}/\Omega\bigr)$ together with
\eqref{phys:eq:Einsteinscalar}, and the areal-radius gradient norm is checked
as well. \TA\ These give explicit constraints and invariants to test; they are
not a promise that arbitrary characteristic data admit a Hahn solution. For a
candidate expansion, publish the coefficient domain, the scale ordering, the
equations solved at each order, the residual valuation, and real error
information where available --- and check both evolution equations \emph{and}
constraints, since a proof that one reduced ODE is satisfied does not establish
a solution of the coupled field equations.

\subsection{Stage D: prove a realization bridge}
\label{phys:sub:stageD}

The most useful single result would connect a formally certified solution in a
named asymptotic class to an ordinary solution with a controlled remainder. In
an ODE reduction this might use an existing extension or summation theorem
after verifying its hypotheses; in a PDE setting it would need appropriate
analytic estimates and a common domain, with the selection of boundary or
Stokes data explicit. A substantive target is a selected Fuchsian or quiescent
Einstein--matter construction formulated in a differential Hahn setting, with a
proof that its shadow and remainder bounds reproduce the known classical
solution, including the constraints, gauge conditions, and compatibility of
evolution with the constraints. Only then should one ask whether a larger
support class yields a genuinely stronger result.

The concrete obstruction of \cref{phys:sub:varyingexponent} --- position-dependent
exponents --- is exactly the kind of thing a prototype should expose. And the
success criterion is not a longer formal expansion; it is an expansion whose
interpretation can be checked against a real family or against an independently
defined generalized dynamics. A failure of realization can itself reveal
missing sectors or inconsistent assumptions.

\subsection{Stage E: only then, test a singularity-resolution hypothesis}
\label{phys:sub:stageE}

Only after the conservative benchmarks should one propose a modified action, a
constitutive law, a quantum effective equation, or a nonclassical observable
algebra. State which classical assumptions are changed. Then test regularity in
a specified sense, continuation of physical probes, preservation or replacement
of constraints, stability, and ordinary-limit agreement.

For a regular-core hypothesis, specify whether the core scale is a fixed
positive real, an asymptotic parameter, or a literal infinitesimal, and check
the distinction of \cref{phys:warn:twostatements}. An alternative target is a
controlled inner-layer model with a finite physical scale whose source comes
from a stated action or effective equations rather than being chosen solely to
cancel a curvature pole; examine causal structure, regularity, energy
conditions, stability under perturbations, and matching to the exterior. A
finite central scalar is one benchmark among several, not the complete result.
For a generalized spacetime hypothesis, define its topology, curves and notion
of maximal extension \emph{before} claiming completeness. For a quantum
hypothesis, define the state space and its observable map before claiming
unitary evolution through the singular regime. For a compatible quantum
dynamics, begin with finite systems whose Hamiltonian, real-time evolution and
phase calculus are specified together; prove reduction to the ordinary
Schr\"odinger equation under common-domain limited-coefficient hypotheses;
treat large-phase and long-time scalings as separate asymptotic regimes; and
include the explicit test \eqref{phys:eq:constantphase} in any proposed phase
interface.

The decisive deliverable would be a model with a definite observable
consequence, or a stronger mathematical existence theorem --- not an expression
in which the old divergent quantity has been assigned an infinite name. The
program permits both positive and negative results and states criteria for
recognizing each.

\subsection{Four concrete projects, and what would count as an advance}
\label{phys:sub:projectsAD}

\begin{description}[style=nextline,leftmargin=0pt,labelsep=0pt]
\item[Project A: certified multiscale infall.]
Develop the crossover problem \eqref{phys:eq:crossoverexact} in a controlled
Puiseux or Hahn setting, prove uniform remainder estimates on matching
domains, and compare with high-precision real integration. Success is a theorem
connecting the symbolic sectors to actual geodesics, not a list of formal
dominant terms.
\item[Project B: a regular-core singular perturbation problem.]
Use \eqref{phys:eq:hayward} or another explicitly specified model to track the
outer Schwarzschild and inner-core expansions, their overlap, and invariant
curvature bounds as functions of the regulator. Success distinguishes
regularity at fixed regulator from uniform control as the regulator vanishes,
and identifies the required stress tensor.
\item[Project C: a wave and resurgence benchmark.]
Select one black-hole wave equation with known boundary conditions and
well-tested frequencies, construct its transseries representation, prove the
needed support and summation statements, and recover the physical solution.
Success is reproducible accuracy or a new theorem about the asymptotic
structure, not a numerical coincidence from an unspecified infinite sum.
\item[Project D: a genuine non-Archimedean modification.]
Only after the preceding interfaces are understood, choose an explicit new
dynamical or operational postulate. Prove a well-posedness result in a reduced
model, show what replaces the classical endpoint, and determine whether the
ordinary observable prediction differs from the corresponding real model.
Success would be a substantive physical proposal; \emph{no such postulate is
forced by surreal arithmetic alone.}
\end{description}

\TA\ A mathematically convincing advance would be a new controlled solution
class, a sharper singular-asymptotic estimate, a support-based algorithm that
detects invalid expansions, a rigorous matching theorem that is materially
easier or stronger in the new language, an extension theorem for a relevant
gravitational solution class compatible with gauge transformations and
controlled remainders, or a spectral result for a precisely defined surcomplex
wave operator. Such results would justify surreal methods \emph{even with no
new physical prediction}.

A physically convincing advance would additionally identify a well-defined
finite observable, show how it is computed, prove recovery of the appropriate
classical or quantum regime, and explain why the result is independent of
unphysical choices. For a singularity-resolution claim it would specify what
continues through the formerly singular regime and prove the relevant
continuation or dynamical theorem. This report offers testable mathematical
targets. It offers no completed quantum-gravity model.

\section{Final assessment}
\label{phys:sec:conclusion}

\paragraph{\TA\ As a language for near-singularity behaviour: yes, with
substantial potential.}
Surreal and surcomplex methods record distinct unlimited and infinitesimal
scales exactly. Fractional powers, logarithms, exponentials and suitably
supported series expose structure that a single divergent limit suppresses: the
infall exponents \eqref{phys:eq:radmodes}--\eqref{phys:eq:transmodes}, the
competing angular-momentum scales of
\cref{phys:comp:sectors}, the logarithmic weights of
\cref{phys:comp:weaktide}, the exponentially small sector of
\cref{phys:comp:beyond}, and the inner/outer criteria of
\cref{phys:prop:matching} are all concrete gains. The established differential
and integration theories provide a serious foundation for selected extensions
beyond elementary examples, and the repository's separation of algebra,
summability, domains and topology is an asset rather than a technical
inconvenience.

\paragraph{\TA\ As an automatic cure for black-hole singularities: no.}
This is the report's conclusion and it is negative. The Schwarzschild invariant
remains exactly $48m^{2}/r^{6}$; under a formal radial change $r=s^{q}w(s)$
with $w(0)\ne0$ the pole order becomes exactly $6q$ with the leading
coefficient unchanged, verified for $q=1,2,3$, so no radial reparametrization
removes the blow-up. The proper-time invariant remains $64/(27u^{4})$; the
proper time from rest at $R$ to $r=0$ is $(\pi/2)\sqrt{R^{3}/2M}$, finite ---
and for infinitesimal $R$ it is infinitesimal, hence still finite, so the
infall terminates either way. The tidal component $2M(2M-r)/r^{4}$ diverges,
and the tidal equations retain their singular coefficients. The identity
$r^{6}\Kret=48m^{2}$ has no field-valued solution at $r=0$ in any
characteristic-zero field. A regular-core model changes the metric and its
effective source, and making its core scale infinitesimal does not make its
central curvature finite in real units. What is true is that surreal numbers
make the hierarchy near a singularity \emph{exactly representable} while every
invariant obstruction survives. The sentence to refuse, in any form, is that
surreal numbers resolve the black-hole singularity.

\paragraph{\TA\ As a fundamental replacement for all real and complex physics:
a possible research direction, not an established physical theory.}
Finite algebra transfers well: \cref{phys:thm:spectral} and
\cref{phys:thm:quantumshadow} are real results. Topology, differentiation,
integration, infinite-dimensional quantum theory and empirical interpretation
require explicit choices that cannot be inferred from the field alone. In the
full fine topology, ordinary real-time continuous curves into $\SC$ are
constant; coefficientwise dynamics or a different topology avoids that, but
must be declared. There are already rich positive mathematical constructions,
and their mutual compatibility must be \emph{proved} rather than assumed ---
\cref{phys:prop:phase} and \cref{phys:cor:constantphase} illustrate that
requirement sharply, and they do so about the most canonical construction
available.

\paragraph{\TA\ The best-supported next step.}
Use surreal ideas to strengthen \emph{the analysis of physical models},
especially multiscale and singular perturbation problems, while retaining a
clear ordinary observational limit: retain ordinary spacetime and operational
observables, use a carefully chosen formal or Hahn coefficient structure,
connect it to actual solutions by proved realization and matching results, and
formalize the algebraic interfaces. Such a program could produce valuable
mathematics --- reusable even if surreals never become fundamental physical
scalars --- and might eventually help a physical theory of singular regimes. A
more radical theory remains possible, but its singularity-resolving content
must lie in explicit new dynamics or a rigorously defined replacement of the
classical state space. The larger numerical universe supplies expressive
resources. Dynamics and evidence must supply the physics.
\clearpage
\appendix
\part{Appendices}
\label{phys:part:app}

\section{Derivation of the curvature formulas}
\label{phys:app:curvature}

\subsection{The general static spherical metric}
\label{phys:app:sub:static}

For
\[
\dd s^{2}=-f(r)\dd t^{2}+\frac{\dd r^{2}}{f(r)}+r^{2}\dd\Omega^{2},
\]
work first in a region $f>0$ and choose the orthonormal coframe
\[
e^{0}=\sqrt f\,\dd t,\quad e^{1}=f^{-1/2}\dd r,\quad
e^{2}=r\,\dd\theta,\quad e^{3}=r\sin\theta\,\dd\phi .
\]
With the curvature convention
\begin{equation}\label{phys:eq:riemannconv}
R^{a}{}_{bcd}=\partial_{c}\Gamma^{a}_{db}-\partial_{d}\Gamma^{a}_{cb}
+\Gamma^{a}_{ce}\Gamma^{e}_{db}-\Gamma^{a}_{de}\Gamma^{e}_{cb},
\qquad R_{bd}=R^{a}{}_{bad},
\end{equation}
the independent orthonormal components, up to the Riemann symmetries, may be
taken as
\begin{align*}
R_{0101}&=\tfrac12 f'', &
R_{0202}=R_{0303}&=\frac{f'}{2r},\\
R_{1212}=R_{1313}&=-\frac{f'}{2r}, &
R_{2323}&=\frac{1-f}{r^{2}}.
\end{align*}
A simultaneous overall sign change from another Riemann convention does not
change $\Kret$. Each listed squared component has multiplicity four in the full
contraction: four entries of magnitude $f'/(2r)$, one of magnitude $f''/2$, and
one of magnitude $(1-f)/r^{2}$. Therefore
\[
\Kret=4\Bigl[\frac{(f'')^{2}}{4}+4\frac{(f')^{2}}{4r^{2}}
+\frac{(1-f)^{2}}{r^{4}}\Bigr],
\]
which is \eqref{phys:eq:Kgeneral}; contracting before squaring gives
\eqref{phys:eq:Rgeneral}.

\TA\ These coordinate expressions are rational differential identities in $f$
and $r$. They continue to regions $f<0$ where the metric is nondegenerate, and
to a regular horizon when a nonsingular chart such as \eqref{phys:eq:ef} is
used; the static orthonormal coframe itself is not asserted to exist with that
real square root inside the horizon. Substitution into a Hahn field preserves
the resulting rational identities wherever the denominators remain nonzero,
and no transfer theorem for arbitrary real functions is invoked. Each of the
three companion programs performs the coordinate calculation from
\eqref{phys:eq:riemannconv} without relying on a static orthonormal frame
across the horizon.

For $f=1-2m/r$ the three contributions are $16m^{2}/r^{6}$ each, giving
\eqref{phys:eq:Kschwarz}; for $f=1-r^{2}/\ell^{2}+O(r^{5})$ they tend to
$4/\ell^{4}$, $16/\ell^{4}$ and $4/\ell^{4}$, giving
\eqref{phys:eq:haywardK}.

\subsection{Kasner and the tidal coefficients}
\label{phys:app:sub:kasner}

For the Kasner metric \eqref{phys:eq:kasner} the analogous components have
magnitudes $p_{j}(p_{j}-1)/\tau^{2}$ for the time--space planes and
$p_{j}p_{k}/\tau^{2}$ for the spatial planes; counting the same multiplicity
four gives \eqref{phys:eq:kasnerK}, and the vacuum constraints in
\eqref{phys:eq:kasner} follow by contracting these components.

For Schwarzschild, the time--radial component yields the radial stretching
coefficient $2m/r^{3}$ and the transverse components yield $-m/r^{3}$ in the
geodesic-deviation equation. A radial boost preserves these principal tidal
coefficients, which is why \eqref{phys:eq:jacobi} holds in a radially freely
falling frame and not only in the static frame.

\subsection{Dimensional conventions}
\label{phys:app:sub:dimensions}

All curvature powers are stated with $[\Kret]=\text{length}^{-4}$. The
proper-time variable in \eqref{phys:eq:properr} is a length; if
$\Delta\tau_{\rm sec}$ is measured in seconds, replace $u$ by
$c\,\Delta\tau_{\rm sec}$, so that for the $E=1$ radial infall
\[
\Kret=\frac{64}{27\,c^{4}(\Delta\tau_{\rm sec})^{4}} .
\]
This convention prevents dimensional constants from being mistaken for new
asymptotic exponents.

\begin{warning}[\TA\ $\Kret$ is not a universal regularity criterion]
\label{phys:warn:notuniversal}
The Kretschmann scalar is a convenient diagnostic for the Schwarzschild and
regular-core examples, where it is positive. It is \emph{not} a universal
positive norm of curvature in Lorentzian geometry. More general spacetimes can
require frame components or other invariants to detect a singular regime, and
Kretschmann finiteness alone constrains no higher derivatives. \emph{No
conclusion in this report equates a bounded Kretschmann scalar with general
spacetime regularity.}
\end{warning}

\section{The verification programs: what they check and what they do not}
\label{phys:app:verification}

\subsection{The three suites}
\label{phys:app:sub:suites}

The three source programs are preserved unmodified in \texttt{code/} and their
recorded outputs in \texttt{data/}. All three use exact symbolic arithmetic
with a computer algebra system, compare residuals exactly against zero ---
never against a numerical tolerance --- and raise on any nonzero residual, so a
failure aborts rather than being reported as a pass.

\begin{longtable}{@{}L{0.16\textwidth}L{0.09\textwidth}L{0.68\textwidth}@{}}
\caption{The three verification suites, as recorded and as re-run.}
\label{phys:tab:suites}\\
\toprule
Member & Checks & What it covers \\
\midrule
\endfirsthead
\toprule
Member & Checks & What it covers \\
\midrule
\endhead
\bottomrule
\endfoot
12 & 45 &
Independent coordinate Riemann, Ricci and mixed Einstein reconstruction for the
general static spherical metric, matching \eqref{phys:eq:Kgeneral} and
\eqref{phys:eq:Rgeneral} (checks 01--03); Schwarzschild $\Kret$, vanishing
$\Scal$ and $G^{t}{}_{t}$, horizon curvature, the dimensionless
$48\eps^{-6}$; the ingoing radial--time block determinant; the $E=1$ first
integral; the proper-time coefficients $64/27$, $4/9$, $-2/9$; both indicial
root sets; the Kasner constraints, coefficient $64/27$ and flat case;
regular-core central curvature $24/\ell^{4}$ and $\Scal=12/\ell^{2}$ by limit,
$\rho$ and $\rho(0)$, the Einstein source identity, the expansion through
$r^{5}$, the infinitesimal-core $24\eps^{-4}$; the outer, inner and
transition-scale rational rewrites; eight exact geometric remainders
$N=1,\dots,8$; the boundary-layer $b''(0)=-2$ and derivative scaling; the
Planck threshold; the two-scale exponent; both Gaussian regulator integrals;
two-level Born normalization and the rare-outcome limit. \\
\addlinespace
13 & 50 &
The generic static-spherical Riemann contraction; all sixteen Schwarzschild
Ricci components; $\Kret$; horizon curvature $3/(4m^{4})$; the ingoing
Eddington--Finkelstein radial determinant; the infinitesimal-radius identity;
the radial proper-time cubic relation and the $E=1$ geodesic equation; the
proper-time curvature and both tidal coefficients; all four Jacobi powers
$-1/3$, $4/3$, $1/3$, $2/3$; the Kasner linear and quadratic constraints, the
$64/27$ coefficient and \emph{all three} flat exceptional cases; six Hayward
identities; the curvature power-counting radius; the two-scale classification;
the regulator finite-part reparametrization; the Borel pole residue; two
weak-tide primitives; and the double-null $R_{uu}$ normalization and
areal-radius gradient norm. \\
\addlinespace
14 & 37 &
The Kretschmann invariant from a full four-index Riemann loop and the mixed
Einstein component; Schwarzschild $\Kret$ and its Ricci components;
Eddington--Finkelstein determinant; horizon curvature; the $E=1$ first
integral; $\Kret=64/(27s^{4})$; the radial tidal coefficient and all four
Jacobi exponents; the nonradial $2/5$ leading balance; the \emph{exact
crossover integrand denominator factor} of
\eqref{phys:eq:crossoverexact}; Hayward core curvature and density and the
$\ell\to0$ Schwarzschild limit at $r>0$; the finite-part shift $a\mapsto a-c$;
factorial ODE residuals at truncations $N=1,\dots,8$; and the exponentially
small homogeneous ambiguity. \\
\end{longtable}

All three recorded runs pass. They were re-run for this merge, on
\emph{copies}, and all three pass again on a newer interpreter than the one
recorded: $45$, $50$ and $37$ checks, $132$ in total, every suite exit $0$.
\emph{None of the three writes any file beside itself}; the build helpers
shipped with one member do create a build directory and copy a PDF, which is
why the suites were run on copies as a matter of discipline.

\subsection{What the programs establish, and what they do not}
\label{phys:app:sub:scope}

\begin{warning}[The scope disclaimers are part of the result]
\label{phys:warn:disclaimers}
Each program prints its own scope disclaimer, and those disclaimers are
reproduced here because they are exactly the honest limitations a merge must
carry.
\begin{itemize}[leftmargin=1.6em]
\item Member 12: \emph{algebraic identities only; no claim of a verified
surreal field, general PDE existence, geodesic completion, or a physical
singularity cure.} The program does not implement a surreal number field at
all: substitutions such as $r=m\eps$ use an ordinary positive symbol.
\item Member 13: \emph{general mathematical proofs, asymptotic realization, and
physical claims were not machine-verified.} It does not verify set-theoretic
foundations, global support lemmas, the standard-part reduction theorem,
analytic existence, or the interpretation of quantum probabilities, and it runs
no black-hole evolution numerically. Its role is to catch concrete algebraic
mistakes in the examples, not to confer machine-checked status on the general
proofs.
\item Member 14: \emph{finite symbolic identities only; no general analytic or
physical validation.} It does not verify \cref{phys:prop:discrete},
\cref{phys:thm:einstein} or \cref{phys:thm:spectral} --- those have
ordinary text proofs --- does not construct the full surreal field, does not
prove convergence of any transseries, does not validate a physical model, and
does not simulate across a physical singularity.
\end{itemize}
In the tier language of \cref{phys:conv:tiers}: \emph{the programs certify tier
\TE\ and nothing else.} No statement in tier \TC\ and no statement in tier \TA\
is machine-checked anywhere in this report, and no Lean formalization of this
article is included or claimed.
\end{warning}

\subsection{Independent verification performed for this merge}
\label{phys:app:sub:independent}

Because the central conclusion is negative, and because a merge is most likely
to damage exactly that, the load-bearing facts were re-derived from the metric
rather than read from the sources. Specifically: the Christoffel symbols and
the Riemann tensor of \eqref{phys:eq:schwarz} were built from scratch and
contracted, giving $\Kret=48M^{2}/r^{6}$ exactly, hence $p=6$ in
\cref{phys:prop:pole}; substituting $r=s^{q}(1+s)$ gives a pole of order
exactly $6q$ with leading coefficient $48M^{2}$ unchanged, for $q=1,2,3$;
the proper-time quadrature \eqref{phys:eq:taufinite} evaluates in closed form
to $(\pi/2)\sqrt{R^{3}/2M}$; and the tidal component
$R^{r}{}_{trt}=2M(2M-r)/r^{4}$ diverges as $r\to0$. All four agree with the
manuscripts. The negative conclusion of \cref{phys:sec:conclusion} therefore
rests on independently reproduced computations and not on a source's summary.

\section{Notation and assumptions at a glance}
\label{phys:app:notation}

\begin{longtable}{@{}L{0.21\textwidth}L{0.73\textwidth}@{}}
\toprule
Symbol or phrase & Meaning in this report \\
\midrule
\endfirsthead
\toprule Symbol or phrase & Meaning in this report \\ \midrule
\endhead
\bottomrule
\endfoot
\TE, \TC, \TA, \TI &
Exact symbolic identity; conditional theorem proved in the text; assessment;
imported from the literature. \Cref{phys:conv:tiers}. \\
\addlinespace
$\No$, $\SC$ &
The full proper-class surreal field and its complexification. \emph{Not} an
ordinary set-sized metric space. \\
\addlinespace
$F_{\Gamma}$, $K_{\Gamma}$ &
Set-sized real and complex Hahn fields with well-ordered support and small
monomials $t^{\gamma}$ for $\gamma>0$. \\
\addlinespace
$\val$ &
Least exponent of a nonzero Hahn series; $\val(0)=+\infty$ as a bookkeeping
symbol, never an adjoined field element. \\
\addlinespace
$\Ofin_{F}$, $\mm_{F}$ &
The limited elements (bounded by an ordinary real bound) and the infinitesimals
(smaller in magnitude than every positive ordinary real). \\
\addlinespace
limited / infinitesimal / unlimited &
\Cref{phys:def:limited}; used in place of ``finite/infinite'' to block the
slide from ``is a field element'' to ``is a finite physical magnitude''. \\
\addlinespace
$\st$ &
Real or complex standard part, defined \emph{only} on limited elements. \\
\addlinespace
strong Hahn summability &
Well-ordered union of supports \emph{and} finitely many contributions per
exponent (\cref{phys:def:strongsum}). A support condition, not convergence. \\
\addlinespace
$\eps$ &
A positive dimensionless infinitesimal when working in a Hahn field; a positive
\emph{real} regulator where a real family is explicitly under discussion. The
two uses are always flagged. \\
\addlinespace
$h$, $\pi_{0}$ &
Formal parameter fixed by the coordinate derivatives, and reduction modulo $h$
in Framework F1. \\
\addlinespace
$\Alg(U)$, $\Alg_{\Gamma}(U)$ &
$C^{\infty}(U,\R)[[h]]$ and the Hahn coefficient-function algebra with global
well-ordered support; Frameworks F1 and F2 of \cref{phys:sub:frameworks}. \\
\addlinespace
$\partial_{\mu}$, $\dd/\dd\eps$, $D$ &
Coefficientwise ordinary spacetime derivative; parameter derivative; a scalar
derivation on $\No$. \emph{Never interchangeable}; see
\cref{phys:tab:derivatives}. \\
\addlinespace
$\theta$, $\eta$, $B$ &
Finite angle; infinitesimal phase residue; accumulated phase. Imported
verbatim from the differential-equations report
(\cref{phys:conv:phases}). \\
\addlinespace
$\Oz$, $\ExpEK$, $\Emm$, $\cis$ &
Canonical initial integer part of $\No$; the integer-part-based global
surcomplex exponential; the infinitesimal Hahn exponential; the finite-angle
phase. \\
\addlinespace
$m$, $\ell$, $\lP$, $L$ &
Geometrized mass $GM/c^{2}$; regular-core length scale; Planck length; an
ordinary reference length. \\
\addlinespace
$\tau$, $u$, $s$, $S$ &
Proper time in length units; remaining proper time $u=\tau_{*}-\tau$ (written
$s$ where a source's notation is followed); and the dimensionless remaining
clock of \cref{phys:comp:crossover}. \\
\addlinespace
$\Kret$ &
The Kretschmann scalar $R_{abcd}R^{abcd}$; \emph{not} the coefficient field
$K$. \\
\addlinespace
$\FP$ &
Extraction of the exponent-zero coefficient on \emph{all} Laurent series, as
opposed to $\st$ on $\Ofin$. Non-multiplicative and prescription-dependent
(\cref{phys:comp:FP}). \\
\addlinespace
formal solution / formal certificate &
An identity coefficient by coefficient, with support admissibility and residual
vanishing through a valuation threshold. It asserts \emph{no} convergence,
summability to a real solution, or physical validity
(\cref{phys:sub:certificates}). \\
\end{longtable}

\section{Consolidated register of what is not claimed}
\label{phys:app:nonclaims}

The three manuscripts state $36$ explicit non-claims between them. All of them
survive here, consolidated into $31$ numbered items and not one dropped. This
register is part of the result.

\subsection{Provenance and status}
\label{phys:app:sub:nc-provenance}

\begin{enumerate}[label=\textup{N\arabic*.},leftmargin=3.0em]
\item Every member self-labels as prepared with AI assistance, not refereed,
and makes no research-priority claim for any proposition, computation or
benchmark. The proved results are supplied as explicit mathematical deductions
for an assessment, and this merge adds no priority claim of its own. The
physical equations used in the examples are standard models, not newly
discovered black-hole solutions.
\item No Lean formalization of any member, or of this report, is included or
claimed. No member ran the repository's Lean build, and no declaration-level
audit was performed; neither was one performed for this merge. A theorem in an
unrefereed report and a successfully checked generic Lean declaration have
different evidential status.
\item Repository inspection was targeted: between them the members read the
root README, the documentation catalogue, the surcomplex analysis README (one
member cites the analysis source file) and the trigonometry README, at one
pinned commit. Not every report in the collection was read; no file in the
repository was changed by any member or by this merge.
\item The literature search was targeted, not systematic. Absence of a
predictive model in the inspected sources is not proof that no other proposal
has ever been written, and the survey does not establish that no future
non-Archimedean theory could resolve a singularity --- only that the reviewed
sources do not.
\item The Berarducci--Mantova derivation, the Aschenbrenner--van den
Dries--van der Hoeven universal $H$-field results, van den Dries--Ehrlich
surreal exponentiation, the Ehrlich--Kaplan exponential and trigonometric
constructions, and the Costin--Ehrlich integration theory are all
\emph{imported} with their stated hypotheses and function-class restrictions
intact. None is reproved and none is strengthened. The Ehrlich--Kaplan erratum
concerns in-text bibliography reference numbering, not a withdrawal of the
global phase construction.
\item Penrose incompleteness, the null Raychaudhuri focusing bound, the
Dafermos and Dafermos--Luk interior results, Donoghue's effective-field-theory
corrections, Hollands--Wald, Steinbauer--Vickers, Unruh--Sch\"utzhold,
Horowitz--Marolf, Hatsuda, Hayward, Carballo-Rubio et al.,
Andersson--Rendall, Damour--Henneaux--Nicolai, Benci--Horsten--Wenmackers,
Loeb, Gubser et al.\ and the Nieto papers are cited, not proved here, and
source hypotheses and regularity qualifications are retained. No current
universal result about all black-hole interiors is claimed.
\item One member's package README lists a checksum file that is not present in
the delivered directory. The discrepancy is recorded here rather than silently
repaired, and nothing in this report depends on it.
\end{enumerate}

\subsection{What the programs do not verify}
\label{phys:app:sub:nc-programs}

\begin{enumerate}[label=\textup{N\arabic*.},leftmargin=3.0em,start=8]
\item The three suites check \emph{finite symbolic identities only}: $45+50+37
=132$ checks. They do not construct or verify a surreal field --- indeed they
implement no surreal arithmetic, using ordinary positive symbols for
infinitesimal substitutions --- and they do not verify Hahn summability in
general, set-theoretic foundations, global support lemmas, the standard-part or
formal Einstein reduction theorems, the discreteness proposition, the finite
spectral theorem, analytic existence, PDE existence, geodesic completion,
convergence of any transseries, quantum probability interpretation, or any
proposed physical interpretation.
\item They run no black-hole evolution numerically and do not simulate across a
physical singularity. Their role is to catch algebraic transcription errors in
the worked examples.
\item Re-running them establishes reproducibility on a newer interpreter and
nothing further. One member's README notes that a given Python version is the
intended interface but only the recorded environment was tested, and one
member's build records one minor underfull-box LaTeX diagnostic.
\end{enumerate}

\subsection{Scope of the mathematical results}
\label{phys:app:sub:nc-math}

\begin{enumerate}[label=\textup{N\arabic*.},leftmargin=3.0em,start=11]
\item \Cref{phys:prop:discrete} is proved only for \emph{set-sized} subsets. It
explicitly does not claim that the proper class $\SC$ is a discrete space,
since no radius can be chosen below a proper class of positive distances; it
does not apply to the intrinsic valuation topology of a set-sized Hahn field;
and it does not rule out coarser topologies or coefficientwise smooth
structures.
\item \Cref{phys:prop:pole} and \cref{phys:prop:unlimited} are statements about
Hahn-field valuation and about order in an ordered extension respectively.
Neither is a statement about real limits, and neither rules out a modified
metric or a rescaled observable.
\item \Cref{phys:prop:noendpoint} is conditional on preserving the same
invariant identity at the endpoint. It is not a theorem excluding a modified
metric, a different equation, nonclassical observables, a quantum replacement
of the endpoint, or a geometry with no zero-areal-radius event.
\item \Cref{phys:thm:einstein} and \cref{phys:thm:reduction} assume a
nondegenerate smooth leading metric on one common domain, with strictly
positive order or positive global support for the perturbation. Singular
perturbations are explicitly outside them; \cref{phys:ex:boundarylayer} is
offered as evidence that the hypotheses cannot be dropped. Both are conditional
and neither is machine-verified.
\item \Cref{phys:cor:shadow} is a conditional obstruction and not a universal
no-go. Five escape routes are enumerated in its status note: a degenerate or
nonexistent shadow metric, negative-order fields, singular boundary layers, new
equations, and nonclassical observables or a different notion of spacetime.
\item \Cref{phys:prop:matching} is a criterion on Hahn-field valuation data and
not on real convergence. It detects misuse of an expansion; it invents no
solution and selects no metric.
\item \Cref{phys:thm:quantumshadow} is conditional on adopting real standard
parts as measured probabilities. It excludes conditioning on
standard-probability-zero events, and excludes infinite and hyperfinite
experiments. It is not a theorem that non-Archimedean models are
experimentally indistinguishable in general, and it does not cover
non-standard-part probability interpretations.
\item Infinite-dimensional quantum theory, functional integration, measures and
unbounded self-adjoint operators are \emph{not} supplied by algebraic closure.
A proper-class path sum is not licensed by Hahn summability.
\item \Cref{phys:prop:phase}, \cref{phys:cor:constantphase} and
\cref{phys:prop:periods} concern one specific commutative construction and one
specific normalized derivation. They do not assert that global surcomplex
exponentials cannot be defined --- one is constructed in
\eqref{phys:eq:EKexp} --- and the Ehrlich--Kaplan constructions must not be
described as nonexistent.
\item The realization bridge of \cref{phys:sub:stageD} is open, and one
concrete obstruction is flagged: differentiating $t^{p(x)}$ produces
$p'(x)\log t$, so the constant-exponent coefficient algebra is not closed under
position-dependent exponents.
\item The transseries display \eqref{phys:eq:schematictrans} is schematic. An
arbitrary triple-indexed family is not automatically admissible, and the
schematic $\ee^{\kappa v}v^{-p}$ form is not a universal theorem about all
Kerr interiors.
\end{enumerate}

\subsection{Scope of the physical statements}
\label{phys:app:sub:nc-physics}

\begin{enumerate}[label=\textup{N\arabic*.},leftmargin=3.0em,start=22]
\item \textbf{A physical resolution of black-hole singularities is not supplied
by scalar extension and is not claimed anywhere in this report.} This report is
a mathematical assessment and a research proposal, not a quantum-gravity
resolution of black-hole singularities, and not a claim of a new empirically
validated theory, a new black-hole solution, or a new prediction.
\item All negative conclusions are scoped as \emph{missing implications}, not
impossibility theorems. A future surreal-variable theory could still make
nontrivial predictions; a quantum bounce is not excluded; and
non-Archimedean probability interpretations other than standard-part are not
covered.
\item $r_{\mathrm Q}$ of \eqref{phys:eq:rq} is dimensional power counting, not
a derived universal threshold, not a sharply defined transition radius, and not
a derived minimum radius. Additional scales or state-dependent effects can
invalidate a simple curvature criterion sooner. The $\eps^{\alpha}$,
$\eps^{\beta}$ and $\eps^{a}$, $\eps^{b}$ assignments describe a \emph{family
of models} and do \emph{not} assert that the measured Planck length is
literally infinitesimal; the Planck length is read throughout as an ordinary
small positive real.
\item The regular-core effective stress-energy $T=G/8\pi$ is a consistent
effective description with no microphysical Lagrangian and no stability proof.
Finite central curvature is not global geodesic completeness, unique evolution,
or stable evaporation, and the model's global viability, stability and quantum
derivation are left open.
\item The regular-core ansatz is \emph{not} advertised as $C^{\infty}$ at the
centre: the $r^{5}$ term of \eqref{phys:eq:haywardcenter} becomes $\abs x^{5}$
along a Cartesian line through the centre. ``Regular'' there means finite
curvature plus $C^{2}$-level regularity, nothing more.
\item Kretschmann finiteness alone constrains no higher derivatives, and
$\Kret$ is not a universal positive norm of Lorentzian curvature
(\cref{phys:warn:notuniversal}).
\item Generic oscillatory Kasner interiors may admit no single well-ordered
expansion; Hardy-field and transseries technology targets eventually ordered
nonoscillatory behaviour, so BKL-type interiors may need extra variables,
sectors, or a different function category. A Schwarzschild Puiseux law
licenses no inference about generic interiors, and charged and rotating
interiors are explicitly not Schwarzschild copies.
\item Generic Jacobi rates in
\eqref{phys:eq:radmodes}--\eqref{phys:eq:transmodes} do not apply to every
Jacobi field; special initial conditions change them. The crossover limit of
\cref{phys:prop:crossover} is uniform only on bounded positive $R$-intervals,
and the large-$R$ behaviour of the limiting function is a matched-asymptotic
statement, not uniform convergence on an unbounded interval. The near-extremal
bookkeeping of \cref{phys:comp:rn} is not a claim that extremality or a Cauchy
horizon is physically harmless.
\item The conservative architecture \eqref{phys:eq:architecture} is called the
most defensible starting point, not a derived necessity. Taking real standard
parts as measured probabilities is a possible modelling choice, not a
consequence of the field axioms. Declaring a dimensionless ratio infinitesimal
is a stated convention; the report insists that ``the radius is
$\omega^{-1}$'', without a unit, is not well defined. The $s$-weighted tidal
integral is a model criterion for singularity strength, and
$\lP^{4}\Kret\sim1$ is a large-correction indicator; signature $(-,+,+,+)$ and
$G=c=1$ are conventions, and treating $\eps$ as a scalar constant with
$\partial_{\mu}\eps=0$ is a choice.
\item Everything in \cref{phys:sec:program} --- the architecture, the five
stages, the four benchmarks, the four projects, the proposed Lean module names,
the five-item extension plan, the suggestions that a reduced gravitational ODE
might fall in a resurgent class or that surreal methods might describe the
incoming and outgoing regions of a bounce, the five-level hierarchy of success
and the four-question test --- is a \emph{proposal} or conceptual framing.
None of it is an implemented or experimentally validated theory, and none of it
is a description of an existing module.
\end{enumerate}

\section{Provenance: member by member}
\label{phys:app:provenance}

\begin{longtable}{@{}L{0.30\textwidth}L{0.11\textwidth}L{0.52\textwidth}@{}}
\caption{Where each statement came from. ``All three'' means the shared spine,
printed once.}
\label{phys:tab:provenance}\\
\toprule
Statement & Members & Note \\
\midrule
\endfirsthead
\toprule
Statement & Members & Note \\
\midrule
\endhead
\bottomrule
\endfoot
Static-spherical $\Kret[f]$, Schwarzschild $\Kret$, horizon determinant,
vacuum $\Ric$, proper-time infall, Jacobi exponents, Kasner coefficient,
Hayward central limits, $r_{\mathrm Q}$, the two-scale exponent, the
finite-part shift, the framework comparison, the research program &
all three & Printed once. Each member ships its own independent symbolic
re-derivation of the connection and curvature. \\
\addlinespace
\Cref{phys:prop:st}, \cref{phys:prop:discrete},
\cref{phys:thm:spectral} & all three & Same proofs; printed once, with the
union of statements and caveats. \\
\addlinespace
\Cref{phys:thm:quantumshadow} & all three & Union of three overlapping
versions; the conditioning boundary (d) is from member 13, the projection and
joint-probability clauses from member 12, the Born-weight clause from member
14. \\
\addlinespace
\Cref{phys:prop:pole}, ramification invariance &
12 & Route 1. Unique to member 12. \\
\addlinespace
\Cref{phys:prop:unlimited} & 14 & Route 2. Unique to member 14, and a different
theorem; both are kept (\cref{phys:rem:tworoutes}). \\
\addlinespace
\Cref{phys:prop:noendpoint} & 13, 14 & Same one-line proof; printed once.
Member 13 names the published proposal as its target. \\
\addlinespace
\Cref{phys:thm:einstein} (Framework F1) & 12, 14 & Same theorem in two
parametrizations; printed once. \\
\addlinespace
\Cref{phys:thm:reduction}, \cref{phys:cor:shadow} (Framework F2) & 13 &
Genuinely different hypotheses and a different proof ingredient (the Neumann
support lemma); kept separately. \\
\addlinespace
\Cref{phys:ex:boundarylayer} & 12 & The boundary-layer warning about
differentiating standard parts. \\
\addlinespace
\Cref{phys:ex:oscillatory} & 14 & A purely real counterexample to the
interchange. \\
\addlinespace
\Cref{phys:prop:matching}, \cref{phys:comp:remainder} & 12 & The exact
inner/outer Hahn support criteria, sharp at $\alpha=2/3$. \\
\addlinespace
\Cref{phys:comp:crossover}, \cref{phys:prop:crossover},
\cref{phys:comp:sectors} & 14 & The small-angular-momentum crossover. \\
\addlinespace
\Cref{phys:sec:control} as a headline & 14 & Member 12 makes the same argument
inside its effective-field-theory section; member 13 in its
ultraviolet section; member 14 promotes it, and this report follows member 14
on placement. \\
\addlinespace
\Cref{phys:sub:finitephase}, \cref{phys:prop:phase},
\cref{phys:cor:constantphase}, \cref{phys:rem:naivetaylor} & 12 &
The global-phase section. \\
\addlinespace
\Cref{phys:prop:periods} & 13 & Consistent with member 12's construction;
\cref{phys:rem:periodsagree} records why. \\
\addlinespace
\Cref{phys:prop:frame}, \cref{phys:lem:balance},
\cref{phys:comp:weaktide}, \cref{phys:comp:borel},
\eqref{phys:eq:doublenull}--\eqref{phys:eq:nullconstraint},
\cref{phys:comp:FP} first half & 13 & Bounded frames, dominant balance, weak
singularities, the Borel residue, the double-null identities, and
non-multiplicativity of $\FP$. \\
\addlinespace
\Cref{phys:comp:beyond}, \cref{phys:comp:rn},
\cref{phys:sub:qnm}, \cref{phys:comp:FP} second half,
\cref{phys:warn:hayward}(a) & 14 & Beyond-all-orders ambiguity, near-extremal
bookkeeping, the wave benchmark, noncanonical subtraction, and the
$\abs x^{5}$ remark. \\
\addlinespace
\Cref{phys:comp:gaussian}, \cref{phys:ex:rare}, the units-and-falsifiability
material of \cref{phys:sub:units,phys:sub:novelty} & 12 & \\
\addlinespace
\Cref{phys:comp:distrib}, \cref{phys:ex:amplify},
\cref{phys:sec:measurement} recoverability material,
\cref{phys:sub:stageA}--\cref{phys:sub:stageE} stage structure with reported
failure conditions & 13 & \\
\addlinespace
\Cref{phys:sub:projectsAD}, \cref{phys:sub:minimalspec},
\cref{phys:sub:fourquestions} & 14 & \\
\addlinespace
The claim tiers of \cref{phys:conv:tiers} and
\cref{phys:tab:ledger} & merge & Synthesized from the three members' own claim
ledgers, which they state separately; the tier tags and their use throughout
are this report's contribution to the merge and are not a new mathematical
claim. \\
\end{longtable}

\TA\ No statement in this report rests on a member that does not prove what it
is cited for. Where the three disagreed, they disagreed about \emph{emphasis
and placement} and not about content: the one substantive editorial conflict
was whether the loss-of-control argument is a headline or a late aside, and
\cref{phys:warn:placement} records its resolution.
\clearpage
\phantomsection
\addcontentsline{toc}{section}{References}
\begingroup
\raggedright
\begin{thebibliography}{99}
\small
\setlength{\itemsep}{0.35em}

\bibitem{phys:RepoRoot}
V.~Reshetnikov, \emph{Surreal}, repository root \texttt{README.md}, pinned
snapshot \Commit.
\url{https://github.com/VladimirReshetnikov/Surreal/blob/39f2be6667ade51bca2b45daa47e289d69c09764/README.md}.
Accessed 21 September 2026. Scope and stated limitations of the implemented
Lean prerequisites.

\bibitem{phys:RepoDocs}
V.~Reshetnikov, \emph{The surreal and surcomplex reports}, repository
documentation catalogue \texttt{docs/README.md}, same snapshot.
\url{https://github.com/VladimirReshetnikov/Surreal/blob/39f2be6667ade51bca2b45daa47e289d69c09764/docs/README.md}.
AI-assisted research documentation; includes the coefficient-category warnings
and the status of the reports.

\bibitem{phys:RepoAnalysis}
\emph{Surcomplex Analysis: Germs, Coherent Hahn Sections, and Infinitesimal
Zero Geometry}, research report in \texttt{Surreal},
\texttt{docs/surcomplex/analysis/}, same snapshot. Distinguishes Hahn
summability from topological convergence, and separates germ, common-domain
and formal coefficient classes. Not treated here as a verified physics
foundation.

\bibitem{phys:RepoTrig}
\emph{Trigonometry on the Surcomplex Plane: Canonical Phases, Arbitrary-Scale
Geometry, and Infinitesimal Degeneration}, research report in \texttt{Surreal},
\texttt{docs/surcomplex/trigonometry/}, same snapshot. Includes the
Ehrlich--Kaplan integer-part global normalization and explicitly selects no
scalar derivation.

\bibitem{phys:RepoDiffEq}
\emph{Differential Algebra and Differential Equations over the Surreal and
Surcomplex Numbers: the finite-primitive criterion, the phase obstruction, and
support-certified coordinate systems}, merged research report in
\texttt{Surreal}, \texttt{docs/surcomplex/differential-equations/}. Owns the
oscillation obstruction and the three phase symbols adopted in
\cref{phys:conv:phases}.

\bibitem{phys:RepoFound}
\emph{Foundations for Surreal and Surcomplex Mathematics: Sets, Classes,
Universes, Type Theories, and a Pinned Lean Audit}, merged research report in
\texttt{Surreal}, \texttt{docs/foundations-and-computation/foundations/}.
Owns the workspace-localization theorem and the size obstructions cited in
\cref{phys:sub:workspace}.

\bibitem{phys:Conway}
J.~H.~Conway, \emph{On Numbers and Games}, second edition, A K Peters, 2001;
first edition, Academic Press, 1976.

\bibitem{phys:Gonshor}
H.~Gonshor, \emph{An Introduction to the Theory of Surreal Numbers}, London
Mathematical Society Lecture Note Series 110, Cambridge University Press, 1986.

\bibitem{phys:Ehrlich}
P.~Ehrlich, The absolute arithmetic continuum and the unification of all
numbers great and small, \emph{Bulletin of Symbolic Logic} \textbf{18} (2012),
1--45. \url{https://doi.org/10.2178/bsl/1327328438}.

\bibitem{phys:vdDE}
L.~van den Dries and P.~Ehrlich, Fields of surreal numbers and exponentiation,
\emph{Fundamenta Mathematicae} \textbf{167} (2001), 173--188.
\url{https://doi.org/10.4064/fm167-2-3}.

\bibitem{phys:BM}
A.~Berarducci and V.~Mantova, Surreal numbers, derivations and transseries,
\emph{Journal of the European Mathematical Society} \textbf{20} (2018),
339--390. \url{https://doi.org/10.4171/JEMS/769};
\url{https://arxiv.org/abs/1503.00315}.

\bibitem{phys:BMgerms}
A.~Berarducci and V.~Mantova, Transseries as germs of surreal functions,
\emph{Transactions of the American Mathematical Society} \textbf{371} (2019),
3549--3592. \url{https://doi.org/10.1090/tran/7428};
\url{https://arxiv.org/abs/1703.01995}.

\bibitem{phys:ADH}
M.~Aschenbrenner, L.~van den Dries and J.~van der Hoeven, The surreal numbers
as a universal $H$-field, \url{https://arxiv.org/abs/1512.02267}.

\bibitem{phys:EK}
P.~Ehrlich and E.~Kaplan, Surreal ordered exponential fields, \emph{Journal of
Symbolic Logic} \textbf{86} (2021), 1066--1115.
\url{https://doi.org/10.1017/jsl.2021.59}; \url{https://arxiv.org/abs/2002.07739}.
In particular the section on trigonometric fields and surcomplex
exponentiation.

\bibitem{phys:EKErratum}
P.~Ehrlich and E.~Kaplan, Surreal ordered exponential fields --- Erratum,
\emph{Journal of Symbolic Logic} \textbf{87} (2022), 871.
\url{https://doi.org/10.1017/jsl.2022.5}. Concerns in-text bibliography
reference numbering.

\bibitem{phys:CE}
O.~Costin and P.~Ehrlich, Integration on the surreals, \emph{Advances in
Mathematics} \textbf{452} (2024), article 109823.
\url{https://arxiv.org/abs/2208.14331}, version 5 consulted. Used only with the
stated function-class restrictions.

\bibitem{phys:Penrose}
R.~Penrose, Gravitational collapse and space-time singularities,
\emph{Physical Review Letters} \textbf{14} (1965), 57--59.
\url{https://doi.org/10.1103/PhysRevLett.14.57}.

\bibitem{phys:Raychaudhuri}
A.~Raychaudhuri, Relativistic cosmology.~I, \emph{Physical Review} \textbf{98}
(1955), 1123--1126. \url{https://doi.org/10.1103/PhysRev.98.1123}.

\bibitem{phys:Witten}
E.~Witten, Light rays, singularities, and all that, \emph{Reviews of Modern
Physics} \textbf{92} (2020), 045004.
\url{https://doi.org/10.1103/RevModPhys.92.045004};
\url{https://arxiv.org/abs/1901.03928}.

\bibitem{phys:Senovilla}
J.~M.~M.~Senovilla, A critical appraisal of the singularity theorems,
\url{https://arxiv.org/abs/2108.07296}; \url{https://doi.org/10.1098/rsta.2021.0174}.

\bibitem{phys:Landsman}
K.~Landsman, Penrose's 1965 singularity theorem: from geodesic incompleteness
to cosmic censorship, \url{https://arxiv.org/abs/2205.01680}, 2022.

\bibitem{phys:Dafermos}
M.~Dafermos, The interior of charged black holes and the problem of uniqueness
in general relativity, \url{https://arxiv.org/abs/gr-qc/0307013}.

\bibitem{phys:DafermosLuk}
M.~Dafermos and J.~Luk, The interior of dynamical vacuum black holes I: the
$C^{0}$-stability of the Kerr Cauchy horizon, \emph{Annals of Mathematics}
\textbf{202} (2025), 309--630.
\url{https://doi.org/10.4007/annals.2025.202.2.1};
\url{https://arxiv.org/abs/1710.01722}. The use here is restricted to the
hypotheses and interior conclusion of this work.

\bibitem{phys:Donoghue}
J.~F.~Donoghue, General relativity as an effective field theory: the leading
quantum corrections, \emph{Physical Review D} \textbf{50} (1994), 3874--3888.
\url{https://doi.org/10.1103/PhysRevD.50.3874};
\url{https://arxiv.org/abs/gr-qc/9405057}.

\bibitem{phys:DonoghueReview}
J.~F.~Donoghue, Quantum general relativity and effective field theory,
\url{https://arxiv.org/abs/2211.09902}. Review contribution to the
\emph{Handbook of Quantum Gravity}.

\bibitem{phys:AR}
L.~Andersson and A.~D.~Rendall, Quiescent cosmological singularities,
\emph{Communications in Mathematical Physics} \textbf{218} (2001), 479--511.
\url{https://doi.org/10.1007/s002200100406};
\url{https://arxiv.org/abs/gr-qc/0001047}.

\bibitem{phys:DHN}
T.~Damour, M.~Henneaux and H.~Nicolai, Cosmological billiards, \emph{Classical
and Quantum Gravity} \textbf{20} (2003), R145--R200.
\url{https://arxiv.org/abs/hep-th/0212256}.

\bibitem{phys:Hayward}
S.~A.~Hayward, Formation and evaporation of non-singular black holes,
\emph{Physical Review Letters} \textbf{96} (2006), 031103.
\url{https://doi.org/10.1103/PhysRevLett.96.031103};
\url{https://arxiv.org/abs/gr-qc/0506126}.

\bibitem{phys:CR}
R.~Carballo-Rubio, F.~Di~Filippo, S.~Liberati and M.~Visser, Geodesically
complete black holes, \emph{Physical Review D} \textbf{101} (2020), 084047.
\url{https://doi.org/10.1103/PhysRevD.101.084047};
\url{https://arxiv.org/abs/1911.11200}.

\bibitem{phys:Bhattacharjee}
S.~Bhattacharjee, S.~Sarkar and A.~Virmani, Internal structure of charged AdS
black holes, \url{https://arxiv.org/abs/1604.03730}, 2016.

\bibitem{phys:Unruh}
W.~G.~Unruh and R.~Sch\"utzhold, On the universality of the Hawking effect,
\url{https://arxiv.org/abs/gr-qc/0408009}.

\bibitem{phys:HW}
S.~Hollands and R.~M.~Wald, Quantum fields in curved spacetime, \emph{Physics
Reports} (2015). \url{https://doi.org/10.1016/j.physrep.2015.02.001};
\url{https://arxiv.org/abs/1401.2026}.

\bibitem{phys:HM}
G.~T.~Horowitz and D.~Marolf, Quantum probes of spacetime singularities,
\emph{Physical Review D} \textbf{52} (1995), 5670--5675.
\url{https://doi.org/10.1103/PhysRevD.52.5670};
\url{https://arxiv.org/abs/gr-qc/9504028}.

\bibitem{phys:Hatsuda}
Y.~Hatsuda, Quasinormal modes of black holes and Borel summation,
\emph{Physical Review D} \textbf{101} (2020), 024008.
\url{https://doi.org/10.1103/PhysRevD.101.024008};
\url{https://arxiv.org/abs/1906.07232}.

\bibitem{phys:SV}
R.~Steinbauer and J.~A.~Vickers, The use of generalised functions and
distributions in general relativity, \emph{Classical and Quantum Gravity}
\textbf{23} (2006), R91--R114. \url{https://arxiv.org/abs/gr-qc/0603078}.

\bibitem{phys:Resurgence}
I.~Aniceto, G.~Ba\c{s}ar and R.~Schiappa, A primer on resurgent transseries and
their asymptotics, \url{https://arxiv.org/abs/1802.10441};
\url{https://doi.org/10.1016/j.physrep.2019.02.003}.

\bibitem{phys:Kontsevich}
M.~Kontsevich, Deformation quantization of Poisson manifolds, I,
\url{https://arxiv.org/abs/q-alg/9709040}.

\bibitem{phys:BHW}
V.~Benci, L.~Horsten and S.~Wenmackers, Non-Archimedean probability,
\emph{Milan Journal of Mathematics} \textbf{81} (2013), 121--151.
\url{https://arxiv.org/abs/1106.1524}.

\bibitem{phys:Loeb}
P.~A.~Loeb, Conversion from nonstandard to standard measure spaces and
applications in probability theory, \emph{Transactions of the American
Mathematical Society} \textbf{211} (1975), 113--122.
\url{https://www.jstor.org/stable/1997222}.

\bibitem{phys:Padic}
S.~S.~Gubser, J.~Knaute, S.~Parikh, A.~Samberg and P.~Witaszczyk, $p$-adic
AdS/CFT, \emph{Communications in Mathematical Physics} (2017).
\url{https://doi.org/10.1007/s00220-016-2813-6};
\url{https://arxiv.org/abs/1605.01061}.

\bibitem{phys:Nieto2016}
J.~A.~Nieto, Some mathematical and physical remarks on surreal numbers,
\emph{Journal of Modern Physics} (2016).
\url{https://doi.org/10.4236/jmp.2016.715188};
\url{https://arxiv.org/abs/1611.09699}. The black-hole discussion appears near
the end of the 19-page preprint. Cited as a prior proposal, not as an assumed
singularity-resolution theorem.

\bibitem{phys:Nieto2018}
J.~A.~Nieto, Duality, matroids, qubits, twistors and surreal numbers,
\url{https://arxiv.org/abs/1810.04521}, 2018.

\end{thebibliography}
\endgroup
\end{document}
